\documentclass[11pt]{article}
\usepackage[a4paper,margin=2.4cm]{geometry}
\usepackage[T1]{fontenc}
\usepackage{lmodern}
\usepackage{amsmath,amssymb,amsthm}
\usepackage{booktabs}
\usepackage{longtable}
\usepackage{array}
\usepackage{hyperref}
\hypersetup{bookmarksdepth=2,hidelinks}
\usepackage{enumitem}
\setlist{nosep,leftmargin=1.6em}
\setcounter{tocdepth}{1}

\newtheorem{definition}{Definition}
\newtheorem{theorem}{Theorem}
\newtheorem{proposition}[theorem]{Proposition}
\theoremstyle{remark}
\newtheorem{remark}{Remark}
\newtheorem{example}{Example}

\usepackage{xstring}
\newcommand{\lean}[1]{{\ttfamily\def\_{\textunderscore\allowbreak}\noexpandarg\StrSubstitute{#1}{.}{.\allowbreak}}}
\newcommand{\inmod}[1]{in~{\ttfamily\def\_{\textunderscore\allowbreak}#1}}
\newcommand{\Solv}{\mathsf{Solvable}}
\newcommand{\Ent}{\mathsf{Entails}}
\newcommand{\Mod}{\models}
\newcommand{\mt}[1]{{\ttfamily\frenchspacing #1}}
\newcolumntype{L}[1]{>{\raggedright\arraybackslash}p{#1}}
\emergencystretch=3em

\title{The Transform Calculus\\
\large Contextual knowledge, occurrence transactions and their interpretations}
\date{20 November 2026}
\author{}

\begin{document}
\sloppy
\maketitle

\begin{abstract}
The transform calculus gives communication, constraint solving and resource updates a common
operational step. A firing opens a rule, selects occurrences, and adds matching equations and
tells to its contextual knowledge. Admission requires a solution and asks the rule's guards at
every solution. An admitted firing publishes its knowledge and occurrence edits together.
We develop a design in which queries declare what they observe, simulations can run privately
on snapshots, and a live commit can validate a stronger claim such as \emph{this worker is still
best}. Grades can describe or weight executions, while selection and probability laws require
an explicit interpretation of the enabled alternatives. Scope, occurrence multiplicity and
continuation observations give concrete tests for MeTTa profiles and for comparisons with
reflective communication and constraint rewriting. Footnotes identify established Lean results
within this design. The research agenda gives proof obligations and distinguishing examples for
the remaining theory.
\end{abstract}

\begingroup\raggedright
\paragraph{Conventions.} Lean declarations are cited by name. Unless another namespace is given,
names are relative to
\lean{Mettapedia.Languages.ProcessCalculi.\allowbreak ConstraintMultisetRewriting}. A rule's
\emph{patterns} say what a firing takes, reads and makes, what it tells the knowledge, and what it
guards; a state's present occurrences form its \emph{space}, and occurrences may carry a
\emph{handle} naming the space they belong to. Programs are written in MeTTa notation: variables
\mt{\$x}, program equations \mt{(= lhs rhs)} as schemas, typing \mt{(: \$x A)}; equations of the
knowledge are written $\doteq$. Estimates in the qualification record are historical.
The Lean audits use only the standard axioms \lean{propext}, \lean{Classical.choice} and
\lean{Quot.sound}, or fewer (Section~\ref{sec:receipts}).\par
\endgroup

\tableofcontents
\clearpage

\section{States and firings}\label{sec:design-state}

A process may wait for a provable condition, use several resources together, or emit new
information that another process can use. A transform gives these actions a common transition
rule. Fix a frame action with terms $T_\Gamma$, equality $\approx_\Gamma$, and substitution
$f_*:T_\Gamma\to T_\Delta$ along $f:\Gamma\to\Delta$. Substitution respects binding, composition
and equality. A state $s=\langle\Gamma;K;W\rangle$ carries a store $K$ of equations and a finite
inventory $W$ of labelled occurrences. A solution of $K$ is an assignment into a frame at which
the transported equations hold. The same label can occur under two identities and be used
twice. Observing just the label's presence forgets that distinction.

\paragraph{Binding and equality.}
Each opening of a rule or stored schema gives its private variables fresh identities.
Occurrences of a private variable within the opening still denote one variable. A schema may
also capture names in the caller's frame and retain their identities. Thus two reads of a
schema for $p(x,x,y)$ can use $p(x_1,x_1,y)$ and $p(x_2,x_2,y)$, with distinct private variables
and the same captured $y$. Binding a capture then constrains subsequent work in the world
that adopts the binding. Each alternative keeps its own knowledge until a declared
publication.\footnote{\lean{ContextualStoredRows.viewAt} gives a common opening for schemas and
shared values. \lean{ContextualViewNaturality.viewAt\_transport} commutes opening with
transport of captures. \lean{ContextualStoredRowCompilation.compile\_openAt} checks the compiled
slot representation for conditional rows.}
An authored equation \mt{(= lhs rhs)} supplies a rewrite schema. Matching its left side
contributes equations about terms in the frame. The frame can interpret these equations
syntactically, by a declared equational theory, or by typed conversion. A directed rewrite
can change a term even when the two terms are distinguishable by equality.

\paragraph{Admission and publication.}
A candidate firing $f$ opens its schema at an extension of $\Gamma$. It selects present
occurrences for its takes and reads, reserves fresh identities for its makes, and forms $I_f$
by adjoining the match equations and tells to the transported $K$. For a well-shaped,
enabled selection, admission and the successor have the form
\[
 \operatorname{Sol}(I_f)\ne\varnothing,\qquad
 \forall \sigma\in\operatorname{Sol}(I_f),\ H_f(W,\sigma),
 \qquad
 (\Gamma',\,I_f,\,(W\setminus\mathrm{take}_f)\cup\mathrm{make}_f).
\]
Here $H_f$ is the opened guard and query at the cut given by the old inventory and the
firing's information. A read leaves its occurrence available. Each take spends one
occurrence. An admitted firing publishes one successor including the new knowledge and the
row edits. In a versioned realization, validation and publication must agree on the same
pre-state, including the views used by the guards. Spending a selected occurrence or
invalidating a required observation between preparation and commit defeats that
preparation.\footnote{\lean{RuleSyntax.admits\_iff} and \lean{RuleSyntax.result\_knowledge}
give the semantic rule of Definition~\ref{def:program-step}.
\lean{Mettapedia.Machines.NativeStoreTransactions.revision\_validation\_refines\_atomic\_match}
checks an occurrence-store model of validation.}

Information is conjoined within a firing. Match alternatives ordinarily explore different
extensions of the same knowledge. A solver can supply a principal answer, equivalent residual
equations, or a refutation. A budget-limited attempt can leave work pending. This gives an
evaluator distinct outcomes for a completed empty search, a blocked guard, and a resource
limit. It also lets an index propose candidates independently of the solver, provided the
candidate set covers every match of the declared theory. An index using raw heads, for
example, may miss a match when the theory identifies two different heads.\footnote{The
structural pattern fragment has a terminating decision with a principal answer or refutation
at sufficient fuel by \lean{ContextualPatternTermination.solves\_or\_refutes}.
\lean{Mettapedia.GSLT.Contracts.EquationalIndex.rigid\_exact\_iff} characterizes exact rigid
buckets under its feature and matching hypotheses.}

\section{Queries and simulation at commit}\label{sec:design-queries}

\subsection{What a query observes}

For an opened region $R$ and solution $\sigma$, let
$C_R(W,\sigma)$ be the set of present occurrence identities selected by $R$.
A predicate query asks for a proposition about this set, for example existence, absence,
satisfaction of a predicate by all members, or maximality of a proposed member. An aggregate
query computes a value such as its size, the bag of its labels, or a commutative fold over its
members. A firing may ask that a specified result hold at every solution of $I_f$. A
\emph{count} can be fixed across solutions even when member identities vary. An
\emph{exact collect} of identities makes the stronger claim that the whole selected set is
fixed. Two equal labels contribute two copies to a bag or two units to a count when the
occurrence identities differ.\footnote{\lean{RegionQueryBridge.region\_eq\_seen} and
\lean{RegionQueryBridge.observe\_region\_eq\_observe\_seen} transfer predicates and aggregates
through a validated exact query on a fixed region.
\lean{ContextualGuardSemantics.Ask.transport} transports a contextual guard with its captures.}

The read footprint includes the range in which a new occurrence could change the answer.
An empty lookup of idle workers has such a footprint even when it reads no worker identity.
A grounded count can be invalidated by a new matching row. With open knowledge a change
to a captured unknown can change membership without a row edit. The footprint therefore
includes the equality view, captured assignments and authority used by the query. A
region can be confined to a space, an index range or a declared family of terms, allowing
validation to use those boundaries. Unchanged input rows alone give an insufficient
certificate for predicates affected by arrivals.\footnote{\lean{RuleSyntax.admits\_of\_outsideRegions}
gives row locality under its freshness and region hypotheses. It holds at the stated
knowledge cut. An authority policy can permit inspection of the needed region while
reserving takes and makes to its owner. \lean{NamedSpaces.no\_name\_no\_observation}
checks a name-availability premise for observation in one located profile.}

\subsection{Snapshot advice, eligibility and best now}

A simulation can take a snapshot, explore private worlds and return a grade, proof,
ranking or proposed update. Its result has a meaning at that snapshot. A subsequent live
rule can use it with any of three contracts:
\begin{enumerate}
\item \emph{Snapshot advice.} Carry the result and its observation cut as data. The current
  transaction can use the result as a proposed choice, subject to its own live admission.
\item \emph{Current eligibility.} Commit the proposal if the selected occurrences still exist and
  the proposal still satisfies its own live conditions.
\item \emph{Current comparison.} Commit only if a predicate over the competing alternatives
  also holds at the publication cut. For a claim of maximality this is absence of an eligible
  candidate that beats the proposed one under the declared score and tie policy.
\end{enumerate}
Replaying a pure simulation from a cache is a fourth question about its result. Agreement on
the actual observations, including negative lookups, can certify reuse. The execution profile
can validate a covering version or a narrower certificate. Such reuse must respect any sampled
inputs that the result depends on.\footnote{\lean{Mettapedia.Machines.ReadOnlyQuery.run\_eq\_of\_agrees}
proves replay of the same result and read path for a pure adaptive query over a store of
keyed values. The publication protocol must establish the agreement used by that theorem.
\lean{TMeTTa.Frontier.PublicEdit.accepted\_exact} proves a versioned publication
correspondence for its covered regions.}

\begin{samepage}
\paragraph{Worker example.}
Suppose a simulation sees one job and an idle worker $a$ of score $5$. It recommends $a$.
A live rule takes the job occurrence and the occurrence for $a$, and makes their assignment.
Consider five possible changes before validation. Having maximal score admits any tied
worker in this example. The ranking in the last column includes every eligible occurrence.
\begin{center}
\small
\begin{tabular}{@{}L{0.31\linewidth}L{0.19\linewidth}L{0.19\linewidth}L{0.22\linewidth}@{}}
\toprule
Change after the simulation & Take $a$ if eligible & Take $a$ if maximal & Exact ranking is current \\
\midrule
Idle $b$ arrives with score $9$ & yes & no & no \\
Idle $c$ arrives with score $1$ & yes & yes & no \\
Idle $d$ arrives with score $5$ & yes & yes & no \\
An unrelated record arrives & yes & yes & yes \\
The occurrence for $a$ is spent & no & no & no \\
\bottomrule
\end{tabular}
\end{center}
\end{samepage}
After $b$ arrives the eligibility and maximality tests disagree. The selected job and worker are identical, but a test of
$\{w\in C_{\mathrm{idle}}(W,\sigma)\mid \mathrm{score}(w)>\mathrm{score}(a)\}=\varnothing$
has changed truth value. The same local grade for $a$ therefore licenses different outcomes.
A policy promising a unique winner, a particular tie-break, or a sample among tied workers
must observe the tie information its promise uses.\footnote{\lean{AnnotatedRuns.local\_stable}
and \lean{AnnotatedRuns.argmax\_changed\_by\_arrival} separate per-candidate filtering
from maximum selection. \lean{RegionQueryBridge.isBest\_region\_iff\_best\_seen}
gives the maximum test using a validated exact region.}

The communication between simulation and transaction is consequently a proposal with an
explicit validation contract. A simulation can be expensive and happen away from the live
commit. The live rule can use a short certificate of the comparison, or perform the predicate
query in the same firing as its takes. Checking selected row versions alone accepts the
late-$b$ control. Checking an entire space revision can reject the unrelated arrival.
Coordination is needed over whatever can falsify the promised property. The contract leaves
room for locks, reservation and version protocols, maintained aggregate indexes, and
space-local validation. It also exposes which of these require stronger evidence when data
or constraints span spaces. This is the predicate-read problem familiar from transaction
isolation~\cite{ports-grittner-ssi}.

\section{Grades, choice and causal observations}\label{sec:design-grades}

A grade can be assigned to a firing using the data it matched and a specified view. Its
purpose determines its algebra and how it enters execution. A policy that compares candidates
first fixes a family by the matching and eligibility rules at the observation cut, then
uses the grades of that family for the comparison or sampling law.
\begin{itemize}
\item A support test or fixed threshold on a candidate's own grade is a local admission test
  once its inputs are fixed. A maximum, quantile or normalized share depends on a family of
  competitors. A firing using one of these quantities needs an observation of that family
  sufficient for its claim.
\item A cost ledger or a provenance annotation can compose along a run. For grades with a
  multiplicative composition, the grade of a run is the product of the successive grades.
  This assigns a weight to a derivation. Sampling, metering and logical permission each
  require their own interpretation.\footnote{\lean{AnnotatedRuns.siteAccount\_factors\_static}
  gives descent of a commutative site annotation through its specified interchanges.
  \lean{AnnotatedRuns.meteredRun\_iff} accounts for an explicit purse.}
\item A stochastic execution declares a law of choice. For a finite race of admitted
  firings with positive total rate, the probability of the next firing $f$ is
  $\lambda_f/\sum_g\lambda_g$. A continuous-time race with independent exponential clocks
  also uses that total rate for the holding time. A probabilistic program may instead draw
  an explicit random value, keep its source identity and branch on it. Waiting, divergence,
  an all-zero race, and unresolved scheduling choices need a meaning in the specified
  profile.\footnote{\lean{Race.prob} defines finite normalization in the module
  \lean{Step.RaceProbability}. \lean{CausalHistories.history\_factorizes} sums shuffles of
  maximal runs for a product of independent positive-rate systems.}
\end{itemize}
Each admitted firing and its chosen interpretation determine the contribution to the
law. Equal printed results can receive contributions from different firings. Evidence about
a hypothesis requires a model of the observations and their dependence. For example, two reads
of one shared fair draw agree with probability $1$, while two independent fair draws agree
with probability $1/2$. Those programs can have the same available outcomes in each read.

Runs compose as paths of admitted firings. A candidate interchange between consecutive
firings $f$ and $g$ needs both residual firings $g/f$ and $f/g$, enabled in their respective
successor states, and a comparison of the resulting contexts:
\[
 s \xrightarrow{f} s_f \xrightarrow{g/f} s_{fg}
   \ \cong\ s_{gf} \xleftarrow{f/g} s_g \xleftarrow{g} s .
\]
The comparison transports the occurrence and frame identities usable by the continuation.
It must also respect observations the client is allowed to retain. If each of two
templates takes its own row and counts the remainder, both orders can leave the store
empty. The first taken row is tagged with $1$ and the second with $0$ in each order. A predicate can spoil a residual when the
first firing changes the other's observation. Even pairwise squares can fail to assemble
into a cube: the three tells $x\doteq y$, $x\doteq 1$, and $y\doteq 2$ are pairwise consistent
over distinct $1,2$ and jointly inconsistent.\footnote{\lean{Interface.Transport.Independent.cell}
constructs the interchange cell for the stated transport and independence.
\lean{AnnotatedRuns.CellControls.hollow\_cube} gives the three-tell obstruction.}

The choice of observation determines whether to distinguish paths or use causal histories.
A sum over two linear orders of one history can count its contribution twice. An
interchange of resource updates also says nothing by itself about shared stochastic
sources. A probability calculation on histories needs both a law for their schedules and
the source-dependence information used by the sampling or evidence model.\footnote{\lean{AnnotatedRuns.GradeControls.denote\_counts\_interleavings} gives the
commuting-square example. \lean{Interface.KnowledgeTensorCore.SeparatedControls.separated\_yet\_correlated}
gives a correlation counterexample under operational separation.}

\section{Language profiles and comparison}\label{sec:design-profiles}

\paragraph{MeTTa from the step.}
A MeTTa profile assigns meaning to calls and scope, argument demand, lookup and guard
timing, firing grain, and the observation of answers and effects. Rules for calls can select
an authored definition occurrence, match its schematic left side, and continue with the
opened body, keeping any captured caller identities used there. A profile can give
different demands to arguments while sharing that core. Explicit holding and activation
let a program state when syntax is data and when it runs. If an argument has two possible
results, evaluating once and reusing it gives two correlated pairs. Evaluating it
independently at two positions gives four.\footnote{In
\lean{Mettapedia.Languages.MeTTa.ContextualTransform.TMeTTa},
\lean{annotated\_agreement} checks agreement on the explicitly annotated fragment
and \lean{activate\_held\_call} checks activation of a held nullary call.}

A collection observes an alternative search at a declared cut. Its result is a bag of
constrained answers, each with its own store and any retained captures. Reopening one answer
conjoins its constraints with the caller's. Commutative folds consume the bag directly.
An order-sensitive consumer asks for a sequence by a declared linearization. A
structural canonical order and a choice among permutations have different observations.
Multiplicity belongs to the authored occurrences and their firings. For example, two
authored copies of an equation for $q(x)$ that reads $\mathsf{edge}(x,y)$ and returns $y$,
against two copies of $\mathsf{edge}(a,b)$, give four answers $b$. With the caller's $x$
subsequently set to $z\ne a$, reflecting those answers gives none. Projecting to the
printed $b$ loses the constraint that explains this result.\footnote{\lean{TMeTTa.Query.Edges.versioned\_occurrences},
\lean{TMeTTa.Query.Edges.count\_four} and \lean{TMeTTa.Query.Edges.reflect\_count\_zero}
check this example in \lean{Mettapedia.Languages.MeTTa.ContextualTransform.TMeTTaQuery}.
\lean{TMeTTa.Frontier.runs\_iff} compares two readings of the shared scheduler with alternative
implementations of read, publication and collector close.}

The desired connection is a family of profiles with demonstrated translations.
A good translation gives repeatable contextual rules for the source operations, and
preserves and reflects the appropriate observations, including scope, multiplicity,
pending work and public effects. Profiles can differ, for example in when templates read
a world or whether their conflicting effects run in sequence. A bag of answers forgets
reporting order while retaining differences in what executed templates observed. The
frame-level fibre theorem in Section~\ref{sec:join} has a particular mathematical domain.
A fibration of language profiles would additionally require coherent lifting of the
profile changes and their observations.

\paragraph{Reflective communication and constraints.}
In Meredith and Radestock's $\rho$-calculus, quotation builds channel names from processes,
name equivalence determines synchronization, and semantic substitution explains the
activation of the received process~\cite{meredith-radestock-rho}. This is a natural
specialization of the transform idea: communication uses a selected input and output, a
channel comparison, and publication of the instantiated continuation.\footnote{\lean{FragmentsOfTheStep.ReflectiveCalculus.CleanFragment.quotientEquiv}
identifies the closed, sealed-quotation source quotient with the determinate closed-row
domain of one four-rule program, modulo its readback equivalence.
\lean{ContextualCommunicationAdequacy.Payload.authored\_bisimulation} checks a separately
declared scope-sensitive, quote-safe source profile.}
Rholang and F1R3Lang profiles can add joins, peeks, persistent resources, and \emph{where}
conditions. These are direct comparison points. A guard on one candidate tuple, a proof
that no other tuple meets a condition, and a normalized quantitative selection each have
a different observation contract. The pinned F1R3Lang compiler draft gives the hardware
image the support of graded clauses and sends graded programs to the simulator
\cite[decisions 19--22]{f1r3lang-compiler-draft}. The worker example states what a
further promise of live quantitative selection would require.

Concurrent constraint programming gives a close account of ask and tell over partial
information~\cite{saraswat-rinard-ccp}. CHR's multiheaded simplification, propagation and
simpagation supply a comparison for taking, keeping and producing constraints
\cite{fruehwirth-chr}. MM2's sources and sinks distinguish work during matching from
work during rewriting~\cite{mork-sources-sinks}. A comparison of concurrent updates
must specify the view seen by each source and the joint effects of its sinks.
The useful mathematical comparisons specify the store discipline,
constraint theory, guard and reading rules, and observations of a surrounding program.
Even without predicate queries, solutions of shared open equations and atomic
multi-occurrence interactions give tests beyond simple binary communication. Theorems
about expressiveness or minimality must state the allowed encodings, composition,
observations and cost or atomicity criterion.\footnote{\lean{NullsOnTheStep.ParallelClosed.no\_uniformTranslation\_presentation}
gives an absence separation for uniform parallel translations with its stated success
observation and parallel-context premise. Section~\ref{sec:addendum-separation} gives the
criterion and its scope.}

\clearpage
\appendix

\section{Research agenda}\label{sec:design-agenda}

\paragraph{Contextual equality and solver portfolios.}
State an interchange contract for frame-specific theories in terms of the solutions of
each query, including the projection of private helpers and conversion beneath binders.
Prove that a regional solver returns an equivalent solved or residual problem, a
refutation, or a pending attempt. Identify the hypotheses for combining regions with
shared variables. Check reordered distinct-variable spines, a repeated spine that
remains pending in the pattern profile, and an equation that crosses a dependent frame
boundary. Require an index to preserve all admissible occurrences even when a tell
changes equivalence classes or a capture changes the lookup key.

\paragraph{Validation of predicates and aggregates.}
Lift the versioned publication result to dependent query views and certificates of the
exact claim used at commit. Instantiate it first with the worker example, comparing
a snapshot ranking, eligibility, and live maximality. Check late better, lower, tied
and unrelated arrivals, a spent selected identity, and a score changed by a shared
binding. Derive a sufficient validation footprint for each contract before comparing
coarse revisions with predicate-aware protocols. Extend the same argument to an
aggregate whose value is stable while its contributing identities change.

\paragraph{Compositional MeTTa profiles.}
Give the demand, effect and observation choices one formal interface. Prove translations
from independently stated source judgments for scoped calls, held code, answer
collection and reflective activation. Close a nested collector only after every child
is resolved, and exercise the suspended-child, incompatible-capture and duplicate-origin
controls. Give pending calls and wakeups an admitted-rule treatment and compare the
resulting schedules, including effects. Establish what morphisms of profiles preserve
before claiming a language-level fibration.

\paragraph{Witnessed concurrency.}
Assemble transport of actual residuals, squares and higher cells under frame change,
including fresh helper identities and observer-dependent continuation readback.
Establish conditions for composing the squares, including the joint solvability used
by a larger wave. Use both a disjoint pair and conflicting remove-and-count continuations,
then use the hollow cube to delimit the hypotheses of a higher-dimensional law.

\paragraph{Stochastic programs and inference.}
Extend the finite race and history models with explicit stochastic inputs and their
reuse. Determine when a probability law descends to a chosen quotient of runs, and
when global contention requires keeping a scheduler or event order. Start with a
finite source whose exact posterior is calculable by enumeration. Compare a second
inference route on the normalizer, evidence replay, shared versus independent draws,
and resumed work. A rule for an approximate or bounded comparison should specify the
error claim that the live commit can validate.

\paragraph{Comparative adequacy and execution.}
Compare pinned profiles of $\rho$, F1R3Lang, CHR, MM2 and the MeTTa dialects on
quote substitution, guarded joins, non-consuming reads, open knowledge, and the
concurrent worker example. Require independent source reductions and an observation
at the publication boundary. Separately refine one common prepared admission
representation to an interpreter, bytecode and eligible native routes. Compare exact
answers, effects and pending work across routes. Bound candidate and validation work
in a stated access model and measure index preparation, reuse and invalidation
separately.

\clearpage
\part*{Qualification record}
\addcontentsline{toc}{part}{Qualification record}
\addtocontents{toc}{\protect\setcounter{tocdepth}{0}}

The following sections retain the detailed qualification through 13 November 2026.
Their estimates and open rows describe those checkpoints.
Appendix~\ref{sec:design-agenda} states the current research questions in terms of
the design above. Section~\ref{sec:receipts} records the corresponding Lean audits.

\section{The question and the answer}

The question is whether one transform calculus coherently accounts for communication, querying and
execution, together with ordinary Prime programs and checked inference, and if not, why not. The
completion rule is strict: lists of results do not settle it, and an impossibility needs a proved
obstruction under explicit adequacy criteria.

\paragraph{Answer.} Retain the step. In detail:
\begin{enumerate}
\item \emph{Retain.} Every behaviour in the seven decisive tests is decided by theorems about
  steps that define their firings by one rule (Section~\ref{sec:recommendation}). Those steps are
  instances of one interface, whose own laws are only that admission implies preparation, that
  knowledge refines and that the space changes by the firing's commit; what is admitted is fixed by
  each instance (Section~\ref{sec:core}).
\item \emph{One calculus.} The rule is defined once, as the step of a program
  (\lean{RuleSyntax.step}): a program is a set of rules, a firing is a fresh opening of one of them
  with a selection, and admission is computed from the rule, its opening and the selection
  (Section~\ref{sec:core}). That step and the older steps are instances of one interface,
  \lean{Interface.Step} over a frame action \lean{Interface.Frames}: an effect core with an
  admission predicate that implies preparation. The interface is where laws are proved once. It is
  not the calculus, since every transition system whose knowledge is solvable at its sources and
  refines along its transitions is an instance (\lean{RuleSyntax.realizable\_iff}). Guard
  semantics is an optional certificate, and for guarded steps admission is characterized once
  (\lean{Interface.Step.admits\_iff}). The first-order step, the contextual steps and the
  dependent full step are guarded instances, and every other step the comparisons use is placed as
  an instance or a restriction of one, the contextual guarded firings at admission. The first-order
  and the dependent full step, each restricted to the fresh openings of a program, have the
  transition relation of that program's step, with the restrictions stated after
  Definition~\ref{def:program-step}. Each feature is proved at the interface, or on a named
  restriction with a separating control and a bridge, or named as an open gap
  (Section~\ref{sec:joined}).
\item \emph{The join.} Over a package whose names are all rigid, the second-order frames of the
  conflict-bearing steps are the $\eta$-long fibre of the frames of the dependent full step
  \lean{Interface.DependentFull.step} (\lean{ContextualDependentFibre.fibre\_full}). For firings
  whose guards and predicate queries are equations, admission transfers through the fibre wherever
  the dependent unifier decides the translated information
  (\lean{Interface.DependentFull.fibre\_admission\_transfer\_of\_full}); absence and counting
  transfer by separate theorems with their own hypotheses, and results are compared for the
  first-order instance (Section~\ref{sec:join}). Both halves of ``decides''
  are inhabited by actual kernel runs of the unifier: a solved translated pattern store whose firing
  is admitted on both sides, and a refuted one whose firing is refused on both sides
  (\lean{Interface.DependentFull.NonVacuity.solved\_case\_nonvacuous},
  \lean{Interface.DependentFull.NonVacuity.refuted\_case\_refused}).
\item \emph{Open, and one claim withdrawn.} That the dependent unifier decides every store of the
  dynamic pattern fragment is false as stated: with universes included silently, the fragment has a
  solvable store with no principal answer, and every exact solver stops there postponed
  (\lean{UnificationWithLevels.DependentFibre.cross\_postponed}). The earlier rule set fails already
  on two other stores of the fragment
  (\lean{ContextualDependentFragment.Controls.fragment\_not\_decisive},
  \lean{ContextualDependentFragment.Controls.blocked\_not\_decisive}); the current solvers keep
  every outcome of the earlier ones and decide both of those stores
  (Section~\ref{sec:open-unification}). Solvers whose steps keep every solution decide named
  fragments over packages whose names are all rigid, and postpone the rest
  (Section~\ref{sec:addendum-levels}). Everything stated ``where the unifier decides'' stays true.
\end{enumerate}
The rule is not, by itself, the whole calculus, and the reasons are proved:
\begin{enumerate}
\item logical validity is not admission: inference is an \emph{authority restriction} of the step's
  firings, and checked holes need a second, \emph{closed} ask;
\item the frames must be dependent with $\eta$: over a package whose names are all rigid, the
  second-order frames are their $\eta$-long fibre, without $\eta$ the translated pattern equations
  lose their principal answers, and principal
  answers transfer through the fibre in one direction only;
\item what the step may observe is split into ordinary queries, read through solutions, and
  reflection queries, which see syntax;
\item transport of native firings exists only up to dead names, with a move defined on part of the
  firings, and helper firings at dependent frames agree only up to the exchange of their telescopes;
\item concurrent waves need joint solvability of the whole wave;
\item atomic blocks of unbounded size need rules computed per block or a compound occurrence;
\item fresh identities are those never born, not those merely absent;
\item spaces share the knowledge, so a firing at one space can change admission at another through
  a shared unknown, and independence across spaces holds on the restriction that the firing adds no
  information; rights alone do not bound what a space receives;
\item contention between firings includes inhibition through predicate queries;
\item a guard that something bad is never reached has a finite certificate exactly on finite-state
  data, and is lowered to the check of one; reachability of a proof does not make fair search
  complete.
\end{enumerate}
Each is a revision of the candidate as first presented, forced by a proved separation
(Section~\ref{sec:recommendation}). The precisely stated stronger hypotheses are rejected there,
each with its obstruction.

\section{The core calculus}\label{sec:core}

\subsection{Syntax, frames and solutions}

Terms are built over a binding signature: sorts, operators with binding arities, bound variables, and
metavariables declared with arities $(\bar s, s)$ that are applied to terms of the sorts
$\bar s$. A \emph{frame} declares metavariables: the unknowns that exist and their sorts or types.
The canonical frames are declaration families, the foundation's frames are formed dependent
telescopes, and a substitution action has one frame.

\begin{definition}[Frames acting on terms; \lean{Interface.Frames}]
A frame action has frames $a,b,\dots$, terms $T_a$, morphisms $f : a\to b$ with an action
$f_* : T_a \to T_b$, identities and composites, and an equivalence $\approx_a$ on each $T_a$, with
$(\mathrm{id})_*t=t$, $(g\circ f)_*t = g_*f_*t$ and $s\approx t \Rightarrow f_*s\approx f_*t$.
\end{definition}

For an equation store $P$ at $a$, a \emph{solution} is a morphism $f : a\to b$ into any later frame with
$f_*s\approx f_*t$ for every $(s,t)\in P$. $\Solv(P)$ asks for one; $\Ent(P,\varphi)$ asks that
$\varphi(b,f)$ hold at every solution; $u : a\to m$ is \emph{principal} when it is a solution and
every solution factors through it (\lean{Interface.Frames.Principal}), and then a store's solutions
are exactly the factorizations through $u$ (\lean{Interface.Frames.Principal.models\_iff}).
Solvability is witnessed at a weakly terminal frame, and the frame with one parameterless
metavariable per sort is weakly terminal, so existence never needs a higher-order helper
(\lean{Interface.Frames.solvable\_iff\_into}, \lean{Interface.Frames.firstOrderFrame\_weaklyTerminal},
\lean{Interface.Frames.contextual\_solvable\_iff\_consistent}). A substitution action is the one-frame
case, whose solvability is consistency (\lean{Interface.Frames.ofSubstitution\_solvable\_iff}).

\paragraph{Closed asks.} Given a closed frame, a store is \emph{closed-solvable} when it has a solution
into it (\lean{Interface.Frames.ClosedSolvable}). Symbolic and closed solvability agree under a
principal answer whose frame has a closed point
(\lean{Interface.Frames.closedSolvable\_iff\_of\_principal}).

\subsection{States}

A state $\langle\Gamma;K;W\rangle$ has a frame $\Gamma$, its \emph{knowledge} $K$ (an equation store
at $\Gamma$) and its space $W$, a finite set of occurrence identities; each occurrence carries a
label, a term at the frame. Occurrence identities are physical: equal labels at two identities are
two occurrences. Fresh schematic variables of a rule are opened in a fresh scope at each firing;
caller variables shared with the state are not renamed; the identities a firing makes are fresh.

\subsection{The rule}

The calculus has one rule. It is given first as the step of a program, where admission is computed
from the rule. Then comes the interface that this step and the older steps implement. Then come the
older steps, whose firings carry an instance of a rule, and their agreement with the step of a
program.

\begin{definition}[The step of a program; \lean{RuleSyntax.step}]\label{def:program-step}
Fix a frame action and, for each frame $\Gamma$, its \emph{extensions} $x$, each reaching a frame
$\Gamma_x$ with a weakening $w_x:\Gamma\to\Gamma_x$ (the identity at one frame, new declarations at
declaration families, helper telescopes at typed frames). A \emph{rule} $r$ has its own frame
$\rho_r$ of schematic unknowns and, at that frame, take patterns, read patterns, told equations,
guards (equations), predicate queries (the absence of a pattern, a count of a pattern, and patterns
whose exact set of matching rows the firing claims), and made rows; in the Lean the guards, absences
and counts are a rule's checks and the last kind its census patterns. A \emph{program} $P$ is a set of rules. A state $\langle\Gamma;K;W;\ell\rangle$ has a
label $\ell(o)\in T_\Gamma$ for each occurrence. A \emph{firing} of $P$ at that state is a rule
$r\in P$, an extension $x$, an \emph{opening} $\iota:\rho_r\to\Gamma_x$, a selection of the
occurrences it takes, reads and makes, and for each pattern of the last kind a claimed set of rows. With $\mathrm{row}(o)=(w_x)_*\ell(o)$, the firing determines
\begin{itemize}
\item its information $I_f$: $(w_x)_*K$, each opened take and read pattern equated with the row
  selected for it, and the opened tells;
\item \emph{freshness} of the opening: every instantiation of $\Gamma$ and every instantiation of
  $\rho_r$ into a common frame are realized together by one instantiation of $\Gamma_x$, on the
  terms the state mentions and on the rule's terms;
\item its shapes: a fresh opening, distinct taken and distinct made identities, one row for each
  pattern, one made identity for each made row, one claimed set for each such pattern;
\item $\mathrm{Enabled}(W,f)$: the taken and the read occurrences are present, the read ones apart
  from the taken ones, and the made identities fresh;
\item $H_f(\sigma)$ at a solution $\sigma$ of $I_f$: every opened guard holds at $\sigma$, every
  opened absence and count holds there (absence: no present row is equated with the pattern by
  $\sigma$; a count: exactly that many are), and the rows that $\sigma$ equates with a pattern of the
  last kind are exactly its claimed set.
\end{itemize}
Admission is $\mathrm{Shapes}(f)$, $\mathrm{Enabled}(W,f)$, $\Solv(I_f)$ and $\Ent(I_f,H_f)$, and
the result is
$\langle\Gamma_x;\,I_f;\,(W\setminus\mathrm{taken}_f)\cup\mathrm{made}_f;\,\ell'\rangle$, where
$\ell'$ gives each made identity its opened made row and every other identity its row.
\end{definition}

Here admission is computed from the rule, the opening and the selection; it is not given with the
step (\lean{RuleSyntax.admits\_iff}), and a firing of $P$ is by construction a fresh instance of a
rule of $P$.

\begin{definition}[The step interface; \lean{Interface.Step}]\label{def:step}
A step over a frame action gives, at each state $s=\langle\Gamma;K;W\rangle$, a type of firings. A
firing $f$ has a selection (the occurrences it \emph{takes}, \emph{reads} and \emph{makes}), a
target frame $\Gamma'$ with an extension $e_f:\Gamma\to\Gamma'$ (its new unknowns), its
\emph{information} $I_f$, an equation store at $\Gamma'$ every solution of which, composed with
$e_f$, solves $K$, and its shapes, a proposition $\mathrm{Shapes}(f)$. The firing is
\emph{prepared} when $\mathrm{Shapes}(f)$, $\mathrm{Enabled}(W,f)$ and $\Solv(I_f)$ hold; a
prepared firing has a result state, whose space is $(W\setminus\mathrm{taken}_f)\cup\mathrm{made}_f$,
and a settling morphism $c_f:\Gamma'\to\Gamma''$ into the result's frame, along which every solution
of the result's knowledge solves $I_f$. \emph{Admission} is a predicate on firings, given with the
step, that implies preparation.
\end{definition}

These are the fields of \lean{Interface.Step} and of the effect core it extends
(\lean{Interface.EffectCore}). The interface has no labels, patterns or tells, and its admission is
any predicate that implies preparation; its own laws are only that knowledge refines, along the
extension and along the settling morphism, and that the space changes by the selection's commit.
\emph{Guard semantics} is an optional certificate (\lean{Interface.Step.Guarded}): a predicate
$H_f(\sigma)$ at every solution $\sigma$ of $I_f$, the firing's \emph{guards} and \emph{predicate
queries}, such that admission is preparation together with $\Ent(I_f,H_f)$. For a guarded step,
admission is therefore characterized once (\lean{Interface.Step.admits\_iff}): the shapes, the
enabling of the selection's commit at the present occurrences, semantic solvability of the
information, and the ask of $H_f$ at every solution of it. Since $H_f$ is part of the certificate,
this fixes the form of admission and not its content, which is each instance's. The result's space
is $(W\setminus\mathrm{taken})\cup\mathrm{made}$ (\lean{Interface.Step.present\_result}). On a step
whose knowledge is literal, which includes every step built with \lean{Interface.Step.Standard} and
every instance below, the result's knowledge is the information carried along the settling morphism
(\lean{Interface.Step.settle}, \lean{Interface.Step.knowledge\_result}).

\paragraph{What the interface is, and what the rule excludes.} A relation on states with given
frames, knowledge and present occurrences is the admitted relation of an instance of
\lean{Interface.Step} exactly when every transition starts from solvable knowledge and its target's
knowledge refines the source's along some morphism of frames (\lean{RuleSyntax.realizable\_iff},
\lean{RuleSyntax.admitted\_laws}); with no knowledge every relation is one
(\lean{RuleSyntax.ofLTS\_admitted\_iff}). A theorem proved at the interface is therefore a theorem
about all such systems. The step of a program excludes more. Its admission does not change when the
present rows change outside the firing's take, read, make and query regions, provided the opening
stays fresh (\lean{RuleSyntax.admits\_of\_outsideRegions}); a row inside an absence region refuses
the firing (\lean{RuleSyntax.Controls.row\_in\_region\_refused}), and there is an instance of the
interface whose admission depends on a row outside every region
(\lean{RuleSyntax.Controls.nonlocal\_restriction}). This locality is in the rows only: the
regions are read at every solution of the whole knowledge, and the knowledge of a result is, by
definition, the firing's information, that is the state's knowledge along the extension, the
equations of the opened patterns with the selected rows, and the opened tells
(\lean{RuleSyntax.result\_knowledge}). The rule's terms are those of the frame action, whose
equivalence may itself carry an equational theory; a placement can put computation there, as the
placement of rho in Section~\ref{sec:addendum-rho} does.

\begin{definition}[The rule in the older steps]\label{def:rule}
In the first-order step \lean{AnnotatedRuns.firstOrder}, the native step \lean{Interface.native}
and the dependent full step \lean{Interface.DependentFull.step}, a firing is an instance of a rule
with take, read and make patterns, told equations, guards and predicate queries, at an extension of
the state's frame, together with a selection. Its information contains $(e_f)_*K$ together with the
equations matching its take and read patterns to the selected labels and the equations it
\emph{tells}, and $H_f(\sigma)$ asks its guards and predicate queries at $\sigma$. Each of these
steps is guarded and has literal knowledge, so its one rule is
\[
\frac{\mathrm{Shapes}(f)\qquad \mathrm{Enabled}(W,f)\qquad \Solv(I_f)\qquad \Ent(I_f,H_f)}
{\langle\Gamma;K;W\rangle\ \xrightarrow{\ f\ }\ \langle\Gamma'';\,(c_f)_*I_f;\,(W\setminus\mathrm{taken}_f)\cup\mathrm{made}_f\rangle}
\ \textup{[\textsc{Fire}]}
\]
where $c_f:\Gamma'\to\Gamma''$ is the firing's settling morphism.
\end{definition}

In the first-order and the dependent full step a firing carries its rule instance, and no program
restricts which instances are offered. The first-order step, restricted to the firings that open a rule of a program freshly,
has the same transition relation as the step of that program at one frame, restricted to the
firings whose made identities already carry their made rows
(\lean{RuleSyntax.OneFrame.firstOrder\_eq}; the first-order step fixes the label of every identity
in advance). The dependent full step, restricted in the same way, has the transition relation of
the step of the program at typed frames with helper telescopes
(\lean{RuleSyntax.Dependent.dependent\_eq}), and a program of two rules whose takes have types that
depend on terms has a run there (\lean{RuleSyntax.Typing.Programs.dependent\_run}). Unrestricted, both steps also offer rule instances that mention the state's own
unknowns. Two differences from the syntax are proved. A predicate query of the dependent step may
have any region and any answer that read through solutions, where the syntax has pattern regions
answered by none, by a number, or by an exact set. And the helper step of the contextual layer
relates a fresh helper, without reading a row, to an old name that the state mentions
(\lean{RuleSyntax.Contextual.CouplingControls.mentionedHelper\_coupled}, from
\lean{RuleSyntax.Contextual.helper\_coupled}), which no firing with a fresh opening that takes and
reads no row does (\lean{RuleSyntax.information\_amalgamates}). The restriction to names the state mentions is
needed: at a state that mentions nothing, the step of a one-rule program admits a firing that relates
a new name to the old one
(\lean{RuleSyntax.Contextual.CouplingControls.coupling\_admitted},
\lean{RuleSyntax.Contextual.CouplingControls.coupling\_couples}), and refuses it once the knowledge
mentions that name
(\lean{RuleSyntax.Contextual.CouplingControls.coupling\_refused\_when\_mentioned}). No theorem
says that the helper step is the step of no program.

The settling morphism is the identity on the dependent full step and an
inclusion on the native step. Whether settling keeps solvable information solvable is a named
property, \lean{EvidenceInterface.SettlesConservatively}, not part of the rule.

Enabling requires the taken occurrences to be present, the read ones present and apart from the
taken ones, and the made ones fresh. The shapes are each instance's own: the number of patterns, the
distinctness of takes and makes, and the labels of what is made. Two read positions may select one
occurrence.

The guard is asked at every solution of the firing's \emph{own} information, which includes its own
tells and its own matches: that is the \emph{birth/guard cut}. A predicate query fixes the
complete set of present occurrences in a region at every solution; absence is the empty case.

\paragraph{Runs and the cut.} Runs are dependent words of admitted firings
(\lean{Interface.Step.Run}). Knowledge only refines along them: every solution of the last
knowledge pulls back to a solution of the first (\lean{Interface.Step.Run.models\_start}). With
unknowns read as proof obligations and solutions as proofs, a run's final solution is a proof of the
start's obligations. On a guarded step, what an admitted firing's guards and predicate queries
answered at its cut holds at every later solution, read back to its target frame
(\lean{InterfaceTable.holds\_of\_admitted}, \lean{InterfaceTable.cut\_preserved}).

\subsection{Identity of states}

Three identities of states are used, each with its theorem.
\begin{itemize}
\item \emph{Renaming of occurrences and equivalence modulo the knowledge}
  (\lean{Observation.ModuloEquivalent}): a bisimulation of the ordinary step that offers every
  firing under its rule, given a fresh supply of identities for every label
  (\lean{Observation.ModuloEquivalent.semantic\_bisimilar}); on configurations with consistent
  knowledge, and given that supply, it is exactly the bisimilarity of that step
  (\lean{LogicOverTheStep.ProgramFreeBisimilarity.ruleOffers\_bisimilar\_iff},
  Section~\ref{sec:addendum-program-bisimilarity}). Queries that read through solutions
  keep it; raw inspection of quoted syntax does not (\lean{Program.rawQuoted\_not\_semantic}).
\item \emph{Up to dead names} on the native step: states that agree except on names no row,
  knowledge or ledger entry mentions are bisimilar
  (\lean{Interface.NativeUpToDead.bisimilar\_of\_deadEquiv'}).
\item \emph{The races layer}: strong, weak and branching bisimilarity over any instance and any
  transport (\lean{Race.Generic.Bisimilar}, \lean{Race.Generic.WeakBisimilar},
  \lean{Race.Generic.BranchingBisimilar}); transport-equivalent states are bisimilar
  (\lean{Race.Generic.bisimilar\_of\_same}). At dependent frames, firings that make helpers are alike
  up to an isomorphism of telescopes (\lean{Interface.DependentFull.HelperAlike}).
\end{itemize}

\subsection{The fibre tower}

\begin{center}\small
\begin{longtable}{@{}L{0.20\textwidth}L{0.74\textwidth}@{}}
\toprule
\textbf{Level} & \textbf{Frames and the theorem that relates it to the next} \\
\midrule
\endhead
one frame & a substitution action; consistency, entailment, principality (\lean{Interface.Frames.ofSubstitution\_solvable\_iff}) \\
first order & parameterless metavariables of declaration families; weakly terminal, decides existence (\lean{Interface.Frames.contextual\_solvable\_iff\_firstOrder}); separately, first-order terms with their own unifier, which decides (\lean{FirstOrder.accepted\_iff\_consistent}, \lean{FirstOrder.accepted\_principal}); no frame map between the two presentations is proved \\
patterns & metavariables applied to distinct bound variables; the driver solves or refutes (\lean{ContextualPatternTermination.solves\_or\_refutes}); elsewhere residuals stay and lose no solution (\lean{ContextualFrameWorklist.preserves\_every\_solution}, \lean{BetaEtaReadings.FlexFlex.never\_refuted}) \\
second order & declaration families; raised into dependent telescopes, with reflection of solutions and asks (\lean{Interface.Translation.translation\_reflects}) \\
dependent, with $\eta$ & formed telescopes with typed $\beta\eta$; over a package whose names are all rigid, the second-order frames are its $\eta$-long fibre on translated frames (\lean{ContextualDependentFibre.fibre\_full}, \lean{ContextualDependentFibre.fibre\_solutions\_correspond}); a dependent principal answer of a translated store at a translated assignment gives a second-order principal answer, one way only (\lean{ContextualDependentFibre.fibre\_principal\_reflect}); a pattern equation has a principal answer (\lean{ContextualDependentGeneralPattern.Pattern.replace\_principal}); the driver returns principal answers when it solves, refutes only with a certificate, and never refutes a solvable store (\lean{ContextualDependentWorklist.drive\_solved\_principal}, \lean{ContextualDependentWorklist.drive\_not\_refuted\_of\_solvable}); decisiveness on the whole dynamic pattern fragment is open (Section~\ref{sec:open-unification}) \\
\bottomrule
\end{longtable}
\end{center}

The tower ends at the dependent frames with $\eta$ because principality of pattern answers comes from
$\eta$: over the dependent frames without $\eta$ the translated pattern equation $m(x)\doteq n(x)$ has
no principal answer (\lean{Interface.Translation.pattern\_no\_principal},
\lean{ContextualDependentFibre.Controls.point\_pattern\_no\_principal}), and with $\eta$ it has one
(\lean{ContextualDependentFibre.Controls.point\_pattern\_principal},
\lean{Interface.Translation.EtaFrames.eta\_pattern\_principal}). Conservativity stops at frames the
second-order side cannot express: at a frame with a $\Sigma$-typed variable a projection solves a
translated store and is no translation (\lean{ContextualDependentFibre.Controls.offImage\_solution}).
That a second-order principal answer stays principal against solutions into arbitrary dependent
telescopes, with $\Pi$, $\Sigma$ or identity types outside the image, is not proved; admission does
not need it where the unifier decides, since the answers of the solvers in the tree land in
translated telescopes (Section~\ref{sec:join}).

\section{One calculus: joined qualification}\label{sec:joined}

The results we use come from several developments. They are about one calculus only if each
step is an instance of Definition~\ref{def:step}, each feature is proved at the interface or on a
named restriction, and two definitions of one concept agree or are proved to differ.

\subsection{\texorpdfstring{The feature $\times$ instance table}{The feature x instance table}}

The columns are the interface, the first-order step \lean{AnnotatedRuns.firstOrder}, the
contextual steps over declaration frames (the helper step \lean{Interface.Helpers.helperStep}, with
the native step \lean{Interface.native} where the helper step lacks a feature), and the dependent
full step \lean{Interface.DependentFull.step}. In the three instance columns, unqualified names are
in \lean{InterfaceTable}. A \emph{gap} is open.

\begin{center}\scriptsize\setlength{\tabcolsep}{3pt}
\begin{longtable}{@{}L{0.11\textwidth}L{0.25\textwidth}L{0.20\textwidth}L{0.18\textwidth}L{0.20\textwidth}@{}}
\toprule
\textbf{Feature} & \textbf{Interface} & \textbf{First order} & \textbf{Contextual} & \textbf{Dependent full} \\
\midrule
\endhead
frames and solutions & \lean{Interface.Frames.Principal.models\_iff} & \lean{frames\_firstOrder} & \lean{frames\_contextual} & \lean{frames\_dependent} \\
admission & \lean{Interface.Step.admits\_iff} & \lean{admission\_firstOrder} & \lean{admission\_helper}, \lean{admission\_native} & \lean{admission\_dependent} \\
runs & \lean{Interface.Step.Run.models\_start} & \lean{runs\_firstOrder} & \lean{runs\_helper}, \lean{runs\_native} & \lean{runs\_dependent} \\
the cut & \lean{InterfaceTable.cut\_preserved} & \lean{cut\_firstOrder} & \lean{cut\_helper}, \lean{cut\_native} & \lean{cut\_dependent} \\
guards & \lean{InterfaceTable.holds\_of\_admitted} & \lean{guards\_firstOrder} & \lean{guards\_helper} & \lean{guards\_dependent} \\
predicate queries & absence: \lean{Interface.DependentFull.entails\_absent\_iff}; count: \lean{Interface.Frames.Principal.entails\_count\_iff} & \lean{absence\_firstOrder}, \lean{count\_firstOrder} & \lean{absence\_contextual}, \lean{count\_contextual} at frames; the helper step has none; the native step asks them (\lean{admission\_native}) & \lean{absence\_dependent}, \lean{count\_dependent} \\
spaces & \lean{Spaces.run\_executable}, \lean{Spaces.present\_result\_spaces}, \lean{Spaces.fiber\_unchanged}; commuting only on a restriction & \lean{SpacesInterface.located\_independent} & \lean{spaces\_helper}, with no restriction & \lean{SpacesInterface.dependent\_independent}; gap: query regions are not located \\
permissions & \lean{Permissions.refused\_without\_receive}, \lean{PermissionsInterface.witness\_present}; \lean{PermissionsInterface.absence\_preserved} on a restriction & \lean{PermissionsInterface.old\_owned\_absence\_preserved} & \lean{witness\_helper}; native absence: gap & \lean{PermissionsInterface.dependent\_absence\_preserved} \\
disequations & \lean{DisequalityFrames.frames}; \lean{Interface.Step.withDisequations} for steps whose prepared settlement reflects equality & \lean{DisequalityFrames.oneFrame\_admits\_iff} & \lean{diseq\_helper}; \lean{diseq\_native} for the native step with disequations, guards kept (\lean{PreparedStep.nativeGuarded}) & \lean{diseq\_dependent} \\
evidence rows & \lean{Evidence.not\_admits\_of\_not\_solvable}; \lean{EvidenceInterface.run\_solvable} on a restriction & \lean{EvidenceInterface.old\_firstOrder\_run\_consistent}, \lean{Evidence.Family.profiles\_distinct} & \lean{evidence\_helper}, \lean{evidence\_native}; the four-way family: gap & \lean{EvidenceInterface.dependent\_run\_solvable}, \lean{EvidenceInterface.DependentFamily.profiles\_distinct} \\
races and identity & \lean{Race.Generic.bisimilar\_of\_same} & \lean{races\_firstOrder} & \lean{races\_helper} & \lean{Interface.DependentFull.plain\_bisimilar\_of\_same}, \lean{Interface.DependentFull.helper\_bisimilar\_of\_same} \\
programs & gap: one program core over the interface & the first-order core (Section~\ref{sec:programs}) & the binder core; joined with the first-order core by agreement & gap \\
\bottomrule
\end{longtable}
\end{center}

\subsection{Features on a restriction}

For each feature that does not hold on every instance, the restriction theorem, a control that fails
off the restriction, and the bridge by which an instance meets it.

\begin{center}\footnotesize
\begin{longtable}{@{}L{0.14\textwidth}L{0.26\textwidth}L{0.24\textwidth}L{0.28\textwidth}@{}}
\toprule
\textbf{Feature} & \textbf{Restriction} & \textbf{Control off it} & \textbf{Bridge} \\
\midrule
\endhead
spaces commute & \emph{quiet}: the firing adds no information, its information being equivalent to the knowledge (\lean{SpacesInterface.located\_independent}, \lean{SpacesInterface.dependent\_independent}) & a shared unknown carries a binding across spaces (\lean{Spaces.Controls.shared\_variable\_crosses\_spaces}, \lean{SpacesInterface.Controls.dependent\_shared\_variable\_crosses\_spaces}) & a syntactic check implies quiet (\lean{SpacesInterface.quiet\_of\_exact}, \lean{SpacesInterface.exact\_independent}, \lean{SpacesInterface.Controls.quiet\_read\_exact}) \\
absence under permissions & labels move along firings, \lean{PermissionsInterface.LabelLaw} (\lean{PermissionsInterface.absence\_preserved}) & a firing that relabels a row spoils the absence (\lean{PermissionsInterface.Controls.relabel\_spoils\_absence}) & \lean{PermissionsInterface.locatedLabels\_law}, \lean{PermissionsInterface.dependentLabels\_law} \\
disequations & a step whose prepared settlement reflects equality (\lean{Interface.Step.withDisequations}) & a step that settles by guessing passes every admission test and holds $0\neq 0$ (\lean{DisequalityFrames.Restriction.guess\_breaks\_frame}, \lean{DisequalityFrames.Restriction.guess\_not\_reflects}) & \lean{DisequalityFrames.firstOrder\_reflects}, \lean{DisequalityFrames.Dependent.reflects}; on the native step, a firing without checks and queries is admitted with disequations exactly when it is ready and a solution of its information separates them (\lean{PreparedStep.nativeGuarded\_admits\_iff\_of\_no\_guards}) \\
no inconsistent store & \lean{EvidenceInterface.SettlesConservatively} (\lean{EvidenceInterface.run\_solvable}); a retraction suffices (\lean{EvidenceInterface.conservative\_of\_retracts}) & a step that settles by guessing enters $0\doteq 1$ (\lean{EvidenceInterface.Controls.guess\_enters\_inconsistency}, \lean{EvidenceInterface.Controls.guess\_not\_retracts}) & \lean{EvidenceInterface.firstOrder\_retracts}, \lean{EvidenceInterface.helper\_retracts}, \lean{EvidenceInterface.dependent\_retracts}, \lean{EvidenceInterface.withDisequations\_retracts}; native: \lean{EvidenceInterface.Native.native\_conservative} \\
\bottomrule
\end{longtable}
\end{center}

There is no single generic theorem for commuting spaces: a locality hypothesis general enough to
state one would restate the residual law and be met by one instance only, so independence is proved
per instance into the interface's own notion of independence.

\subsection{The join through the fibre}\label{sec:join}

\paragraph{The dependent full step.} \lean{Interface.DependentFull.step} is an instance over the
typed $\beta\eta$ frames. A firing makes a telescope of fresh helpers, has take and read patterns
and made rows at the extended frame, tells and guards equations of typed terms, and asks predicate
queries that read through solutions; absence and counting are instances. Its admission is the
interface's (\lean{Interface.DependentFull.admits\_iff}); its takes really constrain
(\lean{Interface.DependentFull.Firing.admitted\_take\_unifies}); after a firing every row is typed,
a made row at its checked type and every other row as the old label carried along the weakening
(\lean{Interface.DependentFull.Firing.subject\_reduction}). Told equations between different types
are heterogeneous, and are refused only through solvability
(\lean{Interface.DependentFull.Controls.tell\_heterogeneous}); an ill-typed make is refused
(\lean{Interface.DependentFull.Controls.illTyped\_refused}).

\paragraph{Admission transfer.} For candidate steps over the second-order frames whose guards and
predicate queries are equations, admission equals the dependent admission of the translated data
whenever the unifier decides the translated information inside the fibre
(\lean{Interface.DependentFull.fibre\_admission\_transfer}), under the package's standard facts, a
closed realization, injective names, and the hypothesis that every name of the package is rigid, so
that no name computes in conversion; the admission theorems below assume the same. That condition
is derived, not assumed: the full solver
deciding the translated store at some budget gives it
(\lean{Interface.DependentFull.decidedInFibre\_of\_full},
\lean{Interface.DependentFull.fibre\_admission\_transfer\_of\_full}). Absence and counting go through
separate theorems with their own hypotheses. For the first-order step at the second-order
endomorphism action, admission without census queries is translated admission when the unifier
decides the problem and every absence extension by an endomorphic answer or a refutation
(\lean{Interface.DependentFull.admits\_translate\_iff}), with equal results
(\lean{Interface.DependentFull.result\_translate}); counting checks transfer under the same
condition for each extension (\lean{Interface.DependentFull.admits\_translate\_iff\_count}).
Census queries with arbitrary regions are not translated: a second-order region has no meaning at a
dependent solution outside the image. A second-order step over several frames with equation guards
needs no endomorphism condition (\lean{Interface.DependentFull.translateEquation\_admits\_iff\_of\_full}).

\paragraph{Solvers that keep translated telescopes.} Every rule of the pattern and full solvers
keeps the frame, removes the newest metavariable, or replaces two metavariables by a pruned one,
and each keeps translated telescopes (\lean{Interface.DependentFull.full\_keeps\_translated}). The
later solvers do too: rotation only permutes chains of sorts, and preparation normalizes what is
already normal under the checkable condition that the root evaluator is silent on chains of sorts,
which the evaluator without root steps meets (\lean{Interface.DependentFull.normalizing\_keeps\_translated},
\lean{Interface.DependentFull.permuting\_keeps\_translated},
\lean{Interface.DependentFull.fragment\_keeps\_translated},
\lean{Interface.DependentFull.identity\_keeps\_translated},
\lean{Interface.DependentFull.pi\_keeps\_translated}, \lean{Interface.DependentFull.chainsSilent\_none},
\lean{Interface.DependentFull.decidedInFibre\_of\_keeps}).

\paragraph{Non-vacuity.} In the package with one sort constant, the actual full solver, evaluated by
the kernel, solves the translated $m(x)\doteq n(x)$ with a principal answer at a translated
telescope of one metavariable and refutes $F(x)\doteq y$
(\lean{Interface.DependentFull.NonVacuity.pattern\_commit},
\lean{Interface.DependentFull.NonVacuity.take\_transfer}). A second-order firing whose information is
the first store is admitted on both sides
(\lean{Interface.DependentFull.NonVacuity.solved\_case\_nonvacuous}); the firing for the second is
refused on both sides, with an independent second-order proof
(\lean{Interface.DependentFull.NonVacuity.refuted\_case\_refused},
\lean{Interface.DependentFull.NonVacuity.outside\_unsolvable\_direct}). Off the image, a firing with a
$\Sigma$-typed helper is admitted and its result is outside the fibre, without disturbing admission
on it (\lean{Interface.DependentFull.FibreControls.pair\_admitted},
\lean{Interface.DependentFull.FibreControls.pair\_result\_not\_translated}).

\paragraph{Concurrency at dependent frames.} Two helper-free firings are independent when they take
disjoint occurrences, neither takes what the other reads, they make distinct occurrences, their joint
information is solvable, and neither writes into the other's query regions; then both orders are
admitted and the composites have the same occurrences and labels, with equivalent knowledge
(\lean{Interface.DependentFull.local\_diamond}), and every clause
but the query clause is necessary (\lean{Interface.DependentFull.clauses\_of\_residuals}). One-helper
firings close up to the swap of their helpers (\lean{Interface.DependentFull.helper\_diamond}), a
helper firing and a helper-free one close on the nose (\lean{Interface.DependentFull.mixed\_diamond}),
and firings that make whole telescopes close up to the block exchange, the unique map between the two
pullbacks (\lean{Interface.DependentFull.telescope\_diamond},
\lean{Interface.DependentFull.exchange\_unique}). Without the exchange the composites differ
(\lean{Interface.DependentFull.TelescopeControls.composites\_differ}); a made row whose type mentions
the other firing's helper is the residual of no firing
(\lean{Interface.DependentFull.no\_residual\_of\_mentions}). Helper firings are placed in the races
layer up to an isomorphism of telescopes (\lean{Interface.DependentFull.helper\_bisimilar\_of\_same},
\lean{Interface.DependentFull.helper\_branchingBisimilar\_of\_same}).

\paragraph{Open in the join.} Census queries through the join (60\%); the endomorphism condition of
the first-order instance from the solver's answers, which needs inhabitants of dropped metavariables'
types (45\%); absence and counting through the multi-frame instance (65\%); step bisimilarity for
helper firings, which needs a strict transport (60\%); the converse of dead-helper simulation, which
needs strengthening of typed judgments (45\%).

\subsection{Other steps as instances}

\begin{center}\small
\begin{longtable}{@{}L{0.24\textwidth}L{0.70\textwidth}@{}}
\toprule
\textbf{Step} & \textbf{Placing theorem} \\
\midrule
\endhead
occurrence core & a core step is an admitted query-free firing of the first-order step, and core reachability is exactly a query-free run (\lean{OccurrenceCorePlacement.core\_step\_iff\_admitted}, \lean{OccurrenceCorePlacement.coreReaches\_iff\_run}); with predicate queries, \lean{OccurrenceCorePlacement.query\_step\_iff\_admitted}. This places the resource, census, snapshot, Need-machine and register-comparison results, each over its own substitution action \\
compiled clauses & a compiled clause is admitted exactly when its constructed firing is admitted by the native step (\lean{OccurrenceCorePlacement.Native.admitted\_iff\_native}) \\
contextual guarded firings & a guarded firing at an execution endpoint is admitted exactly when its query-free native firing is, at every budget (\lean{OccurrenceCorePlacement.Native.guarded\_valid\_iff\_native}); this covers the contextual $\lambda$, pair, causal and schematic-storage executions at admission. Their results are not compared with the native published state \\
live native step & \lean{Interface.NativeUpToDead.live}, the native step on firings with pattern queries, live outputs and no recipes \\
dependent and simple fibre & \lean{Interface.Dependent.helperStep}; the simple-telescope fibre computes the full instance's admission when the knowledge has a principal answer at a simple frame (\lean{Interface.Dependent.simpleStep\_admits\_iff}) \\
disequation steps & \lean{Disequality.step}, an instance over the disequation frames of a substitution action (\lean{Disequality.admits\_iff\_frames}); over any frames, \lean{Interface.Step.withDisequations}, for steps whose prepared settlement reflects equality (an instance: \lean{PreparedStep.nativeGuarded}) \\
ordinary step & \lean{Semantic.semanticStep}, a restriction of the first-order step (\lean{Semantic.semanticStep\_admits\_iff}) \\
program core & translated runs are runs of the ordinary step, exactly (\lean{Program.program\_interface\_iff}); the scope of fresh rule variables is threaded by the translated run, not by the frame extension \\
binder program core & each translated equation firing is a step of the contextual guarded step (\lean{ProgramBinders.destination\_is\_core\_step}), and the unmatched return is a guarded firing (\lean{ProgramBinders.unmatched\_is\_guarded\_step}) \\
typing & \lean{TypingChecking.goalStep}, \lean{TypingStates.computeStep}, \lean{TypingCheckingMeta.metaGoalStep}: instances by construction; call contracts are first-order firings (\lean{TypingFirings.closed\_call\_admits\_iff}) \\
inference hosts & a step of a schema process is exactly an admitted query-free first-order firing whose rule is an installed schema's activation and whose made occurrences the process may publish to (\lean{OccurrenceCorePlacement.schema\_step\_iff\_admitted}, \lean{OccurrenceCorePlacement.rule\_step\_iff\_admitted}) \\
restrictions by rights & any instance restricted by a predicate on firings is again an instance (\lean{Spaces.restrict\_admits\_iff}) \\
MM2 \mt{exec} & a covered \mt{exec} over a duplicate-free ground representation is one firing (\lean{MeTTaOperationsExec.exec\_is\_one\_firing}) \\
outside & the abstract machines of the Need lane (for instance \lean{Mettapedia.Machines.BranchLocalNeed.CallTimeChoiceInstructions.source\_iff\_executable}) and the register semantics of Meta-MeTTa are comparison sources, not instances \\
\bottomrule
\end{longtable}
\end{center}

\subsection{Reconciliations and mismatches}

\begin{center}\small
\begin{longtable}{@{}L{0.20\textwidth}L{0.74\textwidth}@{}}
\toprule
\textbf{Concept} & \textbf{Two definitions, and the theorem relating them} \\
\midrule
\endhead
a step of occurrences & the core relation and the first-order instance coincide (\lean{OccurrenceCorePlacement.core\_step\_iff\_admitted}) \\
absence & the guard $\mathsf{absent}(p)$ and the predicate query with an empty census for the region of $p$ hold at the same solutions, and moving the guard into the queries keeps admission (\lean{OccurrenceCorePlacement.absenceQuery\_holds\_iff}, \lean{OccurrenceCorePlacement.admits\_absent\_check\_iff\_query}); the two program hosts' absence checks agree on embedded first-order rows, and so do their return firings (\lean{ProgramHosts.absence\_agrees}, \lean{ProgramHosts.return\_admitted\_iff}) \\
closed solvability & the interface's closed solvability and the dependent unifier's are the same at the typed $\beta\eta$ frames, definitionally (\lean{OccurrenceCorePlacement.Closed.closedSolvable\_iff\_frames}) \\
transport & strict transports are transports up to dead names with the same state classes (\lean{Interface.Transport.toDead\_same\_iff}); the strict separation law fails natively (\lean{Interface.NativeUpToDeadLaws.strict\_separation\_fails}) \\
queries & a raw label query is not invariant under knowledge equivalence (\lean{Semantic.SemanticControls.syntactic\_collision}); a query read through solutions is (\lean{Semantic.SemanticControls.noYSemantic\_agrees}); the same split at dependent frames (\lean{Interface.DependentFull.Controls.syntactic\_collision}) \\
freshness & absent from the space versus never born: a firing whose made identity is absent and was born earlier is admitted and is realized by no commit protocol run after that history (\lean{Interface.NativeFreshnessControls.maker\_absent}, \lean{Interface.NativeFreshnessControls.maker\_not\_realized}); the protocol realizes exactly the runs fresh from the born identities (\lean{Interface.native\_run\_protocol\_iff}); a space is born by a firing that makes its marker, and ``fresh'' is never born (\lean{SpacesCreation.created\_never\_born}) \\
solvability & symbolic and closed solvability differ at an obligation of an empty type (\lean{Interface.Dependent.obligation\_trap}, \lean{ContextualDependentClosedBoundary.Controls.point\_not\_closedSolvable}) and agree under a principal answer with a closed point (\lean{ContextualDependentClosedBoundary.principal\_closed\_iff}) \\
independence & the interface's semantic diamond and the static independence of guard-free firings coincide on static firings (\lean{AnnotatedRuns.firstOrder\_independent\_iff\_static}) \\
first-order terms & the one-frame instance over first-order terms and the parameterless fibre of the declaration families are two presentations; the occurrence core is generic over substitution actions and is placed for both, but no map between the two is proved \\
application & the closure machine of the program core computes substitution $\beta$, both ways, on closed terms (\lean{ProgramSubstitution.plotkin\_closed}, \lean{ProgramSubstitution.machine\_iff}), and every such evaluation is a run of the contextual $\lambda$ execution (\lean{ProgramSubstitution.eval\_executes}); $\beta$ as an equation of the binder core agrees with both on runs (\lean{ProgramBinders.Nested.lambda\_three\_way}) \\
compiled and interpreted programs & one compiled firing answers a call exactly when the fixed interpreter's run answers it, up to renaming both ways, for calls and rows below the call's scope and rows that are data (\lean{ProgramInterpreterJoin.interpreter\_join\_iff}) \\
\bottomrule
\end{longtable}
\end{center}

\section{The requirements and their theorems}\label{sec:map}

Each row gives the requirement, the theorems, the instance, and a classification: a
\emph{restriction} of the step, an \emph{encoding} over it, a \emph{necessary extension}, or a
\emph{proved obstruction}. An encoding row expresses a feature by rules or data of the step, with
the correspondence its theorem states; no row states that a translation is compositional. A hosted
source language with a correspondence of whole runs is given for a first fragment in
Section~\ref{sec:addendum-hosting}. The
status is proved unless a row says otherwise. The
separating example of each row is in Section~\ref{sec:examples} or named in the row.

\subsection{Open unification}\label{sec:open-unification}

\begin{center}\small
\begin{longtable}{@{}L{0.20\textwidth}L{0.52\textwidth}L{0.22\textwidth}@{}}
\toprule
\textbf{Requirement} & \textbf{Theorems} & \textbf{Instance; class} \\
\midrule
\endhead
compatible shared constraints & two interactions with shared equations commute when resources are independent, the joint equations are consistent and writes avoid the other's guards (\lean{Core.independently\_ready\_commute}, \lean{Core.Controls.nontrivial\_guarded\_interactions\_commute}); the tells $X\doteq 1$, $Y\doteq 2$ are independent (\lean{AnnotatedRuns.LawfulControls.unlawful\_changes\_independence}) & first order; restriction \\
incompatible shared constraints & tells of $X\doteq\mathrm{true}$ and $X\doteq\mathrm{false}$: each alone, neither after the other (\lean{Interface.NativeResidualExample.tellers\_not\_after}); at dependent frames, $X\doteq \mathrm{Ground}$ against $X\doteq \mathrm{Ground}\to\mathrm{Ground}$ (\lean{Interface.DependentFull.IndependenceControls.tells\_clash\_refused}) & native; dependent full \\
fresh local variables & helper families renamed by isomorphisms keep solvability (\lean{Interface.Frames.Iso.solvable\_applyStore\_iff}); consumption and production need the fresh family (\lean{ContextualExecutionControls.HelperConsumption.actual\_consumption\_and\_production\_require\_the\_fresh\_family}); helper telescopes at dependent frames (\lean{Interface.DependentFull.telescope\_diamond}). The rows on the one-frame first-order step (\lean{OccurrenceCorePlacement.Examples.OpenUnification.fresh\_second\_read\_admitted}, the program core) exercise freshness by scope numbering, not by frame extension & contextual; dependent full; encoding \\
genuine binder crossing & composition under a real binder (\lean{ContextualUnifierComposition.Controls.both\_orders\_act\_under\_a\_real\_binder}); dependent pattern answers (\lean{ContextualDependentGeneralPattern.Pattern.replace\_principal}); principal with $\eta$ only (\lean{ContextualDependentFibre.Controls.point\_pattern\_principal} against \lean{ContextualDependentFibre.Controls.point\_pattern\_no\_principal}); a take through a binder, admitted on both sides of the fibre (\lean{Interface.DependentFull.FibreControls.binder\_take\_admitted}) & dependent $\eta$ frames; extension of the frames \\
intentional caller sharing & a shared second read is refused (\lean{ContextualSchematicStorage.Controls.the\_explicit\_shared\_second\_read\_is\_refused}, on the activated step, whose shared-row, guard-free profile is the occurrence core over the contextual action, \lean{ContextualActiveStepProfile.step\_iff}; that this control lies in the profile is not proved); a shared output keeps one reading, a reusable one has independent readings (\lean{ContextualPublication.Controls.shared\_publication\_keeps\_one\_shared\_reading}, \lean{ContextualPublication.Controls.the\_reusable\_output\_has\_both\_independent\_readings}) & activated step; restriction \\
solver guarantees & first order decides; patterns decide; the dependent driver never refutes a solvable store and loses no solution (\lean{ContextualDependentWorklist.drive\_no\_loss}); decisiveness on one-equation stores where a rule fires (\lean{ContextualDependentDecisiveFragment.single\_decided}); a non-pattern store is suspended by the fragment and rigid solvers at every budget and number of rounds (\lean{ContextualDependentFragment.Controls.repeated\_pending}) & tower \\
the earlier rule set does not decide the fragment & two fragment stores are decided at no budget: the unsolvable $\lambda a\,b.\,\mathrm{Id}\,T\,a\,b\doteq\lambda a\,b.\,\mathrm{Id}\,T\,b\,a$ and the solvable $T\doteq A$ over $A,T,G:(x{:}A)\to T$ (\lean{ContextualDependentFragment.Controls.fragment\_not\_decisive}, \lean{ContextualDependentFragment.Controls.blocked\_not\_decisive}) & proved obstruction for that rule set \\
repairs & identity and $\Pi$ decomposition (\lean{ContextualDependentFragment.Controls.id\_identity\_decided}, \lean{ContextualDependentFragment.Controls.pi\_principal}); solving in place (\lean{ContextualDependentInPlace.Controls.blocked\_decided\_principal}); the occurs check along domains, codomains and carriers (\lean{ContextualDependentOccursFull.occurrence\_not\_typeEq}, \lean{ContextualDependentOccursFull.Controls.codomain\_in\_fragment\_refuted}); a rank that every decomposition lowers (\lean{ContextualDependentRank.weightRank\_wf}, \lean{ContextualDependentRank.rankedSettledOfConst}); rigid clashes under abstractions (\lean{ContextualDependentFragmentDecisive.ClashControls.clashes\_decided}) & dependent $\eta$ frames \\
the census and the stratified class & each fragment equation is classified by the shapes of its opened sides (\lean{ContextualDependentCensus.census\_fragment}); a class beyond the fragment whose annotations agree under every solution of the earlier equations contains it and is kept by the scan and the twin step (\lean{ContextualDependentClosure.stratified\_of\_fragment}, \lean{ContextualDependentClosure.Stratified.splice}, \lean{ContextualDependentAcceptance.scan\_stratified}, \lean{ContextualDependentAcceptance.collapse\_stratified}); the solver chain keeps every outcome of the previous family (\lean{ContextualDependentAcceptance.drive\_conservative}, \lean{ContextualDependentAcceptance.arguments\_decided\_of\_any}) & dependent $\eta$ frames \\
decisiveness on the whole dynamic pattern fragment & \emph{false as stated}. With universes included silently, the fragment admits a solvable store with no principal answer, and every exact solver stops there postponed (\lean{UnificationWithLevels.DependentFibre.cross\_postponed}). Replaced by exact steps that end solved, refuted or postponed, with decision by shape on the uniform pattern fragment of a cumulative model with level variables (Section~\ref{sec:addendum-levels}) & tower; proved obstruction \\
\bottomrule
\end{longtable}
\end{center}

\subsection{Absence}

\begin{center}\small
\begin{longtable}{@{}L{0.20\textwidth}L{0.52\textwidth}L{0.22\textwidth}@{}}
\toprule
\textbf{Requirement} & \textbf{Theorems} & \textbf{Instance; class} \\
\midrule
\endhead
$\mathit{cancel}(0)$ with $u$ unknown & refused (\lean{OccurrenceCorePlacement.AbsenceExample.ask\_refused\_unknown}) & first order; restriction \\
later $u\doteq 1$, $u\doteq 0$ & admitted after $u\doteq 1$, refused after $u\doteq 0$; the run $\mathit{bind}\,1;\mathit{ask}$ (\lean{OccurrenceCorePlacement.AbsenceExample.ask\_admitted\_after\_one}, \lean{OccurrenceCorePlacement.AbsenceExample.ask\_refused\_after\_zero}, \lean{OccurrenceCorePlacement.AbsenceExample.later\_equation\_run}) & first order \\
absence at every solution & absence holds at every solution exactly when each extension that would match is refuted, with no principal answer needed (\lean{Interface.DependentFull.entails\_absent\_iff}); at dependent frames the answer can change through a type-level solution alone (\lean{Interface.DependentFull.Controls.absence\_refused\_after\_ground}, \lean{Interface.DependentFull.Controls.absence\_admitted\_after\_arrow}) & dependent full \\
contextual absence by certified refutation & ``no equation matches'' is one guard per equation on the guarded step, holding exactly when no equation matches any reading (\lean{ProgramBinders.guarded\_checks\_iff}); decided by refuting a joint problem: sound (\lean{ProgramBinders.certified\_guarded}), complete for Miller patterns (\lean{ProgramBinders.certifies\_guarded\_of\_absent}), and never for a matched equation at any budget (\lean{ProgramBinders.guarded\_never\_certified}) & contextual guarded step; restriction \\
absence stable under later bindings & every later instance of an unmatched call is unmatched (\lean{ProgramNested.absence\_stable}); under binders, in every frame (\lean{ProgramBinders.absence\_stable}) & ordinary; contextual \\
incomplete scans & a scan is sound exactly under coverage, and coverage is the apartness of every omitted row (\lean{Check.entails\_absent\_iff\_of\_covered}, \lean{Check.coversMatches\_iff\_omitted\_apart}); the partial scan of the example lacks it (\lean{FirstOrderChecks.Controls.OpenAbsence.the\_partial\_scan\_lacks\_coverage}); an unfinished search never reports none (\lean{ContextualAttemptOutcomes.Attempt.unfinished\_cannot\_report\_none}) & checks over the first-order step; restriction \\
relevant insertions & inserting $\mathit{cancel}(1)$ after $u\doteq 1$ refuses (\lean{OccurrenceCorePlacement.AbsenceExample.relevant\_insertion\_refuses}); in compiled clauses, \lean{ContextualNativeGuardControls.a\_matching\_insertion\_is\_detected} & first order; native \\
unrelated updates & inserting $\mathit{cancel}(2)$ keeps admission (\lean{OccurrenceCorePlacement.AbsenceExample.unrelated\_insertion\_keeps}); writers outside a guard's region keep it (\lean{Interface.Step.Guards.guard\_fire}); \lean{ContextualNativeGuardControls.an\_unrelated\_insertion\_is\_permitted} & first order; native \\
the original cut & the earlier cut refuses, a later one would accept (\lean{ContextualNativeGuardControls.the\_actual\_earlier\_query\_is\_refused}, \lean{ContextualNativeGuardControls.choosing\_the\_later\_cut\_would\_accept}); the firing's own equation counts (\lean{ContextualNativeGuardControls.OwnEquations.the\_clause\_own\_equation\_is\_retained}); an earlier count cannot be rebased (\lean{ContextualCausalExecutionControls.the\_original\_count\_cannot\_be\_rebased\_to\_the\_final\_cut}); at the interface, \lean{InterfaceTable.cut\_preserved} & native; interface \\
conflicting evidence against an inconsistent store & at knowledge with no solution no firing is admitted (\lean{Evidence.not\_admits\_of\_not\_solvable}); rows for $p$ and its opposition coexist with consistent knowledge and still steer admission (\lean{Evidence.Rows.conflict\_consistent}, \lean{Evidence.Rows.accept\_refused\_with\_opposition}); a store holding both tells admits nothing (\lean{Evidence.Rows.both\_tells\_admit\_nothing}) & first order; restriction \\
absence, refutation, incomplete search and disequality & four ways of not finding $p(a)$ give four distinct verdict profiles on four probes (\lean{Evidence.Family.profiles\_distinct}), and a scan that skips the open row would license the absence (\lean{Evidence.Family.incomplete\_scan\_would\_license}); the same profiles on the dependent full step (\lean{EvidenceInterface.DependentFamily.profiles\_distinct}) & first order; dependent full \\
disequality tells & a necessary extension, built: disequation frames over any frame action (\lean{DisequalityFrames.frames}); after $X\neq a$, $X\doteq b$ is admitted and $X\doteq a$ refused (\lean{Disequality.Controls.disequality\_checked}); absence changes with an insertion, a disequality does not (\lean{Disequality.Controls.absence\_vs\_disequality}); conservative over equation-only knowledge (\lean{Disequality.conservative}); principal answers live in constrained frames, and no single substitution presents $X\neq a$ (\lean{Disequality.principal}, \lean{Disequality.Controls.no\_presenting\_answer}) & disequation steps; necessary extension \\
\bottomrule
\end{longtable}
\end{center}

\subsection{Composition}

\begin{center}\small
\begin{longtable}{@{}L{0.20\textwidth}L{0.52\textwidth}L{0.22\textwidth}@{}}
\toprule
\textbf{Requirement} & \textbf{Theorems} & \textbf{Instance; class} \\
\midrule
\endhead
resource independence with information interference & trace-independent footprints, and only one order runs, because a guard reads the other's tell (\lean{AnnotatedRuns.Controls.told\_check\_one\_order}); a query against a writer into its region (\lean{AnnotatedRuns.Controls.census\_writer\_one\_order}) & first order; restriction \\
retained guards & a later binding does not repair an earlier check (\lean{ContextualScheduleControls.EarlierSchema.a\_later\_binding\_does\_not\_repair\_the\_compiled\_earlier\_check}); a firing cannot discard a countermodel (\lean{ContextualPrefixGuardControls.an\_actual\_firing\_cannot\_discard\_that\_countermodel}) & contextual; restriction \\
joint solutions against failure prefixes & every order of a jointly solvable clear family runs and all agree up to dead names (\lean{Interface.NativeIterated.every\_order\_admitted}, \lean{Interface.NativeIterated.every\_order\_agrees}); without joint solvability no order runs (\lean{Interface.NativeIterated.no\_order\_of\_not\_jointly}); a refused joint problem keeps the executed producer prefix (\lean{ContextualCausalExecutionControls.joint\_refusal\_does\_not\_erase\_the\_executed\_producer\_prefix}); a refused block leaves no residue (\lean{StagedSourceEffects.refused\_has\_no\_residue}) & native; first order \\
independence is a property of the step & two lawful transports with the same state classes have the same independent pairs (\lean{Interface.Transport.Lawful.independent\_iff}); an unlawful residual changes independence (\lean{AnnotatedRuns.LawfulControls.unlawful\_changes\_independence}) & interface \\
congruence & with no program: the ordinary step that offers every firing under its rule can fire any context itself, and its rule-labelled bisimilarity is a congruence for disjoint contexts (\lean{Semantic.ruleBisimilar\_tensor}); on configurations with consistent knowledge, given a fresh supply of identities, that bisimilarity is equivalence modulo the knowledge (\lean{LogicOverTheStep.ProgramFreeBisimilarity.ruleOffers\_bisimilar\_iff}; controls \lean{LogicOverTheStep.ProgramFreeBisimilarity.ProgramFreeControls.consistency\_needed}, \lean{LogicOverTheStep.ProgramFreeBisimilarity.ProgramFreeControls.freshSupply\_needed}), so this congruence states that a static equivalence is preserved by a disjoint store and is not a congruence result for programs (Section~\ref{sec:addendum-program-bisimilarity}); for a program, closed bisimilarity is not a congruence (\lean{Observation.CongruenceControls.closed\_not\_congruence}); for a triggered program, without predicate queries and checks, minimal bisimilarity, which is open bisimilarity there, is a congruence for contexts of port atoms (\lean{Race.TriggeredProgram.tensor\_ports}). An equivalence of programs preserved by contexts that add rules, extend the frame or run another program is open (Section~\ref{sec:open}) & first order: the program-free step; triggered programs \\
native diamonds & composites of a clear diamond are bisimilar and step bisimilar (\lean{Interface.NativeDead.composites\_bisimilar}, \lean{Interface.NativeUpToDead.makers\_composites\_stepBisimilar}); a concurrent step is bisimilar to each sequentialization (\lean{Interface.NativeConcurrent.Concurrent.bisimilar}) & live native \\
dependent diamonds & helper-free firings close on the nose under five independence clauses, each but the query clause necessary (\lean{Interface.DependentFull.local\_diamond}, \lean{Interface.DependentFull.clauses\_of\_residuals}); helper firings up to the block exchange (\lean{Interface.DependentFull.telescope\_diamond}, \lean{Interface.DependentFull.TelescopeControls.exchanged\_diamond}, \lean{Interface.DependentFull.TelescopeControls.composites\_differ}); a dependency between the made rows excludes any exchange (\lean{Interface.DependentFull.no\_residual\_of\_mentions}) & dependent full \\
probabilities of causal histories & for positive rate systems and maximal runs of two parts, the jump-chain likelihood summed over all interleavings is the product of the parts' likelihoods (\lean{CausalHistories.sum\_shuffles}, \lean{CausalHistories.history\_factorizes}), with the ternary coherence (\lean{CausalHistories.history\_factorizes\_three}); on two ground races of the step (\lean{CausalHistories.StepControls.step\_history\_product}, \lean{CausalHistories.StepControls.local\_exact}); on a confused contention set the per-firing reading depends on the order (\lean{CausalHistories.Controls.confusion\_control}) & rate systems; first order \\
contention & contention is a shared take, a take of what the other reads, or inhibition: a take or make inside the other's predicate-query region, and contending firings are not independent (\lean{Contention.not\_independent\_of\_contends}); contention as sharing alone gives local likelihoods of total mass $2$, contention as interference gives the step's (\lean{Contention.Controls.inhibition\_control}, \lean{Contention.Controls.maker\_guarded\_dependent}) & first order with queries \\
interventions & a \mt{do} restricts weights, composes by conjunction and commutes (\lean{Intervention.restrictWeight\_comm}), lowers exactly on static alphabets (\lean{Intervention.prob\_restrictWeight}), and at a space reaches only firings with the right to make there (\lean{Intervention.doAt\_of\_no\_send}); one atom through an absence guard performs a hard intervention (\lean{Intervention.Controls.context\_performs\_do}) & first order; encoding \\
\bottomrule
\end{longtable}
\end{center}

\subsection{Resources}

\begin{center}\small
\begin{longtable}{@{}L{0.20\textwidth}L{0.52\textwidth}L{0.22\textwidth}@{}}
\toprule
\textbf{Requirement} & \textbf{Theorems} & \textbf{Instance; class} \\
\midrule
\endhead
aliased reads & one occurrence read twice is admitted (\lean{OccurrenceCorePlacement.Examples.Resources.aliased\_read\_admitted}); two read views count as one stored row (\lean{SchematicReach.Controls.two\_read\_views\_still\_count\_as\_one\_stored\_row}); unrestricted aliasing admits a firing the register semantics has no counterpart for (\lean{MetaMeTTaComparison.Family.unrestricted\_read\_aliasing\_admits\_another\_firing}) & first order; restriction \\
distinct takes & one occurrence taken twice is refused (\lean{OccurrenceCorePlacement.Examples.Resources.take\_same\_twice\_refused}); a consumed cursor cannot be replayed (\lean{NeedSnapshots.Controls.consumed\_cursor\_cannot\_be\_replayed}) & first order \\
occurrence multiplicity & two equal-valued occurrences are both taken (\lean{OccurrenceCorePlacement.Examples.Resources.take\_two\_admitted}); equal payloads keep two cells (\lean{NeedSnapshots.Controls.equal\_payloads\_keep\_two\_cells}) & first order \\
read before consume & reading then taking is a run, the converse order is refused (\lean{OccurrenceCorePlacement.Examples.Resources.read\_then\_take}, \lean{OccurrenceCorePlacement.Examples.Resources.read\_after\_take\_refused}) & first order \\
several spaces & one space whose occurrences carry handles behaves as several: enabling and results split by handle, and the frame is unchanged (\lean{Spaces.spaces\_injective}, \lean{Spaces.enabled\_iff\_spaces}, \lean{Spaces.present\_result\_spaces}, \lean{Spaces.fiber\_unchanged}); every run of every instance is an executable commit word (\lean{Spaces.run\_executable}); on the located first-order step, admission at a handle depends only on its space and the knowledge (\lean{Spaces.admits\_iff\_of\_fiber}) & interface; located first order; encoding \\
independence across spaces & on the restriction \emph{quiet} (Section~\ref{sec:joined}); the knowledge is shared, so isolation between spaces is physical, not informational (\lean{Spaces.Controls.shared\_variable\_crosses\_spaces}) & first order, dependent full; restriction \\
creating a space & a firing that makes a birth marker at an unborn handle creates a space; markers are never taken, and every occurrence lies at a handle already born (\lean{SpacesCreation.created\_never\_born}, \lean{SpacesCreation.wellBorn\_result}) & first order; restriction \\
rights per handle & a firing without \mt{receive}, \mt{peek} or \mt{send} at a space is refused (\lean{Permissions.refused\_without\_receive}, \lean{Permissions.refused\_without\_peek}, \lean{Permissions.refused\_without\_send}); without \mt{receive} a space only grows and a witness survives, on every instance (\lean{PermissionsInterface.witness\_present}); without \mt{send} it only shrinks and an absence survives, on instances whose labels move along firings (\lean{PermissionsInterface.absence\_preserved}) & interface; restriction \\
the authority escape & on runs of ticketed firings the rows made plus the tickets left never exceed the tickets at the start (\lean{Permissions.made\_le\_tickets}); rights alone allow unbounded writes (\lean{Permissions.Controls.rights\_alone\_unbounded}), a ticket per write blocks them (\lean{Permissions.Controls.ticketed\_blocks}) & first order; restriction \\
earlier space and permission layers & permitted runs are non-interfering (\lean{Permissions.permitted\_noninterference}); the reach step is a restriction of the first-order step (\lean{SpacesPlacement.reachStep\_iff\_admitted}); the delegation core and the transaction layers are not placed & first order \\
\bottomrule
\end{longtable}
\end{center}

\subsection{Evaluation and the common operations}

\begin{center}\small
\begin{longtable}{@{}L{0.20\textwidth}L{0.52\textwidth}L{0.22\textwidth}@{}}
\toprule
\textbf{Requirement} & \textbf{Theorems} & \textbf{Instance; class} \\
\midrule
\endhead
expression application & $\beta$ up to terminating administrative firings (\lean{Program.beta\_on\_step}, \lean{Program.administrative\_wellFounded}); an atom applied is stuck (\lean{Program.atom\_application\_stuck}); the closure machine computes substitution $\beta$ (\lean{ProgramSubstitution.plotkin\_closed}); contextual $\lambda$: \lean{LambdaComparison.source\_step\_iff\_executes} & program core; encoding \\
lexical binding & every halted run returns the uncaptured value (\lean{Program.no\_capture\_on\_step}) & program core; encoding \\
correlated against independent choice & shared choice always agrees, independent choices can disagree, over all runs (\lean{Program.shared\_vs\_independent\_on\_step}) & program core; encoding \\
inert quotation, explicit activation & quoted code halts as a value, activation runs it one layer (\lean{Program.quotation\_inert}, \lean{Program.activate\_quote\_on\_step}, a run of the machine, not one firing); activation needs code (\lean{Program.activation\_needs\_code}) & program core; encoding \\
pattern equations & a call is matched against every equation by two-way unification, one branch per matching equation in program order (\lean{ProgramEquations.answers\_exact}); admission is unifiability (\lean{ProgramEquations.call\_valid\_iff}); the call's own variable is bound by the equation (\lean{ProgramEquations.Controls.variable\_bound\_on\_step}) & ordinary step; encoding \\
nested calls & call by value, leftmost-innermost, with held data \mt{val} where a position is declared \mt{Atom} (the elaboration by declared type is not formalized); one schema per machine rule, both ways (\lean{ProgramNested.nested\_execution}); the next redex is unique (\lean{ProgramNested.decomposition\_unique}); reduction in contexts agrees up to renaming of fresh scopes (\lean{ProgramNested.reference\_complete}) & ordinary step; encoding \\
unmatched calls return themselves & one firing with an absence predicate query over the equations of the version read (\lean{ProgramNested.unmatched\_fire}, \lean{ProgramNested.u\_execution}); under binders, a guarded firing of \mt{(= \$C \$C)} admitted exactly when no equation matches, final after it fires (\lean{ProgramBinders.returns\_iff}, \lean{ProgramBinders.returned\_is\_final}, \lean{ProgramBinders.evaluates\_iff}) & ordinary; contextual guarded \\
equations under binders & source rewriting by holes with declared parameters is exactly the readings of runs (\lean{ProgramBinders.execution\_iff}); principal answers for Miller patterns on the step (\lean{ProgramBinders.principal\_answers\_on\_step}, \lean{ProgramBinders.miller\_with\_solvers}); controls: $\eta$, capture refused, a non-pattern pending at every fuel (\lean{ProgramBinders.Controls.eta\_on\_step}, \lean{ProgramBinders.Controls.capture\_refused\_on\_step}, \lean{ProgramBinders.Controls.twice\_pending\_on\_step}) & contextual guarded; encoding \\
strategies under binders & innermost: one guarded firing per position is exactly one reference step, runs agree both ways for determinate equations, and the checks are decided by certified refutation (\lean{ProgramBinders.Strategy.innermost\_step\_iff}, \lean{ProgramBinders.Strategy.innermost\_execution\_iff}, \lean{ProgramBinders.Strategy.innermost\_of\_certified}, \lean{ProgramBinders.Strategy.innermost\_redex\_certified}); leftmost-innermost with \mt{Atom} positions: the next redex is unique, and without declarations one step agrees both ways and runs agree for determinate equations (\lean{ProgramBinders.Leftmost.selects\_unique}, \lean{ProgramBinders.LeftmostStep.leftmost\_step\_iff}, \lean{ProgramBinders.LeftmostStep.leftmost\_execution\_iff}); with declarations sound, and complete when nothing frozen inside the redex matches (\lean{ProgramBinders.LeftmostStep.leftmost\_step\_sound\_atoms}, \lean{ProgramBinders.LeftmostStep.leftmost\_step\_complete\_atoms}) & contextual guarded; complete \mt{Atom} positions on the step \emph{open}, 40\% \\
unused divergent work and call by need & quoted divergent code halts, activated it diverges (\lean{Program.quoted\_loop\_halts}, \lean{Program.activated\_loop\_diverges}); call by need at first order: Launchbury's semantics and its machine agree (\lean{ProgramNeed.need\_adequacy}), a thunk is evaluated at most once and later demands read its value (\lean{ProgramNeed.later\_demand\_reads}), agreement with call by value for determinate programs and pure terminating thunks (\lean{ProgramNeed.need\_value\_agree}), on the ordinary step (\lean{ProgramNeed.need\_on\_step}); a loop bound and unused answers by need and has no answer by value (\lean{ProgramNeed.need\_answers\_loop}, \lean{ProgramNeed.value\_diverges\_loop}); a shared choice is never mixed (\lean{ProgramNeed.choice\_never\_mixed}) & program core; ordinary step; under binders \emph{open}, 45\% \\
effects & immediate effects of a failed branch stay, deferred ones vanish (\lean{Program.immediate\_effects\_leave\_a\_trace}, \lean{Program.deferred\_failed\_branch\_leaves\_no\_trace}); a deferred block is all or nothing (\lean{Program.deferred\_block}) and one firing at the cut (\lean{Program.deferred\_one\_firing}); \mt{defer} as a source construct (\lean{ProgramSourceDefer.defer\_one\_firing}) & program core; encoding with computed rules \\
atomic blocks & a rule family with a uniform bound on made rows cannot commit a larger block in one firing (\lean{OccurrenceCorePlacement.no\_bounded\_family\_commits}, \lean{Program.no\_fixed\_atomic\_commit}) & first order; \emph{proved obstruction} to uniformly bounded rules \\
suspension & a suspended continuation resumes exactly once per demand (\lean{Program.exactly\_once\_per\_demand}); a read demand replays without bound (\lean{Program.replay\_unbounded}) & program core; encoding \\
changing definitions & a conclusion does not carry to a new version without a recheck (\lean{Program.unconditional\_transport\_unsound}); what stays valid (\lean{Program.what\_stays\_valid}); for pattern equations, adding equations keeps the old answers and appends the new (\lean{ProgramEquations.Versions.answers\_append}); an unmatched answer cites the version it read, and installing a later one changes no earlier answer (\lean{ProgramNested.UnmatchedControls.no\_retroactive\_change\_on\_step}) & ordinary step; on the binder core \emph{open}, 45\% \\
interpreting programs & a fixed rule family cannot answer every call with variables in one firing (\lean{ProgramCompiled.no\_fixed\_one\_firing\_interpreter}); most general unifiers are unique up to renaming (\lean{ProgramMGURenaming.mgu\_unique\_up\_to\_renaming}); the fixed interpreter and the compiled program give the same answers up to renaming, both ways, on the step (\lean{ProgramInterpreterJoin.interpreter\_join\_iff}, \lean{ProgramInterpreterJoin.interpreter\_join\_on\_step}, \lean{ProgramInterpreterJoin.interpreter\_join\_converse\_on\_step}) & ordinary step; encoding \\
communication & for rho with ordinary substitution, not its semantic-substitution profile, source paths are exactly paths of firings of one rule instance per source redex, whose made components are the source's own reduct (\lean{ContextualCommunicationAdequacy.source\_path\_iff\_firings}), on configurations whose equations the identity solves, so with no open constraints. Rho with semantic substitution is placed by three rule schemas, weakly bisimilar on all closed processes, with the channels matched by the knowledge and the contractum made by a term former whose meaning is the calculus's substitution (\lean{FragmentsOfTheStep.ReflectiveCalculus.reflective\_calculus\_on\_the\_step}, Section~\ref{sec:addendum-rho}) & occurrence core over the process-term action; encoding by one rule per redex; first-order step, three schemas \\
querying & complete answers are exactly the matches at a solution (\lean{ContextualQueryAnswers.actual\_answer\_iff}); query rows realized by census firings (\lean{CensusSourceCompilation.realize\_query\_row}) & first order; restriction \\
\bottomrule
\end{longtable}
\end{center}

\subsection{Reflection}

\begin{center}\small
\begin{longtable}{@{}L{0.20\textwidth}L{0.52\textwidth}L{0.22\textwidth}@{}}
\toprule
\textbf{Requirement} & \textbf{Theorems} & \textbf{Instance; class} \\
\midrule
\endhead
declared syntax and value readings & ordinary queries depend only on represented rows (\lean{Program.semantic\_admission\_iff}); modulo equivalence is a bisimulation of the ordinary step (\lean{Observation.ModuloEquivalent.semantic\_bisimilar}), and on configurations with consistent knowledge, given a fresh supply of identities, exactly its bisimilarity (\lean{LogicOverTheStep.ProgramFreeBisimilarity.ruleOffers\_bisimilar\_iff}) & ordinary step; restriction \\
counterexamples to identification & raw queries collide (\lean{Semantic.SemanticControls.syntactic\_collision}, \lean{Program.rawQuoted\_collision}); the census controls read syntax (\lean{Semantic.persistent\_not\_semantic}); telling a structural equality of distinct constants is inconsistent (\lean{MetaMeTTaComparison.OperationalEquality.telling\_that\_structural\_equality\_would\_be\_inconsistent}); an admitted object theory makes syntactic refutation unsound and head readouts ill defined (\lean{BetaEtaReadings.Application.refutation\_not\_admitted}, \lean{BetaEtaReadings.StoredRedex.head\_readout\_not\_invariant}); reduction does not identify quoted syntax (\lean{QuotedLambdaRetained.Controls.reduction\_does\_not\_identify\_the\_quoted\_syntax}) & first order; separation \\
preservation under the restrictions & the pure fragment agrees with the original execution (\lean{QuotedLambdaRetained.pure\_fragment\_agrees\_with\_original\_execution}); quoted code ignores the environment (\lean{RhoContinuation.quoted\_code\_ignores\_environment}); corrected nested reflection is sound and the naive step misses it (\lean{Mettapedia.Languages.ProcessCalculi.MeTTaCalculus.Corrected.Controls.corrected\_nested\_sound}, \lean{Mettapedia.Languages.ProcessCalculi.MeTTaCalculus.Corrected.Controls.naive\_step\_misses\_nested\_reflection}) & contextual; encoding \\
reflection as a source construct & \mt{reflect} asks a raw syntax question, and its firing is not one of the ordinary step: two configurations with the same rows answer it differently (\lean{ProgramSourceReflection.reflect\_not\_semantic}, \lean{ProgramSourceReflection.reflect\_collision}) & ordinary step extended; separation \\
absence is an ordinary guard & the unmatched-call query reads through solutions (\lean{ProgramNested.absentQuery\_semantic}); on the binder core absence is asked at every future solution (\lean{ProgramBinders.guarded\_checks\_iff}), and a pending problem is told apart from a refutation by solver results, not by reflection (\lean{ProgramBinders.Controls.twice\_pending\_on\_step}) & ordinary; contextual guarded \\
\bottomrule
\end{longtable}
\end{center}

\subsection{Inference}

\begin{center}\small
\begin{longtable}{@{}L{0.20\textwidth}L{0.52\textwidth}L{0.22\textwidth}@{}}
\toprule
\textbf{Requirement} & \textbf{Theorems} & \textbf{Instance; class} \\
\midrule
\endhead
valid derivations succeed & reachability of a judgment is exactly a derivation of the independently specified checker (\lean{InferenceKernelReachability.reaches\_iff\_derivation}); runs reproduce the exact proof (\lean{PersistentInference.realizes\_derivation}); published judgments have source derivations (\lean{InferenceSchemaKernel.System.published\_judgment\_has\_source\_derivation}); MM0 (\lean{Mettapedia.Logic.TierZero.Adequacy.MM0LocatedReachability.located\_kernel\_reaches\_iff\_derives}); typing runs are derivations (\lean{TypingChecking.Kernel.checking\_run\_sound}, \lean{TypingCheckingMeta.checking\_run\_sound}) & authority restriction of the first-order step; typing steps \\
dependent checking as a run & checking tells its obligations into the store, a solution of the store yields the derivation, and an unsolvable store refuses (\lean{TypingCheckingMeta.checking\_run\_sound\_of\_store}, \lean{TypingCheckingMeta.refused\_of\_unsolvable}); the checking step works at a fixed metavariable frame and allocates no fresh metavariable, and its completeness holds when the goal is typed at the identity (\lean{TypingCheckingMeta.switch\_complete}, \lean{TypingCheckingMeta.app\_complete}) & typing step \\
the certain checks as a calculus & a derivation calculus for the first-order checks is sound for free Henkin models and complete at generic points (\lean{AnnotatedRuns.CertainChecks.derives\_sound\_free}, \lean{AnnotatedRuns.CertainChecks.derives\_complete\_generic}); the Herbrand model is sound and not complete (\lean{AnnotatedRuns.CertainChecks.oneConstant\_separates}); a query-free firing is admitted exactly when its admission is derivable (\lean{AnnotatedRuns.CertainChecks.admitsDerivably\_iff}) & first order \\
forged evidence fails & a public row cannot serve as a protected premise (\lean{InferenceAuthorityControls.wrong\_space\_premise\_refused}); a counterfeit conclusion is no activation (\lean{InferenceAuthorityControls.fabricated\_output\_refused}) & first order \\
wrong premises fail & one row cannot supply distinct judgments (\lean{InferenceProjectionControls.one\_row\_cannot\_supply\_distinct\_judgments}); a ground term is not a function (\lean{TypingCheckingControls.ground\_not\_function}) & first order; typing \\
capture fails & assignments cannot capture bound indices (\lean{ScopedInferenceValueControls.contextual\_assignment\_cannot\_capture\_its\_bound\_index}); resolution refuses capture (\lean{Resolution.Controls.capture\_refused}) & first order \\
missing authority fails & a non-writer creates no protected occurrence (\lean{SchemaPublicationAuthority.System.nonwriter\_creates\_no\_protected\_occurrence}); a client cannot publish (\lean{InferenceAuthorityControls.client\_cannot\_publish}) & authority restriction \\
execution confers no validity & the forged producer is admitted by the step and published by no step of its process (\lean{OccurrenceCorePlacement.forged\_admitted\_but\_not\_published}); runtime acceptance of a void contract the step refuses (\lean{TypingContractsControls.void\_trap}); symbolic admission of a hole of an empty type (\lean{Interface.Dependent.trap\_helper}) & \emph{necessary extension}: authority and closed asks \\
guards by polarity & ``never bad'' has a finite certificate exactly when it holds and the datum has finitely many reachable states (\lean{GuardFragments.exists\_invariant\_iff}); the budgeted closure verdict certifies exactly on finite-state data and is suspended at every budget on infinite-state data (\lean{GuardFragments.always\_iff\_exists\_commit}, \lean{GuardFragments.suspended\_of\_infinite}); on open data a reachability guard is decided at a principal answer and a ``never'' guard is not stable (\lean{GuardFragmentsOpen.eventually\_entails\_iff}, \lean{GuardFragmentsOpen.Controls.always\_not\_stable}) & interface asks; restriction \\
lowering by checked certificates & where an invariant kept by every admitted rule makes the specification equivalent to the existence of a certificate, the runs with the specification are exactly the runs with the check, certificates erased (\lean{GuardLowering.runs\_spec\_iff\_lowered}, \lean{GuardLowering.certified\_runs}); a checked certificate is sound whoever made it (\lean{GuardLowering.admitted\_of\_certificate}); a fragment not closed under successors admits wrongly (\lean{GuardLowering.Controls.closure\_free\_admits\_wrongly}); a guard priced by the store it scans breaks interleaving invariance (\lean{MeteredGuards.Controls.scan\_price\_order}) & first order; encoding \\
fair search makes certificates & a monotone search firing, enabled and weakly fair, fires or is preempted (\lean{FairSearch.fires\_or\_preempted}); on finite-state data every weakly fair run's store becomes a certificate the lowered guard checks (\lean{FairSearch.fair\_search\_certifies}); in one store the answer is sound under every schedule and arrives under weak fairness (\lean{FairSearchGuards.answer\_sound}, \lean{FairSearchGuards.answer\_eventually}); search firings that take are not complete under fairness (Example~\ref{ex:takerace}) & commits; restriction to monotone search \\
\bottomrule
\end{longtable}
\end{center}

\section{Minimized examples}\label{sec:examples}

Each example is the smallest we know of on the common structure. All are kernel-checked.

\begin{example}[Open unification; \lean{OccurrenceCorePlacement.Examples.OpenUnification}]
Rows $0:p(1)$ and $1:p(2)$ in the first-order step. After reading row $0$ through $p(X_0)$, reading
row $1$ through the same caller variable is refused, and through a fresh variable is admitted
(\lean{shared\_second\_read\_refused}, \lean{fresh\_second\_read\_admitted}). The tells $X_0\doteq1$ and
$X_0\doteq2$ are each admitted and never both (\lean{tells\_conflict}), and the two failed compositions
leave prefixes that every solution distinguishes (\lean{prefixes\_differ}); a compatible pair runs
(\lean{compatible\_run}). Binder crossing is not first order; its minimal pair is the pattern
equation $m(x)\doteq n(x)$, principal with $\eta$ and not without
(\lean{ContextualDependentFibre.Controls.point\_pattern\_principal},
\lean{ContextualDependentFibre.Controls.point\_pattern\_no\_principal}).
\end{example}

\begin{example}[Absence; \lean{OccurrenceCorePlacement.AbsenceExample}]
Row $0$ is $\mathit{cancel}(0)$, rows $1,2$ are requests. The firing $\mathit{ask}$ takes request $1$
under the guard that no $\mathit{cancel}(u)$ is present. With $u$ unknown it is refused; after the
tell $u\doteq1$ it is admitted; after $u\doteq0$ it is refused. After $u\doteq1$, inserting $\mathit{cancel}(1)$
refuses it and inserting $\mathit{cancel}(2)$ keeps it (\lean{insertions\_separate}). Moving the
guard into an empty-census predicate query gives the same verdicts
(\lean{ask\_as\_query\_refused\_unknown}, \lean{ask\_as\_query\_admitted\_after\_one}).
\end{example}

\begin{example}[Composition]
At the empty state, $\mathit{tell}(X\doteq1)$ and $\mathit{ask}(X\doteq1)$ have trace-independent footprints;
only the order tell-then-ask runs (\lean{AnnotatedRuns.Controls.told\_check\_one\_order}). Three tells
$X\doteq Y$, $X\doteq1$, $Y\doteq2$ are pairwise independent and form no cell
(\lean{AnnotatedRuns.CellControls.hollow\_cube}); natively, three pairwise clear tellers have no order
at all (\lean{Interface.NativeIteratedControls.pairwise\_not\_enough}).
\end{example}

\begin{example}[Resources; \lean{OccurrenceCorePlacement.Examples.Resources}]
Rows $0:c$ and $1:c$. Taking $c$ twice takes both (\lean{take\_two\_admitted}); taking $0$ twice is
refused (\lean{take\_same\_twice\_refused}); reading $0$ twice is admitted
(\lean{aliased\_read\_admitted}); read-then-take runs, take-then-read is refused
(\lean{read\_then\_take}, \lean{read\_after\_take\_refused}).
\end{example}

\begin{example}[Evaluation and reflection]
On the program core: one shared choice used twice agrees with itself in every halted run, two
independent choices disagree in some (\lean{Program.shared\_vs\_independent\_on\_step}). Quoted
$\Omega$ halts and activated $\Omega$ diverges (\lean{Program.quoted\_loop\_halts},
\lean{Program.activated\_loop\_diverges}). Two rows $20:X$ and $21:Y$ under $X\doteq Y$ are equivalent modulo
the knowledge, and a query for the label $Y$ itself separates them, while the query read through
solutions does not (\lean{Semantic.SemanticControls.syntactic\_collision},
\lean{Semantic.SemanticControls.noYSemantic\_agrees}).
\end{example}

\begin{example}[An unmatched call returns itself;
\lean{ProgramNested.UnmatchedControls}]\label{ex:unmatched}
The program is
\begin{center}\ttfamily\frenchspacing (= (g a) b)\qquad (= (f \$x) (out \$x))\end{center}
and the call is \mt{(f (g c))}. No equation matches \mt{(g c)}, so its call returns the data
\mt{(g c)} by one firing whose absence predicate query sees no matching equation at its cut; the
outer call then matches, and the run answers \mt{(out (g c))}
(\lean{inner\_value\_outer\_call\_on\_step}). Without the return firing the machine has no step at the inner call
(\lean{ProgramNested.Controls.unmatched\_inner\_stuck}). Versions: under a version with no equation
for \mt{f}, \mt{(f a)} answers itself and cites that version
(\lean{answered\_as\_itself\_v1\_on\_step}); where \mt{(= (f a) b)} is installed the return is refused
(\lean{unmatched\_refused\_v2}); installing it later changes neither the earlier answer nor its
citation (\lean{no\_retroactive\_change\_on\_step}). Under a binder, with only
\mt{(= (f (\(\lambda\) x. \$N)) \$N)}, the call \mt{(f (\(\lambda\) y. y))} matches nothing, since \mt{\$N} may not
capture \mt{x}, and returns itself
(\lean{ProgramBinders.UnmatchedControls.capture\_returns\_itself}).
\end{example}

\begin{example}[Inhibition: one atom forces a race; \lean{Contention.Controls}]\label{ex:inhibition}
The space holds one occurrence \mt{a}. Two firings race for it, with positive weights $w_r$ and
$w_v$:
\begin{center}\ttfamily\frenchspacing
racer: take a, make out\qquad vetoed: take a, guard (absent p), make out
\end{center}
At this space \mt{racer} wins with probability $w_r/(w_r+w_v)<1$; after one atom \mt{p} is added
from outside, \mt{vetoed} is refused and \mt{racer} wins with probability $1$
(\lean{forcing\_control}). Without inhibition no context forces a race over positive weights
(\lean{Contention.no\_context\_forces}), and no added atoms force \mt{vetoed}
(\lean{cannot\_force\_vetoed}). The same inhibition separates the two readings of contention: a
firing that takes \mt{a} and makes \mt{p}, and one that takes \mt{b} under \mt{(absent p)}, share no
occurrence, yet the first refuses the second; normalizing within shared occurrences gives local
likelihoods of total mass $2$, and normalizing within interference gives the step's own, of total
mass $1$ (\lean{inhibition\_control}).
\end{example}

\begin{example}[A take race that fair search never completes; \lean{FairSearch.Controls}]\label{ex:takerace}
The space holds the goal \mt{goal}, and two search firings race for it:
\begin{center}\ttfamily\frenchspacing
take goal, make dead\qquad take goal, make proof
\end{center}
A run in which the first wins and then idles is strongly fair for every task, has unbounded budget,
and never makes \mt{proof} (\lean{race\_stronglyFair}, \lean{race\_never\_proves}), although
\mt{proof} is reachable (\lean{race\_proof\_reachable}). Reading the goal instead of taking it, every
weakly fair run proves (\lean{read\_race\_proves}). Adding the repair \mt{take dead, make proof}
reaches exactly the same stores (\lean{race\_reachable\_iff}), and every weakly fair run of the
repaired race proves (\lean{repair\_fair\_proves}). So reachability does not decide fair
completeness (\lean{reachability\_does\_not\_decide\_fair\_completeness}). What separates the two
races is that a proof stays reachable from every reachable store, a property bisimilarity preserves
and reachability does not (\lean{FairSearchIdentity.alwaysPossibly\_of\_bisimilar},
\lean{FairSearchIdentity.search\_identity\_is\_branching}); that this property with strong fairness
gives completeness on every finite-state system is open.
\end{example}

\begin{example}[Inference]
The identity is typed by a run of the typing step and a ground term against a function type is
refused (\lean{TypingCheckingControls.identity\_run}, \lean{TypingCheckingControls.ground\_not\_function}).
The producer of a protected occurrence is admitted by the first-order step and reaches its result by
no step of the client's process (\lean{OccurrenceCorePlacement.forged\_admitted\_but\_not\_published}).
\end{example}

\section{Full programs from the core}\label{sec:programs}

\paragraph{The interpreter core.} A source language over S-expressions with closures, lexical
binding, quotation and activation, choice, facts added and removed, suspension and versioned
definitions has an independent reference semantics of 24 rules. Each rule is a schema of take and
read patterns and made terms, copied into a fresh scope at each firing; the two-way unifier solves
the equations. One source step is exactly one firing of the ordinary step, in both directions, for
every program (\lean{Program.step\_iff\_rewrites}, \lean{Program.program\_interface\_iff}), and the
administrative firings terminate (\lean{Program.administrative\_wellFounded}). The laws of
evaluation in Section~\ref{sec:map} are theorems about these runs.

\paragraph{Pattern equations.} A program is a list of equations \mt{(= (f p1 ... pn) rhs)}, each
compiled to one schema that takes the call and makes the instantiated right-hand side; the step's own
knowledge does the unification. Each matching equation gives exactly one branch, in program order,
duplicates included, and a call that matches nothing has no step (\lean{ProgramEquations.answers\_exact},
\lean{ProgramEquations.stuck\_iff}); the reference semantics is an independent inductive relation
over a semantic notion of most general unifier (\lean{ProgramEquations.equations\_execution}).
Calls inside arguments and right-hand sides are evaluated by call by value, leftmost-innermost, the
reading of the HE specification, with positions declared \mt{Atom} written as held data
(\lean{ProgramNested.nested\_execution}). A call that matches nothing returns itself by an absence
predicate query, ordinary, exact at its cut and stable under later bindings, and its answer cites the
version it read (Example~\ref{ex:unmatched}).

\paragraph{Compiled and interpreted.} A fixed interpreter of 13 equations answers calls with
variables by allocating one fresh variable per slot in administrative firings and unifying the
renamed copy in one firing; its answers are the compiled ones up to renaming, in both directions,
at the source and on the step (\lean{ProgramInterpreterJoin.interpreter\_join\_iff},
\lean{ProgramInterpreterJoin.interpreter\_join\_converse\_on\_step}). The allocation is forced: a
fixed rule family cannot answer every call with variables in one firing
(\lean{ProgramCompiled.no\_fixed\_one\_firing\_interpreter}). This is the positive reflective result:
an interpreter on the step that runs stored equations as data.

\paragraph{Equations under binders.} On the contextual guarded step, over any binding signature, an
equation is a pair of terms with holes that list the object variables their solutions may use. One
schema per equation; the equation \mt{lhs} $\doteq$ \mt{call} joins the retained store, and a hole
no equation constrains stays unknown, to be bound by a later firing. Source rewriting is exactly the
readings of runs (\lean{ProgramBinders.execution\_iff}), and for Miller patterns the actual worklist
computes one principal answer at every firing (\lean{ProgramBinders.principal\_answers\_on\_step},
\lean{ProgramBinders.miller\_with\_solvers}). The $\lambda$ calculus and the quoted-$\lambda$ account are
instances: $\beta$ is \mt{(= ((\(\lambda\) x. \$B x) \$A) \$B[\$A])} and activation is
\mt{(= (activate (quote \$C)) \$C)} (\lean{ProgramBinders.Nested.lambda\_three\_way},
\lean{ProgramBinders.Quoted.executes\_iff\_rewrites}). Absence of a matching equation is a guard per
equation, decided by certified refutation, and it gives unmatched calls that return themselves and
innermost and leftmost-innermost reduction under binders (Section~\ref{sec:map}). On embedded
first-order rows the binder core agrees with the first-order core
(\lean{ProgramBinders.FirstOrder.step\_join}, \lean{ProgramHosts.absence\_agrees}).

\paragraph{Call by need.} At first order, a heap cell holds a thunk; the first demand marks it busy,
evaluates it and writes the value, and every later demand reads the value. Launchbury's semantics,
its machine and the ordinary step agree (\lean{ProgramNeed.need\_adequacy},
\lean{ProgramNeed.need\_eval\_on\_step}); each cell is its own row, and a demand of a value reads the row
without taking it.

\paragraph{Computed rules.} A rule family with a uniform bound on made rows bounds what one firing
commits (\lean{OccurrenceCorePlacement.added\_card\_le},
\lean{OccurrenceCorePlacement.no\_bounded\_family\_commits}). Atomic blocks of unbounded size are
therefore realized by a rule computed per block, which is still one firing of the step at the
block's cut (\lean{Program.deferred\_one\_firing}, \lean{ProgramSourceDefer.defer\_one\_firing},
\lean{StagedSourceEffects.refused\_has\_no\_residue}), or by a compound occurrence carrying the block.

\paragraph{Compiled clauses and contextual executions.} Compiled source clauses are admitted exactly
as native firings with their resource and original-cut queries
(\lean{OccurrenceCorePlacement.Native.admitted\_iff\_native}). Contextual guarded firings, which
carry the $\lambda$ and quoted-$\lambda$ executions, are admitted exactly as query-free native firings
(\lean{OccurrenceCorePlacement.Native.guarded\_valid\_iff\_native}).

\paragraph{Inference hosts.} Steps of schema processes are exactly authorized first-order firings
(\lean{OccurrenceCorePlacement.schema\_step\_iff\_admitted}); validity lives in the authority
restriction and in the derivation adequacy of the host, never in admission alone.

\paragraph{What programs still lack.} The two program hosts are joined by agreement, not by one
program core over the interface: a single core needs frames that can adjoin a schema's holes, an
operation \lean{Interface.Frames} does not have (55\%). At dependent frames, definitions by rows over
one package are checked row by row and extend the package, Church--Rosser and conservatively, without
a proof that the extension is qualified (Section~\ref{sec:addendum-definitions}); a program core with
a correspondence of runs there is not built (45\%). On the binder core: versions, which need a state of two occurrences with absence scoped to
the call (45\%); complete \mt{Atom} positions on the step, since whether a position is frozen
depends on heads a solution may still choose (40\%); call by need (45\%); runs from the reference
to the step for non-determinate programs (40\%). Call by need with open data and two-way unification
(55\%); answers up to $\eta$ in matching (50\%); fairness and cost.

\paragraph{Necessary extensions.} A rule generator for atomic blocks; the authority layer for
publication; closed asks for checked holes; a born ledger for identities
(\lean{Interface.nativeUnborn}); frames with $\eta$; disequation frames for disequality tells.
Each is placed above or inside the same step.

\section{Recommendation}\label{sec:recommendation}

\subsection{Retain}

Retain the rule of Definition~\ref{def:rule} as the operational core of the calculus: one firing
that takes, reads and makes occurrences, tells equations, and asks its guards and predicate queries
at every solution of its own information, over a frame action, with the interface of
Definition~\ref{def:step} as its abstraction. Justification: every proved row of
Section~\ref{sec:map} is a theorem about an instance of the interface; the admission
characterization is the same theorem in every guarded instance; each feature holds at the interface
or on a named restriction with a control and a bridge, or is a named gap
(Section~\ref{sec:joined}); over a package whose names are all rigid, admission of the first-order
and second-order steps transfers to the dependent step through the fibre wherever the unifier
decides, and that condition is inhabited on both sides. The interface alone constrains little: its
laws are that admission implies preparation, that knowledge refines and that the space changes by
the firing's commit (\lean{RuleSyntax.realizable\_iff}). The rule has one definition as syntax, the
step of a program (Definition~\ref{def:program-step}), with admission computed, and unchanged by changes
of the present rows outside the firing's regions when the opening stays fresh; restricted to the
fresh openings of a program, the first-order and dependent full steps have that step's transition
relation, under the restrictions stated after the definition.
The first rows of Section~\ref{sec:open} list what the retained rule does not yet have: that
definition at every instance, the hosted source language beyond its first fragment
(Section~\ref{sec:addendum-hosting}), placements of the dialects by the step's own
matching, an equivalence of hosted
programs that keeps their bags of answers, and types generated from the calculus's own presentation.

\paragraph{An assumption that was not compared: the size of a firing.} That one firing commits
\emph{one} match of one rule is how the rule of Definition~\ref{def:rule} and the step of a program are
written. It is not a result of this report, and no alternative was compared with it. An earlier
formalization of the transform has the quantifier in the firing itself: one row, every row at once, the
unique row (\lean{TransformQuantifiers.Footprints.EachStep},
\lean{TransformQuantifiers.Footprints.EveryStep}); there, committing every row has one result
(\lean{TransformQuantifiers.Footprints.everyStep\_functional}) and is neither the box nor the diamond of
committing one (\lean{TransformQuantifiers.Controls.every\_not\_box\_each},
\lean{TransformQuantifiers.Controls.every\_not\_dia\_each}). Three results of this report are costs of the
one-match firing: a bounded rule family does not commit a larger block in one firing
(\lean{OccurrenceCorePlacement.no\_bounded\_family\_commits}, \lean{Program.no\_fixed\_atomic\_commit}); no
finite program rewrites ``every copy at once'' over a store of any size
(\lean{RuleSyntax.Representation.allAtOnce\_not\_finite}); and an \mt{exec} is one firing only of a rule
built from its effect, whose admission does not check that every match was used
(\lean{DialectsAsStrategies.MM2.Controls.incomplete\_matches\_admitted}). The candidates (one match; every
match of a rule; any bag of matches that can fire together; a maximal such bag) have not been compared on
this step. Until they are, the retention above is of the interface and of the admission characterization,
and not of the size of a firing; ``a rule computed per block'' (Section~\ref{sec:programs}) is a way to
host a block on the one-match step and not an argument for that step.

\subsection{Revise}

\begin{center}\small
\begin{longtable}{@{}L{0.32\textwidth}L{0.62\textwidth}@{}}
\toprule
\textbf{Revision} & \textbf{Forcing theorem} \\
\midrule
\endhead
transport up to dead names, with a partial move & \lean{Interface.NativeUpToDeadLaws.strict\_separation\_fails}, \lean{Interface.NativeResidualSeparation.residual\_not\_separated}; partiality: \lean{Interface.NativeDead.Controls.deadQuery\_realized} with \lean{Interface.NativeDead.Controls.carrying\_not\_surjective} \\
helper firings agree up to the exchange of their telescopes & \lean{Interface.DependentFull.TelescopeControls.composites\_differ}, \lean{Interface.DependentFull.SwapControls.composites\_differ} \\
waves need joint solvability & \lean{Interface.NativeIteratedControls.pairwise\_not\_enough}, \lean{Interface.NativeIterated.no\_order\_of\_not\_jointly}, \lean{AnnotatedRuns.CellControls.hollow\_cube} \\
principal pattern answers need $\eta$ & \lean{Interface.Translation.pattern\_no\_principal}, \lean{ContextualDependentFibre.Controls.point\_pattern\_no\_principal} \\
ordinary and reflection queries split & \lean{Semantic.SemanticControls.syntactic\_collision}, \lean{Program.rawQuoted\_collision}, \lean{ProgramSourceReflection.reflect\_collision} \\
atomic blocks need computed rules or a compound occurrence & \lean{OccurrenceCorePlacement.no\_bounded\_family\_commits}, \lean{Program.no\_fixed\_atomic\_commit} \\
an interpreter allocates fresh variables in administrative firings & \lean{ProgramCompiled.no\_fixed\_one\_firing\_interpreter} \\
freshness needs the born ledger & \lean{Interface.NativeFreshnessControls.maker\_not\_realized}, \lean{Interface.native\_run\_protocol\_iff} \\
consistency and inhabitation are two asks & \lean{Interface.Dependent.trap\_helper} (relative to an empty type), \lean{ContextualDependentClosedBoundary.Controls.point\_not\_closedSolvable} with \lean{ContextualDependentClosedBoundary.Controls.ground\_uninhabited} \\
publication is an authority restriction & \lean{OccurrenceCorePlacement.forged\_admitted\_but\_not\_published} \\
reusable outputs are snapshots & \lean{Interface.NativeSnapshot.snapshot\_obstruction}; they commute under conservativity (\lean{Interface.NativeSnapshot.snapshots\_commute}) \\
open systems need labels for the context & \lean{Observation.CongruenceControls.closed\_not\_congruence} \\
disequality is a frame extension, told only on restrictions & \lean{DisequalityFrames.Restriction.guess\_breaks\_frame}; principal answers in constrained frames: \lean{Disequality.Controls.no\_presenting\_answer} \\
independence across spaces is a restriction, since the knowledge is shared & \lean{Spaces.Controls.shared\_variable\_crosses\_spaces} \\
rights alone do not bound writes; a ticket per write does & \lean{Permissions.Controls.rights\_alone\_unbounded}, \lean{Permissions.made\_le\_tickets} \\
contention includes inhibition & \lean{Contention.Controls.inhibition\_control}, \lean{Contention.Controls.forcing\_control} \\
a ``never'' guard is certified by a checked certificate, which exists exactly on finite-state data & \lean{GuardFragments.exists\_invariant\_iff}, \lean{GuardFragments.suspended\_of\_infinite}, \lean{GuardFragmentsOpen.Controls.always\_not\_stable} \\
reachability does not decide fair completeness; the identity of searches is branching & \lean{FairSearch.Controls.reachability\_does\_not\_decide\_fair\_completeness}, \lean{FairSearchIdentity.search\_identity\_is\_branching} \\
the dependent unifier needs more rules than the earlier set & \lean{ContextualDependentFragment.Controls.fragment\_not\_decisive}, \lean{ContextualDependentFragment.Controls.blocked\_not\_decisive} \\
\bottomrule
\end{longtable}
\end{center}

\subsection{Reject}

Each hypothesis below is rejected by an obstruction under the stated adequacy criterion.

\begin{enumerate}
\item \emph{Native diamonds close up to isomorphisms of states.} Criterion: residual solutions agreeing
  along both carrying morphisms agree. Obstruction: \lean{Interface.NativeResidualSeparation.residual\_not\_separated}.
\item \emph{Transport up to dead names is total on firings.} Criterion: every admitted firing at the
  larger state is carried from the smaller. Obstruction: \lean{Interface.NativeDead.Controls.deadQuery\_realized}, \lean{Interface.NativeDead.Controls.carrying\_not\_surjective}.
\item \emph{Pairwise independence suffices for a wave.} Criterion: some order of the family runs.
  Obstruction: \lean{Interface.NativeIteratedControls.pairwise\_not\_enough}.
\item \emph{Second-order frames without $\eta$, raised into dependent frames, keep principal pattern
  answers.} Criterion: a principal answer in the dependent frames, over a package with no root
  computation and symmetric head equality. Obstruction: \lean{Interface.Translation.pattern\_no\_principal}.
\item \emph{Raw label queries are ordinary.} Criterion: configurations equivalent modulo the knowledge
  admit renamed firings together. Obstruction: \lean{Semantic.SemanticControls.syntactic\_collision}.
\item \emph{A rule family with a uniform bound on made rows commits every atomic block in one firing.}
  Criterion: one block, one firing adding the block's rows. Obstruction:
  \lean{OccurrenceCorePlacement.no\_bounded\_family\_commits}. A family with no uniform bound, such as
  one rule per block size, and a block packaged as one compound occurrence are not excluded.
\item \emph{Absence from the space is freshness.} Criterion: the commit protocol realizes the run.
  Obstruction: \lean{Interface.NativeFreshnessControls.maker\_not\_realized}.
\item \emph{Symbolic admission certifies a checked hole.} Criterion: a closed solution exists.
  Obstruction: \lean{ContextualDependentClosedBoundary.Controls.point\_not\_closedSolvable}.
\item \emph{Admission confers publication of a judgment.} Criterion: the process's step reaches the
  result. Obstruction: \lean{OccurrenceCorePlacement.forged\_admitted\_but\_not\_published}.
\item \emph{The literal prototype step simulates the published register semantics with identical
  outputs.} Obstruction: \lean{Mettapedia.Languages.MeTTa.OSLFCore.MetaMeTTa.Controls.no\_identity\_output\_simulation};
  a different encoding is not excluded.
\item \emph{The earlier rule set of the dependent unifier decides the dynamic pattern fragment.}
  Criterion: every fragment store is solved principally or refuted at some budget. Obstruction:
  \lean{ContextualDependentFragment.Controls.fragment\_not\_decisive},
  \lean{ContextualDependentFragment.Controls.blocked\_not\_decisive}. The question for the current
  solver is open.
\item \emph{A fixed rule family interprets every call with variables in one firing.} Criterion: one
  firing makes the answer row. Obstruction: \lean{ProgramCompiled.no\_fixed\_one\_firing\_interpreter}.
\item \emph{Firings at different spaces are independent.} Criterion: a firing at one space leaves
  admission at another unchanged. Obstruction: \lean{Spaces.Controls.shared\_variable\_crosses\_spaces}.
\item \emph{A rights table alone bounds what a space receives.} Criterion: a bound on the rows made by
  permitted runs. Obstruction: \lean{Permissions.Controls.rights\_alone\_unbounded}.
\item \emph{Contention is the sharing of an occurrence.} Criterion: local normalization gives the
  step's probabilities, of total mass $1$. Obstruction: \lean{Contention.Controls.inhibition\_control}.
\item \emph{No context forces a race over positive weights,} with predicate queries. Criterion: added
  atoms leave every competitor below probability $1$. Obstruction:
  \lean{Contention.Controls.forcing\_control}; without inhibition the hypothesis holds
  (\lean{Contention.no\_context\_forces}).
\item \emph{Reachability of proofs decides completeness of fair search.} Criterion: two searches with
  the same reachable stores are both fairly complete or both not. Obstruction:
  \lean{FairSearch.Controls.reachability\_does\_not\_decide\_fair\_completeness}.
\item \emph{A guard may be priced by the store it scans.} Criterion: the price of a history is
  invariant under its interleavings. Obstruction: \lean{MeteredGuards.Controls.scan\_price\_order}.
\end{enumerate}

\subsection{What stays open}\label{sec:open}

The first five rows are what the calculus needs before it can be said to host a higher-order MeTTa.
The first is proved for two instances under the restriction its row states; the other four are not
proved, and Section~\ref{sec:addendum-nov13} narrows the third and the fourth. The other rows refine
rows of the verdict.

\begin{center}\small
\begin{longtable}{@{}L{0.33\textwidth}L{0.09\textwidth}L{0.52\textwidth}@{}}
\toprule
\textbf{Open row} & \textbf{Est.} & \textbf{What is missing, and the effect on the verdict} \\
\midrule
\endhead
the rule as syntax at every instance & 70\% & the step of a program is defined with admission derived, and the first-order and dependent full steps, restricted to the fresh openings of a program, have its transition relation under the restrictions stated after Definition~\ref{def:program-step}. A program of two rules whose takes have types that depend on terms has a run on the restricted dependent full step (\lean{RuleSyntax.Typing.Programs.dependent\_run}). Missing: predicate queries with arbitrary regions and answers; the helper step as a program whose values are carried by rows; the converting and templating steps; adequacy of the modal logic on the step of a program \\
one hosted source language with a whole-run correspondence & 55\% & one language with binders, equations with second-order patterns, types and space operations, translated into rules of the step with a correspondence of whole runs in both directions. Proved so far (Section~\ref{sec:addendum-hosting}): the language is defined, and its fragment of values, calls, sequencing, choice and failure is translated, one rule for each equation and six fixed rules, with runs corresponding both ways and the answers of a configuration read as a bag on a frontier of reached states. Missing: spaces, types, need and text read as code on the step; the bag of answers from the step alone beyond the first fragment and where a rule is repeated (Section~\ref{sec:addendum-bags}); the correspondence beside rows and rules that are not the program's own; a call through a function variable that evaluates its arguments. The program cores of Section~\ref{sec:programs} each have their own source language and correspondence, and two of them are joined by agreement \\
placements of the dialects and of rho in which the step's own matching does the work & 50\%, 85\% & the model of HE is placed by one ground rule per machine transition, with matching a premise of the transition (Section~\ref{sec:addendum-nov12}), and rho with ordinary substitution by one rule instance per redex (Section~\ref{sec:map}). Rho with semantic substitution, name equivalence and reduction modulo the equations is placed on the first-order step by three rule schemas, weakly bisimilar on all closed processes, with the channels matched by the knowledge and the contractum made by a term former whose meaning is the calculus's substitution (Section~\ref{sec:addendum-rho}), and on the step of a program by four rules whose communication rule makes the contractum by ordinary instantiation, for the carrier with sealed quotes, with both directions under one criterion (Section~\ref{sec:addendum-both-directions}). Missing: one schema per rule of a dialect's specification, with matching done by the knowledge; for rho, quotes that hold a bound name, the renaming clause of the encoding criteria, and a comparison on finite programs up to the target's own bisimilarity \\
an equivalence of hosted programs that keeps their bags of answers & 40\% & an equivalence preserved by every context exists and is characterized at one first-order frame (Section~\ref{sec:addendum-congruence}): a reduction-closed barbed congruence for contexts that add rules, tell equations, extend the frame, add rows or run beside another system. It is strong, it makes every opened unknown public, and it does not count, so congruent translations of hosted programs can have different bags. Missing: an equivalence that is sensitive to multiplicity and gives unknowns a scope; the characterization at frames with function symbols, with binders and with types, and in a weak form; a finite check for the equivalence of two programs. Earlier data: Section~\ref{sec:addendum-program-bisimilarity}, Section~\ref{sec:map}, Section~\ref{sec:addendum-nov11} \\
types generated from the calculus's own presentation, with a substitutability theorem & 50\% & a presentation of the step by sorts, constructors and rewrites, from which OSLF would generate spatial and behavioural types without asking the step, and a theorem that bisimilar states are exactly those with the same types. The types proved so far are modal types over the step (Sections~\ref{sec:addendum-nov12} and~\ref{sec:addendum-types-classes}) \\
decisiveness of the dependent unifier on the whole dynamic pattern fragment & --- & \emph{false as stated} under silent universe inclusion (Section~\ref{sec:addendum-levels}). The tree's own driver now decides several fragments by shape (Section~\ref{sec:addendum-levels}). What stays open is narrower: a class theorem for Miller patterns mixed with the join (60\%), initiality from phases whose answers keep unknowns declared (65\%), which stores the phased solver decides (60\%), and level variables inside the tree's own type theory (55\%). The rule is unaffected, since residuals stay pending without loss \\
pruning under an older metavariable & 50\% & base-type instantiations keep the normal form of the instantiated term; outside the fragment \\
census queries and multi-frame queries through the join & 60\%, 65\% & a meaning for arbitrary regions at dependent solutions; absence and counting through the multi-frame instance \\
the endomorphism condition for the first-order instance & 45\% & endomorphic answers need inhabitants of the dropped metavariables' types, which can be empty \\
principality against arbitrary dependent telescopes & 40\% & second-order principal answers against solutions outside the image; admission does not need it \\
step bisimilarity for helper firings; the dead-helper converse & 60\%, 45\% & a strict transport; strengthening of typed judgments and query tests natural under weakening \\
located labels at dependent frames & 70\% & query regions that respect spaces, removing the remaining hypotheses of the dependent spaces theorem \\
guarded native firings with disequations; a label law for native rows; the four-way family on the native step & 55\%, 60\%, 65\% & checks asked at separating solutions only; term labels through publication \\
one program core over the interface; programs at dependent frames & 55\%, 45\% & frames with schema extensions; at dependent frames, definitions by rows over one package are checked and extend it, Church--Rosser and conservatively (Section~\ref{sec:addendum-definitions}); missing there: a program core with a correspondence of runs, a row checker over a presented package, and the qualification of a package extended by an admitted definition \\
the binder program core: versions, complete \mt{Atom} positions, call by need, non-determinate runs & 45\%, 40\%, 45\%, 40\% & see Section~\ref{sec:programs}; no effect on retain or reject \\
results of contextual guarded executions against native published states & 70\% & admission is placed; results are not. Would not change the verdict \\
native transport for recipes, reusable outputs and closure queries & 50\% & refines the revisions (patterns as queries; conservative reuse); a closure query does not determine its residual (\lean{Interface.NativeClosingControls.closure\_extensions\_disagree}) \\
static full abstraction (equivalence modulo the knowledge against probes) and congruence, for contextual and native firings & 50\% & would strengthen observation; no effect on retain or reject \\
causal Bayesian networks on the step & 70\% & the confusion-free race identity for cells enabled by joins, then the staged network program (65\% after it) and replay (60\%) \\
``never'' guards on infinite-state data; fairness in general & 60\%, 70\% & search firings discharging a symbolic closure obligation; strong fairness with ``a proof stays reachable'' giving completeness on finite-state systems \\
grants as firings & 70\% & delegation, the sealed case and revocation; the permission vocabulary is unsettled \\
a checker generated from the calculus's own rules & 50\% & \texttt{checkProof = accept} exactly for derivations, generated from admitted rules; located kernel adequacy covers MM0 only \\
behavioural modalities over the rule & 50\% & OSLF's $\Diamond$ and $\Box$ over the firings of every \lean{Interface.Step} carry the types closed under bisimilarity; without a finite program, finitary formulas give those types for some observation sets and not for others (Section~\ref{sec:addendum-types-classes}); open: fixed points with negation or alternation under a budget (Section~\ref{sec:addendum-when-types}) \\
HE semantics and a uniform register-semantics encoding & 50\% & comparisons; no effect on the rule \\
\bottomrule
\end{longtable}
\end{center}

\section{Addendum: results after the revision}\label{sec:addendum}

Two results arrived after the revision above. They change no row of the recommendation; they
narrow two open rows.

\subsection{Firing of the pattern rules}
The open-unification row of Section~\ref{sec:map} named firing on Miller-pattern equations as the
main gap. The pattern rules now provably make progress, at a large enough budget, on three rows of
the census:
\begin{itemize}
\item \textbf{inversion} for the newest metavariable, $F\,\bar x \doteq t$ with $t$ over $\bar x$
  and older declarations (\lean{ContextualDependentPatternFiring.firingRound\_fires\_inversion});
\item \textbf{inversion in place}, for an older metavariable solved past the later declarations
  (\lean{ContextualDependentPatternFiring.placedRound\_fires});
\item \textbf{refutation by a rigid local outside $\bar x$}, even with flexible spines inside $t$
  (\lean{ContextualDependentPatternFiring.firingRound\_fires},
  \lean{ContextualDependentPatternFiring.refuted\_rigidOccurrence}).
\end{itemize}
The firing theorems are instantiated on tower stores
(\lean{ContextualDependentPatternFiring.Controls.inversion\_fires},
\lean{ContextualDependentPatternFiring.Controls.placed\_fires},
\lean{ContextualDependentPatternFiring.Controls.rigid\_fires}); the typed refutation
\lean{ContextualDependentPatternFiring.refuted\_rigidOccurrence} is applied to no store in the tree. One reader of elimination spines,
along the literal $\Pi$/$\Sigma$ types of the prefixes, reads the nested and applied projections
(\lean{ContextualDependentEliminations.elim\_exact},
\lean{ContextualDependentEliminations.elim\_weight}), and the stratified class is closed under an
answer for the newest metavariable
(\lean{ContextualDependentStratifiedClosure.stratified\_instantiate\_newest}). Each new layer acts
only where the one below it is stuck, so every earlier decision is kept
(\lean{ContextualDependentEliminations.elims\_decided\_of\_placed}).

\emph{Still open.} Four flexible rows (intersection, pruning under rigid structure, flex--flex with
a declaration between, and the cumulative argument) and the spine at a flexible base type are in
the fragment, solvable, and stuck at budget 16
(\lean{ContextualDependentPatternFiring.Controls.unread\_flexible\_rows},
\lean{ContextualDependentEliminations.Controls.flexibleBase\_elims\_suspended}). The four flexible
rows need a rank under which a step that keeps every declaration still decreases; closure under
in-place answers is open. Decisiveness on the fragment, as stated, is false
(Section~\ref{sec:addendum-levels}).

\subsection{Causal networks on the step}
The composition row now carries Pearl's calculus. When every race, once open, persists until it
fires, normalizing each choice within its own race gives exactly the step's trace probabilities
summed over orders, also through joins
(\lean{CausalCells.CellSystem.History.race\_identity}). A causal Bayesian network is then a
program: each variable's schema takes \mt{(pend v)}, reads \mt{(reg v r)} and its parents'
\mt{(val p b)}, and makes \mt{(val v b)}, weighted by its table; its likelihood is the network's
joint (\lean{CausalNetworksStep.step\_likelihood\_eq\_jointWeight}). A restriction at the
variable's handle gives the truncated factorization under positivity
(\lean{CausalNetworksStep.step\_do\_likelihood}); telling a forcing regime atom gives it always,
and equals conditioning on the regime (\lean{Mettapedia.GSLT.Causality.StepBridge.step\_regime\_likelihood},
\lean{Mettapedia.GSLT.Causality.StepBridge.regime\_is\_conditioning},
\lean{Mettapedia.GSLT.Causality.StepBridge.do\_is\_regime}). Replaying recorded noise, with
abduction over the rows the run did not visit, gives the counterfactual relative to the network's
noise law (\lean{CausalReplay.replay\_likelihood}, \lean{CausalReplay.replay\_average},
\lean{CausalReplay.extends\_record\_iff}). It replays the exogenous inputs, never the previous
winners. The network alone does not determine counterfactuals, so they are stated relative to that
law. The network's graph is read off the step's independence
(\lean{CausalSeparation.history\_staticIndependent\_iff}).

\subsection{The C runtime and proof-system integration evidence}
The standalone \texttt{tcalc} runtime executes rule tables used by Micromath,
MM0/MMB, HOL Light proof replay, Megalodon and Dedukti implementations.
Micromath's differential operation suite is not a \texttt{set.mm} verification.
HOLR records are primitive proof nodes in a local format; their counts are not
counts of distinct exported mathematical theorems.

The Megalodon implementation checks explicit natural-deduction proof terms,
with additional typing and signature rules. The Dedukti implementation has
explicit LF/CoC and beta/eta profiles. Named C conversion, substitution and
context helpers remain in their trusted bases. Reference/corpus agreement and
mutated-proof rejection are regression evidence, not general adequacy or
conservativity theorems. Different source-elaboration and proof-replay timings
are not matched-input kernel speed comparisons.

The prepared-effect/admission/event C bridge follows the mathematical layering.
Its source-state guard evidence and cursor-resumption controls are relevant
realization tests, but the complete C-to-Lean correspondence and the live CeTTa
cursor integration remain open. Runtime permission records must preserve the
original event and observation contract; a new profile cannot manufacture
missing evidence.

The exact 211-byte \texttt{flat\_fire.S} reference function is proved by
\texttt{FLAT\_FIRE\_CORRECT} in HOL Light over s2n-bignum's x86 model, with no
new axioms. The theorem's memory/ABI assumptions include scratch separation,
valid modular 32-bit offsets and the 56-byte stack/return exclusion. The
pointer-based runtime packing adapter is not covered by that theorem.

The Lean declarations cited in this addendum were audited separately: all use
\lean{propext}, \lean{Classical.choice} and \lean{Quot.sound}, or fewer, and none uses
\lean{sorryAx}.

\section{Addendum: results of 11 November}\label{sec:addendum-nov11}

These results correct one row and add theory the core needed. They change the recommendation
in one place: decisiveness of the dependent unifier on the whole fragment is no longer an open
goal but a proved obstruction, with a replacement.

\subsection{Unification with level variables}\label{sec:addendum-levels}
The earlier promise, decision on the whole dynamic pattern fragment, is false when universes are
included silently. The fragment admits a solvable store with no most general answer, and its
missing principal answer is exactly a level constraint. In a cumulative model with level
variables, at a valuation with $a\le b$ that store has a principal answer exactly when also
$b\le a$, that is, when $a=b$ (\lean{UnificationWithLevels.CumulativeCodes.Cross.principal\_iff},
\lean{UnificationWithLevels.CumulativeCodes.Cross.principal\_at\_equal\_levels}).

Where $a<b$, no answer is principal
(\lean{UnificationWithLevels.CumulativeCodes.Cross.no\_principal}), and every exact solver stops
there postponed (\lean{UnificationWithLevels.CumulativeCodes.Cross.forced\_postponement}). Knowing
only $a\le b$, the constraint $b\le a$ is neither entailed nor refuted
(\lean{UnificationWithLevels.CumulativeCodes.Cross.postponed\_level\_constraint}). On the tree's own
fibre, the cross store stops postponed
(\lean{UnificationWithLevels.DependentFibre.cross\_postponed}).

The replacement is a solver whose every step keeps the solution set exactly
(\lean{UnificationWithLevels.Solver.Reaches.exact}). A run stops in one of three ways:
\begin{itemize}
\item solved, with a principal answer (\lean{UnificationWithLevels.exact\_nil\_iff\_principal});
\item refuted, with its reason;
\item postponed, with its residual kept as knowledge. On a solvable store without a principal
  answer this outcome is forced (\lean{UnificationWithLevels.postponed\_of\_solvable\_no\_principal}).
\end{itemize}
Preparation holds exactly when the residual is solvable
(\lean{UnificationWithLevels.prepares\_iff\_residual}).

Decision holds where progress is proved by shape. In a second-order model with cumulative
universes and level variables the solver terminates
(\lean{UnificationWithLevels.CumulativeCodes.rules\_terminate}). On the uniform pattern
fragment, every store ends with a principal answer or no solution
(\lean{UnificationWithLevels.CumulativeCodes.principal\_or\_unsolvable}). Level constraints
over $\mathbb{N}$ are decided by a bounded search and by lowering
(\lean{UnificationWithLevels.LevelStore.satisfiable\_iff\_bounded},
\lean{UnificationWithLevels.lowerSat\_iff}). The lowering's count of constraint evaluations is
polynomial in the unary size of the store (\lean{UnificationWithLevels.lowerSat\_cost}).

The occurs check is its own null. $X\doteq\mathrm{arr}\,X\,\mathrm{gnd}$ has no well-founded
solution (\lean{UnificationWithLevels.OccursNull.well\_founded\_refutation}), and it has a
rational one (\lean{UnificationWithLevels.OccursNull.rational\_solution}), which is unique
(\lean{UnificationWithLevels.OccursNull.rational\_solution\_unique}).

\emph{Since then.} In the models of codes, compatible groups are decided by meet narrowing
(\lean{UnificationWithLevels.CumulativeCodes.Meet.decided}), and Miller patterns with pruning are
decided (\lean{UnificationWithLevels.MillerCodes.decided}). In the tree's own fibre, the tree's own
driver decides, from some budget on, families and fitting groups
(\lean{UnificationWithLevels.TreeFibre.family\_decided},
\lean{UnificationWithLevels.TreeFibre.fit\_decided}), the true meet and several parameters
(\lean{UnificationWithLevels.TreeFibre.meetGroup\_decided},
\lean{UnificationWithLevels.TreeFibre.meetGroupK\_decided}), flex--rigid stores
(\lean{UnificationWithLevels.TreeFibre.flexRigid\_decided}) and every equation between arrow codes
(\lean{UnificationWithLevels.TreeFibre.arrows\_decided}). Every normal type has a least universe
(\lean{UnificationWithLevels.TreeFibre.least\_universe}), which gives an exact
greatest-lower-bound rule (\lean{UnificationWithLevels.TreeFibre.glbMeet\_exact}).

Miller patterns with arbitrary distinct locals are decided by the tree's own driver
(\lean{UnificationWithLevels.TreeFibre.Miller.miller\_decided}). This also holds when their codomains are
incomparable, by lowering both unknowns to a computed greatest lower bound
(\lean{UnificationWithLevels.TreeFibre.Miller.millerGlb\_decided}). The greatest-lower-bound rule runs
inside a decided fragment (\lean{UnificationWithLevels.TreeFibre.glbGroup\_decided}), and
flex--rigid stores need no chain of universe declarations
(\lean{UnificationWithLevels.TreeFibre.flexRigidGlb\_decided}). A normal term typed over two
telescopes is typed over their join (\lean{UnificationWithLevels.TreeFibre.typed\_join}). This
decides groups whose domains need not be comparable
(\lean{UnificationWithLevels.TreeFibre.joinMeetK\_decided}).

A single exact rule lowers codomains to a greatest lower bound and widens domains together
(\lean{UnificationWithLevels.TreeFibre.glbJoinK\_exact}). Its condition is stated with the order
alone. Wherever both the identity and a constant function solve the equation, the condition is
also necessary for a principal answer
(\lean{UnificationWithLevels.TreeFibre.no\_principal\_of\_not\_crossAt}). Unknowns at
codomain positions are decided (\lean{UnificationWithLevels.TreeFibre.arrowsCod\_decided}).

One solver decides the union of these fragments
(\lean{UnificationWithLevels.TreeFibre.union\_decided}); a store it cannot place stays pending,
never refuted. Every solvable store in the union, over a telescope of universe families, has an
initial solution (\lean{UnificationWithLevels.TreeFibre.union\_initial}), and the answers become
initial once solved unknowns are deleted
(\lean{UnificationWithLevels.TreeFibre.glbGroup\_initial},
\lean{UnificationWithLevels.TreeFibre.Miller.millerGlb\_initial}).

Under one rank that every rule decreases, a single solver also decides stores that mix these
fragments (\lean{UnificationWithLevels.TreeFibre.one\_decided}). A codomain that uses its argument
is decided by a computed case split that is complete on its class
(\lean{UnificationWithLevels.TreeFibre.depDecide\_complete}).

These runs need no round budget. Each solver's run is defined by well-founded recursion on its
rank, and agrees with every budgeted run that finishes
(\lean{UnificationWithLevels.TreeFibre.drive\_eq\_driveWF}). The mixed solver decides its
fragment without any budget (\lean{UnificationWithLevels.TreeFibre.one\_decidedWF}). When each
phase's answers cover their frame, stores decided in phases have initial solutions whatever their
telescope (\lean{UnificationWithLevels.TreeFibre.phases\_initial}). The cross store stays postponed
throughout.

\emph{Still open:} a class theorem for Miller patterns mixed with the greatest-lower-bound join
(three exact rank-lowering rules decide the tested stores, but no solver yet decides the class);
initiality from phases whose answers keep unknowns declared; which stores the phased solver
decides; level variables inside the tree's own type theory.

\subsection{Names}
A step may carry a ledger of issued names. A firing that issues names is admitted only when
they are new and every name it writes was issued, earlier or by the firing itself. On any step
whose results mention only present or written names, every admitted firing preserves the
invariant that every name in the store and the knowledge was issued
(\lean{NamesOnTheStep.invariant\_result}), so every run keeps it
(\lean{NamesOnTheStep.invariant\_last}). An issued name is new to everything
(\lean{NamesOnTheStep.issued\_new}), and along a run no name is issued twice
(\lean{NamesOnTheStep.issuedAlong\_nodup}).

Both conditions are necessary, as controls show, and nothing more is required. Structural
codes, which quote terms, are a separate constructor (\lean{NamesOnTheStep.quote\_ne\_name})
and need no issuing. A theory whose rules mention only codes loads into any store without
issuing and cannot forge a name
(\lean{NamesOnTheStep.CodeOnlyLoading.loading\_needs\_no\_issuing},
\lean{NamesOnTheStep.CodeOnlyLoading.no\_forgery}). Micromath is such a theory: every theorem
stored by a run that mentions no issued name is valid in every model
(\lean{NamesOnTheStep.MicromathSoundness.mm\_valid}).

\subsection{Nulls and guards over stored schemas}
Seven kinds of emptiness are defined on the step, alongside refusal, unresolved work and an
exhausted budget. The pairs that matter here are separated on the first-order step; not every
pair has a theorem. A guard over a stored schema reads the schema as all its instances and has
three outcomes. No verdict is unresolved work, never an admission or a refusal
(\lean{NullsOnTheStep.StoredSchemas.noVerdict\_both\_ways}), and a better evaluator only
resolves unresolved attempts (\lean{NullsOnTheStep.ThreeValued.refines\_keeps\_attempt}).

\subsection{Absence and parallel composition}
Consider a firing with takes, reads, makes, tells and equality checks, and no predicate
queries. It stays admitted beside a store that shares no made identity exactly when its
equations have a solution that also solves that store's knowledge
(\lean{NullsOnTheStep.Composition.positive\_admission\_iff}). Absence-guarded admission
survives exactly when the added store has no match
(\lean{NullsOnTheStep.Composition.guarded\_union\_iff}).

No relation that a parallel component can switch off is the reduction of a rho rule set, that is,
of rewrite schemas on rho's processes applied in every parallel context, as communication and drop
with ordinary substitution are (\lean{NullsOnTheStep.Composition.rho\_step\_not\_join\_sensitive}).
In particular, the absence step on rho processes is not
(\lean{NullsOnTheStep.Composition.rho\_step\_not\_absence}). So absence is not the reduction of such
a rule set on the same processes. Rho with semantic substitution is not covered by these two
theorems. The statement for every presentation with a parallel context rule, and a separation for
translations homomorphic on parallel composition, which a map defined on whole configurations need
not be, are in Section~\ref{sec:addendum-separation}, with the tree's earlier encoding results for bag
rules into pattern-level rho. On an owned store seen whole at each firing, absence is meaningful.

The union of stores is the parallel composition:
\begin{itemize}
\item congruence of terms is heating (\lean{StoreComposition.Term.congruent\_iff\_heat});
\item reduction is a bisimulation (\lean{StoreComposition.reduces\_bisimulation});
\item rho's processes, with communication and drop by ordinary substitution, are represented on
  the step, under one rule instance per source redex, by a functional bisimulation
  (\lean{RhoStoreBisimulation.represents\_bisimulation}), which needs heating
  (\lean{RhoStoreBisimulation.heating\_is\_needed}); rho with semantic substitution is not covered
  here, and its placement by three schemas is in Section~\ref{sec:addendum-rho}.
\end{itemize}
An authored parallel operator that had been added to the calculus's syntax duplicated the union
of stores and was removed.

\subsection{Runs and weights}
A word of prepared firings is a run of the permitted step exactly when each of its firings is
permitted (\lean{Interface.KnowledgeTensorCore.exists\_run\_iff\_permitted}). The weight of a run
is determined by the weights of single firings: any functor from runs to the weight monoid that
agrees with them on firings is the weight functor
(\lean{Interface.KnowledgeTensorCore.weightFunctor\_unique}). Absence needs a complement of sums
in the weights. Boolean and independent probability weights have one; nontrivial commutative
rings have none (\lean{Mettapedia.Logic.LinearLogic.KnowledgeTensorWeight.SumComplement.ring\_isEmpty}).
Separation of runs does not make their survival events independent
(\lean{Interface.KnowledgeTensorCore.SeparatedControls.separated\_yet\_correlated}).

\subsection{The C runtime}
The standalone runtime of the previous addendum was a prototype, and it is retired as an
implementation. A reference implementation written from the Lean statement decides admission
literally: guards must hold at every solution, decided by classifying each occurrence against
the most general unifier, and a refusal changes nothing. Its specializer generates code that
agreed with it on every decision tested, and Metamath Zero runs as rules of the step. These are
differential results, not theorems. The proved 211-byte leaf is unchanged.

Dedukti's own matcher differed from its specification on a pattern variable used twice under
binders. The difference was reported upstream with a fix.

The Lean declarations cited in this addendum use \lean{propext}, \lean{Classical.choice} and
\lean{Quot.sound}, or fewer; none uses \lean{sorryAx}.

\section{Addendum: results of 12 November}\label{sec:addendum-nov12}
These results were proved on 12 November, after an adversarial audit of the previous day's work.
The audit's corrections are already folded into the sections above.

\subsection{Logic over the step}
OSLF's modalities are defined over the step's admitted firings: possibility and necessity over
successors, and necessity over predecessors. OSLF's generation of a logic from a presentation is
applied to an authored event presentation with one rewrite, whose premises ask the step's own
preparation and admission. For any step with an exact code of its states and firings, the
possibility that presentation generates, read after choosing a firing, is the step's possibility
(\lean{LogicOverTheStep.PresentationMeetsStep.diamond\_meets\_step}). On the first-order step the
two meet through query-free firings. When substitutions are infinite, no injective code from
first-order firings into patterns exists, because query regions are functions
(\lean{LogicOverTheStep.PresentationMeetsStep.no\_pattern\_code}). A presentation of the calculus
itself, from which OSLF would generate its types without asking the step, is open
(Section~\ref{sec:open}).

Pulling predicates back along a map of transition systems commutes with possibility exactly for
weak-pullback squares (\lean{LogicOverTheStep.ModalComparison.beckChevalley\_di\_iff}). Along an
open map it commutes with every modal and Boolean operation
(\lean{LogicOverTheStep.ModalComparison.pb\_modal\_hom}). Unlocked hyperpose is not such a map, and
a formula witnesses the failure (\lean{LogicOverTheStep.Narrowings.PeTTa.hyperpose\_step\_fails}).

Predecessors do not lift along the narrowing to rho's communication with ordinary substitution,
because every communication records its matches in the knowledge
(\lean{LogicOverTheStep.RhoPast.rho\_backZag\_fails}). Rho's predecessor
necessity is the store's, quantified over the heating class
(\lean{LogicOverTheStep.RhoPast.rho\_predecessorBox\_eq\_fibre}).

Consider the first-order step with finitely many rules, a supply of fresh identities, and
observations invariant under renaming. There, two states satisfy the same formulas exactly when
they are bisimilar (\lean{LogicOverTheStep.ProgramAdequacy.program\_adequacy}). Counting
modalities, which count firings as events, are adequate for graded bisimulation under the same
hypotheses (\lean{LogicOverTheStep.CountingModalities.program\_graded\_adequacy}). Absence and exact
counts survive an added store exactly when that store lacks the pattern
(\lean{LogicOverTheStep.ContextGrades.absent\_grade\_iff},
\lean{LogicOverTheStep.ContextGrades.count\_grade\_iff}).

\subsection{Runs as proofs}
Runs that record each firing's rule, frame and justification are the initial proof-relevant
model, and the earlier propositional initiality follows from it
(\lean{ProofRelevantRuns.run\_initial\_of\_relevant}). A teller and an asker interchange exactly
when the told fact is already known
(\lean{ProofRelevantRuns.InterchangeControls.teller\_asker\_cell\_iff}). Pairwise compatibility does
not give a joint wave (\lean{ProofRelevantRuns.hollow\_not\_flag}). Median-convex tells do give
one, by a Helly theorem (\lean{ProofRelevantRuns.helly}), which applies on the native step
(\lean{ProofRelevantRuns.every\_order\_admitted\_of\_convex}).

\subsection{Programming at dependent frames}
Programs over vectors, a typed interpreter, dependent pairs and identity types run on the
dependent step. Two limits are proved.
\begin{itemize}
\item A firing's knowledge does not type the rows it makes. A row typed at every solution of the
  firing's information is still refused, because the step settles along the identity
  (\lean{ProgrammingOnTheStep.Refinement.coerce\_refused},
  \lean{ProgrammingOnTheStep.Refinement.coerce\_typed\_at\_solutions}).
\item Functions used in types must compute in conversion. An equation that holds where addition
  computes is refuted where addition runs by firings
  (\lean{ProgrammingOnTheStep.ConversionContrast.Computed.computed\_solvable},
  \lean{ProgrammingOnTheStep.ConversionContrast.Fired.fired\_unsolvable}).
\end{itemize}
The first limit is an artifact of settling. Its core is forced, because typing modulo arbitrary
constraints is undecidable. The second limit is forced. A related consequence: the answer of a
call to a rigid function is not determined by the knowledge, and has to be read from the row the
firing made (\lean{LogicOverTheStep.RigidReadback.append\_call\_undetermined}).

Four tempting laws are refuted on the tree's own objects and kept as permanent tests: a principal
answer for every solvable store, framing in the presence of absence, pairwise compatibility as a
joint wave, and solvability as inhabitation
(\lean{ProgrammingOnTheStep.AcceptanceTests.principalLaw\_fails},
\lean{ProgrammingOnTheStep.AcceptanceTests.frameLaw\_fails},
\lean{ProgrammingOnTheStep.AcceptanceTests.waveLaw\_fails},
\lean{ProgrammingOnTheStep.AcceptanceTests.inhabitationLaw\_fails}).

\subsection{The models of HE and PeTTa, on the step and compared}
The statements of this subsection are about the Lean models of HE and PeTTa;
Section~\ref{sec:addendum-ledger} says which of them hold of the implementations.
The model of HE's retained-control machine is strongly bisimilar to the step restricted to one
ground rule per machine transition, with matching as a premise of the transition: the whole space
is one occurrence, the knowledge stays trivial, and the step's own unification does none of the
matching (\lean{DialectsAsStrategies.HEOnTheStep.he\_on\_the\_step}). With one occurrence per row, an HE
equation query and a PeTTa call request read the same rows exactly when their regions agree
(\lean{DialectsAsStrategies.TriangleFootprints.read\_word\_triangle}). Over a whole run from HE's
entry, HE reads exactly PeTTa's request rows, and the only difference is HE's absence guards over
declarations (\lean{DialectsAsStrategies.TriangleFootprints.WholeRunWords.whole\_run\_rows},
\lean{DialectsAsStrategies.TriangleFootprints.WholeRunWords.whole\_run\_guards}). An add into an
open query's region spoils the query
(\lean{DialectsAsStrategies.TriangleFootprints.add\_into\_region\_race}).

The results from here to the end of the subsection compare the two models directly; the step does
not occur in their statements. Consider programs of equations with symbol heads and no list views,
and a term whose computation settles, which is a condition on the program alone. On every complete
HE run such a term has the
same answer bag as in PeTTa, and a complete run exists
(\lean{DialectsAsStrategies.TriangleFootprints.EveryRun.answer\_bags\_every\_run}). The agreement
extends to bodies that use their variables, nested calls, recursion over compound data and integer
arithmetic
(\lean{DialectsAsStrategies.TriangleFootprints.GroundedEveryRun.Controls.double\_every\_run}).
Every HE run on such a term finishes
(\lean{DialectsAsStrategies.TriangleFootprints.EveryRunFinishesW.every\_run\_finishes\_kg\_with\_he\_lookups}).

The differences are named, each proved on every HE run:
\begin{itemize}
\item A call that no equation matches returns itself in HE and fails in PeTTa
  (\lean{DialectsAsStrategies.TriangleFootprints.EveryRun.uncovered\_every\_run}).
\item HE checks declared argument types and PeTTa does not
  (\lean{DialectsAsStrategies.TriangleFootprints.DeclaredEveryRun.declared\_every\_run}).
\item An \texttt{if} whose condition is neither \texttt{True} nor \texttt{False} returns the call in
  HE and takes the else branch in PeTTa
  (\lean{DialectsAsStrategies.TriangleFootprints.IfDifference.if\_non\_boolean\_differs}).
\item An equation for a name of PeTTa's primitive table: on \mt{(= (cons Yes No) Both)}, whose head
  HE's table of grounded operations lacks, every complete HE run publishes \mt{Both}; on
  \mt{(= (not Yes) No)}, whose head that table executes, every complete HE run publishes the call
  \mt{(not Yes)} itself, although the equation matches it. A complete run exists in both, and the
  PeTTa model stops at a fault with no value
  (\lean{DialectsAsStrategies.TriangleFootprints.PrimitiveNames.Controls.cons\_every\_run},
  \lean{DialectsAsStrategies.TriangleFootprints.PrimitiveNames.Controls.not\_declines\_every\_run}).
\item A free body variable gets a different name on different HE runs
  (\lean{DialectsAsStrategies.TriangleFootprints.FreeVariables.Controls.fresh\_not\_determined}).
\end{itemize}
SWI-PeTTa's lock runs its goal once and keeps one answer
(\lean{DialectsAsStrategies.PeTTaHE.Controls.once\_lock\_loses\_answers}).

\subsection{Dedukti on the reference}
The reference checker types a rule's right-hand side under its left-hand side's typing
constraints, used as hypotheses. This is sound modulo conversion for every solution, well-founded
or not. So its acceptance of \texttt{SR\_unsat.dk} is a policy difference from Dedukti, not a
soundness defect.

Against Dedukti's current master, 181 of the 203 test files agree, 20 are not hosted, and 2 differ
by a deliberate arity choice. A sweep of 4,472 mutants gives no wrong acceptance and no wrong
rejection. These are differential results, not theorems. The matcher fix reported in the previous
addendum was merged upstream.

The Lean declarations cited in this addendum use \lean{propext}, \lean{Classical.choice} and
\lean{Quot.sound}, or fewer; none uses \lean{sorryAx}.

\section{Addendum: further results of 12 November}\label{sec:addendum-nov12-evening}
This addendum settles when an equation set may compute in types, and checks a dependent program
against what its match teaches. It sets out the options for extensional equality, reads a type as a
class that no admitted observer can split, and says when such types apply. It also makes a budget
that runs out its own outcome, models a design space of evaluation modes, checks certificates for
behavioural distance, proves \mt{let} on every run of the model of HE, and says which statements of
Section~\ref{sec:addendum-nov12} about HE and PeTTa hold of the implementations. The
recommendation's retained, revised and rejected rows are unchanged; one open row is corrected.

\subsection{One equation sign, two readings}\label{sec:addendum-two-readings}
An equation set has two readings: operationally its rows fire, and definitionally they are
computing rules of the knowledge's conversion, so that types may compute with them. The set is
\emph{qualified} for the definitional reading (\lean{NativeTypesOnTheStep.Computing.Qualified}) when
its rows form a constructor system (linear constructor patterns, every variable of a right side on
its left, no two left sides overlapping), its right sides are typed, and every declared constant is
semantic for the typed and for the algorithmic equality, which structural recursion on a constructor
argument provides. A qualified package is Church--Rosser with $\beta$, its reductions preserve
typing, and the type-directed algorithm, with $\eta$ for functions and pairs, decides its typed
equality (\lean{NativeTypesOnTheStep.Computing.Qualified.churchRosser},
\lean{NativeTypesOnTheStep.Computing.Qualified.subjectReduction},
\lean{NativeTypesOnTheStep.Computing.Qualified.decide}); $\eta$ is decided by the type, since
$\eta$-reduction does not preserve typing in the cumulative tower
(\lean{NativeTypesOnTheStep.Computing.eta\_reduct\_not\_typed}). Three packages of the tree's type
theory are qualified with every field proved: an executable package whose addition computes in a
type family, the numbers with their recursor, and list accumulators
(\lean{NativeTypesOnTheStep.Computing.executable\_qualified},
\lean{NativeTypesOnTheStep.Computing.ValueType.numbers\_qualified},
\lean{NativeTypesOnTheStep.Computing.ListFunction.lists\_qualified}).

On a Church--Rosser package the two readings agree on the package's own reduction: a call reaches
$v$ by rewriting with the rows, at any position, exactly when $v$ is a conversion normal form of the
call (\lean{NativeTypesOnTheStep.Computing.answer\_iff}). That reduction is the package's term
rewriting, not runs of the step; its tie to firings of the full step is open (see the end of this
subsection). On a
qualified package that answer has the call's type
(\lean{NativeTypesOnTheStep.Computing.Qualified.answer\_typed}). Outside the class they disagree,
with a smallest counterexample for each of five causes
(\lean{NativeTypesOnTheStep.Computing.classification},
\lean{NativeTypesOnTheStep.Computing.RootDeterminism.junk\_disagreement},
\lean{NativeTypesOnTheStep.Computing.RootDeterminism.effect\_identifies}): overlapping rows give one
call two answers; a looping row gives no answer, yet its knowledge solves $t\doteq\mathrm{succ}\,t$
for the call, which the rigid package refutes; an ill-typed right side makes a call convertible with
a term to which it is equal at no type; a partial definition supports a refutation that the
completing row overturns; and rows that read the store agree with it at two states only through a
computation that is not root-deterministic and that, declared, identifies $0$ and $1$. The syntactic
part of the check is one Boolean pass over the rows, sound for Church--Rosser over a package with no
computation steps of its own and a symmetric head equality; it accepts a looping row, so termination
is a separate part (\lean{NativeTypesOnTheStep.Computing.DeclarationCheck.checked\_churchRosser},
\lean{NativeTypesOnTheStep.Computing.DeclarationCheck.loop\_check}). The templating step types each
made row after its match: on it every answer of an admitted firing has its declared type at every
solution of what the firing learned, and a mistyped answer is refused, for every \emph{settler} (the
choice of the solution a firing settles along)
(\lean{NativeTypesOnTheStep.Computing.AnswersTyped.templating\_answers\_typed},
\lean{NativeTypesOnTheStep.Computing.AnswersTyped.mistyped\_answer\_refused}). Without a termination
proof, conversion is checked with a budget and three outcomes: \emph{equal} is sound,
\emph{different} is sound on a Church--Rosser package, \emph{not decided} holds exactly when the
budget ran out on one side, and a decided verdict stays at every larger budget
(\lean{NativeTypesOnTheStep.Computing.Budgeted.check\_equal},
\lean{NativeTypesOnTheStep.Computing.Budgeted.check\_different},
\lean{NativeTypesOnTheStep.Computing.Budgeted.check\_undecided\_iff},
\lean{NativeTypesOnTheStep.Computing.Budgeted.check\_mono}). It is complete in the limit only on
terms that both reach a normal form within some budget
(\lean{NativeTypesOnTheStep.Computing.Budgeted.check\_eventually}); convertible terms without normal
forms stay undecided at every budget (Section~\ref{sec:addendum-budgets}).

\emph{What does not hold.} The step's own declared packages are not qualified: with addition and
\mt{append} declared computing over the rigid numbers, the package has no normalization setting
under its own roles (\lean{NativeTypesOnTheStep.Computing.Vectors.declaredRigid\_no\_setting}). For
it the syntactic check, Church--Rosser, rows typed after their match and unique answers of ground
calls are proved (\lean{NativeTypesOnTheStep.Computing.DeclarationCheck.append\_check},
\lean{NativeTypesOnTheStep.Computing.VectorFunction.append\_churchRosser},
\lean{NativeTypesOnTheStep.Computing.VectorRows.append\_rows\_typed},
\lean{NativeTypesOnTheStep.Computing.VectorFunction.append\_answer\_unique}); subject reduction on
open terms and decidable typed conversion are not. The operational reading is the package's
reduction relation, and no theorem yet ties firings of the full step to it.

\subsection{Indexed families and a dependent program checked by its match}\label{sec:addendum-indexed}
Vectors have a dependent eliminator in the normalization model, with every hypothesis of the model
proved. The family is computed from its index: $\mathrm{vec}\,A\,\mathrm{zero}$ computes to the unit
type and $\mathrm{vec}\,A\,(\mathrm{succ}\,n)$ to $\Sigma\,A\,(\mathrm{vec}\,A\,n)$. The eliminator
computes on open terms, and \mt{append} is typed at $\mathrm{vec}\,A\,(\mathrm{add}\,n\,m)$ in the
context of its variables
(\lean{Mettapedia.TypeTheory.Calculi.ParameterizedPiSigmaId.Presentation.TypedEquality.Normalization.TowerVectors.vecElim\_cons},
\lean{Mettapedia.TypeTheory.Calculi.ParameterizedPiSigmaId.Presentation.TypedEquality.Normalization.TowerVectorsPrograms.append\_typed}).
On the dependent step with this package's conversion, the recursive firing of \mt{append}, which
makes its result at the call's declared type $\mathrm{vec}\,A\,(\mathrm{add}\,(\mathrm{succ}\,k)\,m)$
and tells that the call equals it, is admitted with every variable open
(\lean{ProgrammingOnTheStep.IndexedFamilies.Eliminator.recursive\_admitted}). In the package where
\mt{append}, \mt{cons} and \mt{add} are rigid, the same firing is refused whatever type its made row
is given (\lean{ProgrammingOnTheStep.IndexedFamilies.Rigid.rigid\_refused}); there the two lengths
have no common instance, and here they are equal
(\lean{ProgrammingOnTheStep.IndexedFamilies.Rigid.lengths\_apart},
\lean{ProgrammingOnTheStep.IndexedFamilies.Eliminator.lengths\_meet}).

Fourteen rows, accepted by one Boolean pass, present the package's computation, and the package is
qualified at every level and valuation
(\lean{NativeTypesOnTheStep.Computing.Vectors.vectors\_qualified}); one more row that preserves
types, $\mathrm{add}\,\mathrm{zero}\,b\to\mathrm{zero}$, makes $1$ equal to $0$
(\lean{NativeTypesOnTheStep.Computing.Vectors.overlap\_identifies}). On the converting step, which
settles a firing along its settler's answer and checks each made row there by conversion, the
package's computable settler decides admission with three outcomes, the third claiming nothing
(\lean{NativeTypesOnTheStep.Computing.Vectors.vecAdmits?\_sound},
\lean{NativeTypesOnTheStep.Computing.Vectors.vecAdmits?\_refuted}). By kernel evaluation it admits
four rows typed only after their match: the nil and cons rows of \mt{append}, the cons row of
\mt{head}, and an interpreter's literal row, whose type $\mathrm{val}\,0$ computes to the numbers
(\lean{NativeTypesOnTheStep.Computing.Vectors.nil\_decided},
\lean{NativeTypesOnTheStep.Computing.Vectors.cons\_decided},
\lean{NativeTypesOnTheStep.Computing.Vectors.head\_decided},
\lean{NativeTypesOnTheStep.Computing.Vectors.val\_decided}). The dependent step, which settles along
the identity, refuses the cons row and the literal row, as the first limit of
Section~\ref{sec:addendum-nov12} says
(\lean{NativeTypesOnTheStep.Computing.Vectors.cons\_present\_refused},
\lean{NativeTypesOnTheStep.Computing.Vectors.val\_present\_refused}). Mistyped rows are refused by
every settler: \mt{(= (append \$xs \$ys) \$xs)}, the nil row answering its first argument, and the
literal row told $\mathrm{code}\doteq 1$
(\lean{NativeTypesOnTheStep.Computing.Vectors.keep\_refused},
\lean{NativeTypesOnTheStep.Computing.Vectors.nilFirst\_refused},
\lean{NativeTypesOnTheStep.Computing.Vectors.valOne\_refused}).

\emph{What does not hold.} Row typings are an input of this decision: it takes each made row with a
proof of its type in its own frame, and the kernel decides the settler's answer and the conversion at
it, not that typing; the definition checker of Section~\ref{sec:addendum-definitions} elaborates the
rows and synthesizes their typings, on its class of rows. The family is recursive, not
inductive: its type former is not injective at length zero, and at a successor length its eliminator
unfolds on a variable vector
(\lean{Mettapedia.TypeTheory.Calculi.ParameterizedPiSigmaId.Presentation.TypedEquality.Normalization.TowerVectors.vec\_zero\_not\_injective},
\lean{Mettapedia.TypeTheory.Calculi.ParameterizedPiSigmaId.Presentation.TypedEquality.Normalization.TowerVectors.vecElim\_succ\_unfolds}).
So $\mathrm{vec}\,X\,\mathrm{zero}\doteq\mathrm{vec}\,\mathrm{num}\,\mathrm{zero}$ is solved by
$X\mapsto\mathrm{num}$ and by $X\mapsto\mathrm{unit}$, and the element type of an empty vector is not
inferred (\lean{NativeTypesOnTheStep.Computing.Vectors.emptyElem\_num},
\lean{NativeTypesOnTheStep.Computing.Vectors.emptyElem\_unit}). The motive and the element type live
in the smallest universe only, and a term type with two constructors at one index, such as a typed
interpreter's, needs inductive families proper, with a new indexed case of the model's logical
relation; an extrinsic encoding, and the family pack that such a case would use, are in
Section~\ref{sec:addendum-families}, and the case itself is not built. Unification modulo the package's computation is not unitary: $\mathrm{add}\,x\,y\doteq 1$
has two solutions and no principal answer, and the settler postpones it without refuting it, so the
law of Section~\ref{sec:addendum-nov12} that every solvable store has a principal answer fails here
for a new reason
(\lean{ProgrammingOnTheStep.IndexedFamilies.AcceptanceOnTheRepairedStep.principalLaw\_fails},
\lean{ProgrammingOnTheStep.IndexedFamilies.AcceptanceOnTheRepairedStep.principal\_test\_postponed}).
That store is outside the pattern fragment; whether the law also fails inside the fragment over this
package is open.

\subsection{Extensional equality: the options and what queries see}\label{sec:addendum-equality}
Queries see exactly what the knowledge identifies. An extension of the equality on rows leaves every
count answer unchanged exactly when it relates nothing new
(\lean{EqualityOnTheStep.Options.Observability.count\_conservative\_iff}), and on the dependent step a
counting firing refused under one knowledge is admitted under a solvable knowledge that entails one
more equation (\lean{EqualityOnTheStep.Options.Observability.knowledge\_observed}). Evidence rows are
invisible to queries: a transport firing leaves every count and absence answer as it was, which is
why absence does not transport
(\lean{EqualityOnTheStep.Options.Observability.transport\_answers}). An evidence-free class equation
is visible, and so would be definitional proof irrelevance
(\lean{EqualityOnTheStep.Options.Observability.unwitnessed\_tell\_observed},
\lean{EqualityOnTheStep.Options.Observability.proofIrrelevance\_observed}).

Respect terms are generated from a family's type structure, as observational type theory computes
coercions, but outside the knowledge, which is unchanged. For the class built from membership,
constants, functions, pairs, and $\Pi$ and $\Sigma$ over domains that do not mention the set, the
generated transports are typed in both directions at typed parameters, sets and evidence, and the
transport firing is admitted for every family of the class
(\lean{EqualityOnTheStep.Options.Respect.tr\_typed},
\lean{EqualityOnTheStep.Options.Respect.transport\_by\_code}); the intensional identity family is
outside it, even up to the knowledge's equality
(\lean{EqualityOnTheStep.Options.Respect.idFamily\_outside}). The transports are coherent in the
knowledge along reflexivity and along composite evidence, and the row a transport firing makes
composes too, over a formed context (\lean{EqualityOnTheStep.Options.Respect.tr\_refl},
\lean{EqualityOnTheStep.Options.Respect.tr\_comp},
\lean{EqualityOnTheStep.Options.Respect.made\_row\_comp}); typing alone does not give coherence
(\lean{EqualityOnTheStep.Options.Respect.incoherent\_respect}).

Two options are worked only in models, with no syntax and no step built for them. Propositions in a
proof-irrelevant universe, in the tree's predicate-family model against its set-family model: the
two sets become equal and their members carry no information, so a proof-relevant encoding survives
only through its truncation (\lean{EqualityOnTheStep.Options.Consequences.rowB}), and evidence rows
combine with this universe except for a family that counts members
(\lean{EqualityOnTheStep.Options.Observational.counting\_not\_respected}). A univalent option, in the
groupoid of small types, where a path is defined to be an equivalence, so that univalence holds
there by definition and not by a theorem: transport along the swap computes, equivalent types have
equally many elements, and an identity of types with unique proofs is not univalent
(\lean{EqualityOnTheStep.Options.Consequences.rowC}).

A read can identify class names without evidence: reading the class row $(S,q_1)$ with the pattern
$(S,q_2)$ is admitted, tells nothing, and identifies $q_1$ and $q_2$
(\lean{EqualityOnTheStep.Options.QuotientDiscipline.tells\_alone\_insufficient}), so a discipline over
tells alone is not enough. The discipline that asks for an evidence row, read by the firing, for each
pair of class names the firing newly identifies refuses both evidence-free routes and admits the
quotient's soundness firing
(\lean{EqualityOnTheStep.Options.QuotientDiscipline.discipline\_controls}). It is proved one firing
at a time on the step; over whole runs it is proved only in a registry form with a sole writer
(\lean{EqualityOnTheStep.Options.QuotientDiscipline.identified\_witnessed}).

\emph{Open:} the discipline over runs of the dependent step; which families have respect terms,
since the class is sufficient and not exhaustive; observational or cubical syntax in the knowledge.

\subsection{Types as classes no observer can split}\label{sec:addendum-types-classes}
An observer family reads instances; two instances are apart when some reading tells them apart, and
indistinguishable when every reading agrees. Apartness is an exclusion relation of the tree's
distinction calculus, and its concepts
(\lean{Mettapedia.Cybernetics.DistinctionCalculus.Exclusion.IsConcept}) are exactly the classes
closed under indistinguishability
(\lean{LogicOverTheStep.Distinctions.ObserverConcepts.isConcept\_iff\_saturated}): a type is a class
that no admitted observer can split. One family has every type of another exactly when its
indistinguishability is finer, and an operation belongs to a lens exactly when it sends
indistinguishable inputs to indistinguishable outputs
(\lean{LogicOverTheStep.Distinctions.ObserverConcepts.concepts\_subset\_iff},
\lean{LogicOverTheStep.Distinctions.ObserverConcepts.pullback\_isConcept\_iff}). For tests closed
under negation and conjunction, a class is the extension of one test exactly when finitely many
tests decide it
(\lean{LogicOverTheStep.Distinctions.ObserverConcepts.BooleanTests.denoted\_iff\_finitelyDetermined}).

On the step the observers are modal formulas over its firings. On the first-order step under every
finite program, with the hypotheses of \lean{LogicOverTheStep.ProgramAdequacy.program\_adequacy} (a
supply of fresh identities, finitely many rules, observations invariant under renaming), the types
of the finitary formulas are exactly the classes closed under bisimilarity
(\lean{LogicOverTheStep.Distinctions.StepConcepts.program\_isConcept\_iff\_bisimClosed}); for the
infinitary formulas this holds on every system, each type denoted by one formula
(\lean{LogicOverTheStep.Distinctions.StepConcepts.isConcept\_iformulas\_iff},
\lean{LogicOverTheStep.Distinctions.StepConcepts.exists\_iformula\_of\_bisimClosed}). The finitary
formulas do not denote every type: on the ground first-order step with the one rule that takes one
\mt{T}, the stores with an even number of \mt{T} form a type that no finitary formula denotes, and
with only the knowledge observed every single reach is denoted and the even reaches are not
(\lean{LogicOverTheStep.Distinctions.ParityBoundary.evenCount\_isConcept},
\lean{LogicOverTheStep.Distinctions.ParityBoundary.evenCount\_not\_denoted},
\lean{LogicOverTheStep.Distinctions.ParityBoundary.evenReach\_not\_denoted}). OSLF's possibility and
successor necessity over the firings of any \lean{Interface.Step} carry the types closed under
bisimilarity (\lean{LogicOverTheStep.Distinctions.Survival.step\_possibility\_isConcept},
\lean{LogicOverTheStep.Distinctions.Survival.step\_necessity\_isConcept}); predecessor necessity does
not (\lean{LogicOverTheStep.Distinctions.Survival.Past.predecessor\_control}).

The number lens is a model of measured behaviour, not of an implementation: a stored number is an
integer or a fraction, and division truncates when both arguments are stored as integers. Under the
value lens the stored integer $5$ and the fraction $5$ lie in every type together
(\lean{LogicOverTheStep.Distinctions.NumberLens.five\_same\_concepts}); the storage lens, which also
reads quotients by that division, identifies two numbers exactly when they have the same value and
are stored alike, or the value is zero
(\lean{LogicOverTheStep.Distinctions.NumberLens.indist\_storage\_iff}); exact division belongs to the
value lens and the storage-dependent division does not
(\lean{LogicOverTheStep.Distinctions.NumberLens.exactQuotient\_respects},
\lean{LogicOverTheStep.Distinctions.NumberLens.storedQuotient\_not\_respects}). Typing is stated at
solutions, since the knowledge has no observer of raw rows: typing a made row at every solution is a
type, and the raw typing of the $\eta$ row is not
(\lean{LogicOverTheStep.Distinctions.TypingConcepts.typedAfterMatch\_isConcept},
\lean{LogicOverTheStep.Distinctions.TypingConcepts.etaRow\_not\_isConcept}).

Without a finite program the answer depends on the observation set. On the ground first-order step
with every rule instance admitted, one firing reaches every consistent extension of the knowledge.
With the knowledge observed, or the knowledge and the bag of labels, the finitary types are again
the classes closed under bisimilarity, because those observations already separate what bisimilarity
separates (\lean{LogicOverTheStep.Prices.UnrestrictedStep.knowledgeView\_adequacy},
\lean{LogicOverTheStep.Prices.UnrestrictedStep.storeView\_adequacy\_base}). For other observation
sets invariant under renaming, one of them with four observations, two stores satisfy the same
finitary formulas and are not bisimilar
(\lean{LogicOverTheStep.Prices.UnrestrictedStep.unrestricted\_adequacy\_fails},
\lean{LogicOverTheStep.Prices.StaircaseStep.staircase\_adequacy\_fails}); under every finite program
those sets are adequate, by the theorem above. So the finite program is the hypothesis that cannot
be dropped from the statement over all invariant observation sets.

\emph{Open:} the exact boundary for graded formulas. The number lens sees storage
only through division by nonzero divisors; printing, which also tells $5$ from $5.0$, is not
modelled.

\subsection{When types apply}\label{sec:addendum-when-types}
Fixed-point formulas, the tree's modal $\mu$-calculus read over the step's systems with atoms as free
variables, denote types only when their fixed points have positive bodies. Then states related by a
bisimulation that respects the atoms satisfy the same formulas
(\lean{LogicOverTheStep.WhenTypesApply.FixpointTypes.satisfies\_iff\_of\_bisimulation}); without
positivity a least fixed point splits two bisimilar states of a four-state system
(\lean{LogicOverTheStep.WhenTypesApply.FixpointTypes.Controls.nonpositive\_splits}). On the one-letter
chain these formulas denote exactly the ultimately periodic sets
(\lean{LogicOverTheStep.WhenTypesApply.OneLetterChain.denoted\_iff\_periodic}). On the step, under the
program that takes one \mt{T} with the knowledge observed, the ladder of finitary, fixed-point and
infinitary types is strict (\lean{LogicOverTheStep.WhenTypesApply.Ladder.ladder}): the even reaches
are denoted by $\mu X.\,[\cdot]\bot\vee\langle\cdot\rangle\langle\cdot\rangle X$ and by no finitary
formula, and the stores whose reach is a perfect square by an infinitary formula and by no
fixed-point formula. There a class given by reach is denoted by a fixed-point formula exactly when its
set of reaches is ultimately periodic
(\lean{LogicOverTheStep.WhenTypesApply.Ladder.reachClass\_fixDenoted\_iff}). With the bag of labels
observed as well, no fixed-point formula denotes the even counts
(\lean{LogicOverTheStep.WhenTypesApply.Ladder.evenCount\_not\_fixDenoted}): a store with inconsistent
knowledge cannot fire, and a fixed point counts only what the step unfolds.

Least and greatest fixed points are checked under a budget with three outcomes, member, nonmember
and undecided: member and nonmember are sound, and a decided verdict is kept at every larger budget
(\lean{LogicOverTheStep.WhenTypesApply.Groundable.check\_member},
\lean{LogicOverTheStep.WhenTypesApply.Groundable.check\_nonmember},
\lean{LogicOverTheStep.WhenTypesApply.Groundable.check\_mono}). On image-finite systems, for bodies
with no negation and no nested fixed point, every member of a least fixed point is found and every
non-member of a greatest one is refuted at some budget
(\lean{LogicOverTheStep.WhenTypesApply.Groundable.check\_complete\_lfp},
\lean{LogicOverTheStep.WhenTypesApply.Groundable.check\_complete\_gfp}); through the bisimulation
quotient so on the first-order step under every finite program, with the hypotheses of program
adequacy (\lean{LogicOverTheStep.WhenTypesApply.GroundableUpTo.program\_mu\_groundable},
\lean{LogicOverTheStep.WhenTypesApply.GroundableUpTo.program\_nu\_refutable}). Without
image-finiteness the check stays undecided at every budget on a member
(\lean{LogicOverTheStep.WhenTypesApply.Groundable.stage\_omega\_not\_enough}).

For a lens whose readings a join makes true exactly when one side does, such as which labels are
present and which equations are known, a type survives every growth by another party exactly when it
is built from the readings by unions and intersections, with no complement
(\lean{LogicOverTheStep.WhenTypesApply.Changes.stable\_iff\_positive}). Absence, an exact count
and ``these are all its equations'' are types that one growth breaks
(\lean{LogicOverTheStep.WhenTypesApply.Changes.absentU\_not\_stable},
\lean{LogicOverTheStep.WhenTypesApply.Changes.exactlyOneT\_not\_stable},
\lean{LogicOverTheStep.WhenTypesApply.Changes.onlyEquations\_not\_stable}), and both hypotheses are
needed (\lean{LogicOverTheStep.WhenTypesApply.Changes.identity\_stable\_not\_positive},
\lean{LogicOverTheStep.WhenTypesApply.Changes.storeView\_not\_joinPreserving}). A closure clause,
``these and no others'', is an absence statement about a region of rows, and every such statement
survives every step of another party exactly when no other party creates a row of the region
(\lean{LogicOverTheStep.WhenTypesApply.ThatsAll.closedReadingSurvives\_iff\_owned}). The equivalence
is elementary; its content is in the instances. A publication system with a sole writer owns its
protected region, so every closure statement about protected rows survives every step of another
actor (\lean{LogicOverTheStep.WhenTypesApply.ThatsAll.publication\_closedReadingSurvives}); without
ownership, a later row overturns a refutation drawn from ``these equations are all of it''
(\lean{LogicOverTheStep.WhenTypesApply.ThatsAll.later\_row\_overturns}).

\emph{Open:} groundability for bodies with negation or with nested or alternating fixed points; the
growth theorem on the step for counts ``at least $n$'' and for entailment by the knowledge; a theorem
that exact counts are decided at a store.

\subsection{Budgets as outcomes}\label{sec:addendum-budgets}
The unifier's rounds are driven by its rank (Section~\ref{sec:addendum-levels}); three budgets
remain: the computation inside one round, the depth of phases, and the reader's depth. Each is now a
budgeted decision procedure (\lean{UnificationWithLevels.TreeFibre.Budgets.Procedure}) with three
outcomes at every budget $k$, yes, no and ``checked up to $k$'', and three laws as fields: yes and no
are sound, raising the budget only resolves, and exhausted holds exactly when the computation ran
out. Exhausted is never read as no or as empty: reading it as no, or as yes, breaks the information
order on every instance later answered the other way
(\lean{UnificationWithLevels.TreeFibre.Budgets.Procedure.no\_reading\_breaks},
\lean{UnificationWithLevels.TreeFibre.Budgets.Procedure.yes\_reading\_breaks}). A budgeted decider
becomes such a procedure by taking, at budget $k$, the first answer found up to $k$; ``answered at
some budget'' then holds exactly on the instances that some budget's attempt answers
(\lean{UnificationWithLevels.TreeFibre.Budgets.Decider.answered\_iff}), which defines the class and
is not a completeness theorem for the solvers. Off that class no budget answers: $\Omega$ is
convertible with itself and exhausted at every budget, for every normalizer
(\lean{UnificationWithLevels.TreeFibre.Budgets.omega\_undecided}), and so is a solvable store
without a principal answer, for every decision of stores
(\lean{UnificationWithLevels.TreeFibre.Budgets.undecided\_of\_no\_principal}).

On the step, a firing that only tells is admitted exactly when its information has a solution, for
every settler (\lean{UnificationWithLevels.TreeFibre.Budgets.OnTheStep.onlyTells\_admits\_iff}):
solved means admitted and refuted means refused, and exhausted means neither, since one firing
exhausted at every budget is admitted and another refused
(\lean{UnificationWithLevels.TreeFibre.Budgets.OnTheStep.cross\_exhausted\_admitted},
\lean{UnificationWithLevels.TreeFibre.Budgets.OnTheStep.het\_exhausted\_refused}). The computable
converting step reads exhaustion as refusal: with its settler's families at budget $1$ it refuses a
firing that the semantic step admits, and at budget $2$ it admits it
(\lean{UnificationWithLevels.TreeFibre.Budgets.OnTheStep.lift\_cliff}). The semantic step is
unchanged; the verdict is reported beside it, makes no row, and leaves the state reached, and every
later answer, the same at every budget
(\lean{UnificationWithLevels.TreeFibre.Budgets.OnTheStep.later\_budget\_free}).

\emph{What does not hold.} No budget of the unifier has been removed. For the depth of phases, a
canonical order of declarations does not remove it: the store $F\doteq\lambda x.\,x$ over one
declaration, where no exchange is possible, is exhausted at depth $0$ and solved at depth $1$,
through a first phase that puts it in $\eta$-long form
(\lean{UnificationWithLevels.TreeFibre.Budgets.TowerControls.id\_exhausted\_then\_solved}). Inside a
round, the fuelled comparison succeeds eventually wherever its reducer does, and the round's
conversion scan stops depending on the budget from some budget on, for the vectors package under
strong normalization of its typed terms, which is not proved (Section~\ref{sec:addendum-definitions});
the round's part that runs the executable checker still needs the checker to succeed eventually, and
removing the reader's depth needs a rank under which two exchanges give back the same store; both are
open. One control is partial: an unsolvable store is shown exhausted inside a round only at the
budgets $0$ to $3$.

\subsection{Evaluation modes}\label{sec:addendum-evaluation-modes}
This subsection is about a model of a design space, not about any evaluator. One small core over the
tree's MeTTa atoms, with quotation, running and binding, and one elaboration function parametrized by
six choices give 96 designs (\lean{DialectsAsStrategies.EvaluationModes.Design}): where a held
parameter runs (never, only as the whole body, or at every computation position); whether a
\mt{Proc} result holds it while other computation types run it; whether values run again; whether a
body under an \mt{Atom} result is evaluated; whether heads with no definition evaluate their
arguments; and whether a refused check prints an error. HE, by-name, code value, a hybrid and a draft
of Prime's evaluator are values of this record, and PeTTa is modelled by code value. The models of HE,
PeTTa and the draft are accepted because the kernel evaluates each to its system's answers on eight
programs, measured on hyperon-experimental 0.2.10 with three later fixes, on CeTTa 1.5.0-dev in its
PeTTa mode, and on the draft (\lean{DialectsAsStrategies.EvaluationModes.Acceptance.he\_column},
\lean{DialectsAsStrategies.EvaluationModes.Acceptance.codeValue\_column},
\lean{DialectsAsStrategies.EvaluationModes.Acceptance.primeToday\_column}); nothing more is claimed
about those systems, and HE's column is claimed only on rows without \mt{Proc}, a type HE lacks.

Weakening a result type along the line from \mt{Number} through \mt{Proc} to \mt{\%Undefined\%}
changes no run that refused nothing, in every design where held parameters run alike under the two
types (\lean{DialectsAsStrategies.EvaluationModes.Weakening.weakening}); the hybrid, which holds them
under \mt{Proc}, breaks it, and no design that keeps it passes all three test lines of the draft
(\lean{DialectsAsStrategies.EvaluationModes.Weakening.hybrid\_breaks\_weakening},
\lean{DialectsAsStrategies.EvaluationModes.Separation.primeTests\_break\_weakening}). Weakening,
\mt{let}, final values and identity insertion hold together exactly for the by-name and code-value
rules for held parameters, whatever the other three choices
(\lean{DialectsAsStrategies.EvaluationModes.Separation.four\_laws\_iff}). The choice between them is
one bit: under the four laws, the plain reading of a returned held parameter holds exactly for
by-name, inert held text exactly for code value, and no design keeps both
(\lean{DialectsAsStrategies.EvaluationModes.Separation.plain\_fork},
\lean{DialectsAsStrategies.EvaluationModes.Separation.inert\_fork},
\lean{DialectsAsStrategies.EvaluationModes.Separation.not\_plain\_and\_inert}). The draft's
\mt{\%Undefined\%} test line asks for the first and its two \mt{Proc} lines for the second. With
typed code, an occurrence of \mt{(Code T)}, for \mt{T} a number, Boolean, \mt{Proc} or variable type,
needs a default exactly where \mt{\%Undefined\%} or a type variable not occurring in \mt{T} is
expected (\lean{DialectsAsStrategies.EvaluationModes.Typed.needsDefault\_iff}).

\emph{What does not hold.} ``All 96 designs'' means all values of this record; a design outside it is
not covered. In the model, heads that are not symbols are data, and there is no nondeterminism, no
effect and no sharing by need. The results say which laws each design keeps; they choose no design.
HE's substitution of printed values is shown only on the measured rows, and the plain reading for
by-name only at the body root and at a constructor field; both need a general substitution lemma. A
tie from this core's runs to the step is open.

\subsection{Certificates for behavioural distance}\label{sec:addendum-certificates}
A certificate for a finite chain with rational masses is a rational kernel on pairs of states, with
couplings for some pairs. A Boolean checker decides it in exact rational arithmetic, filling in the
omitted pairs and computing the cost of every coupling it uses; if it accepts, the bisimulation
metric is at most the kernel, at every discount $0\le c$
(\lean{LogicOverTheStep.Distinctions.Certificates.check\_sound}). A test formula, given as a graph
with shared nodes, bounds the metric from below at $0\le c\le 1$
(\lean{LogicOverTheStep.Distinctions.Certificates.testGap\_le\_bisimulationMetric}); when the two
bounds meet the value is exact
(\lean{LogicOverTheStep.Distinctions.Certificates.bisimulationMetric\_eq\_of\_meet}), and otherwise
``distance at most $\varepsilon$'' is not decided exactly when the lower bound is at most
$\varepsilon$ and the upper bound exceeds it
(\lean{LogicOverTheStep.Distinctions.Certificates.within\_undecided\_iff}), which is never read as a
refutation. A dual certificate, a post-fixpoint, bounds the metric from below at $0\le c<1$
(\lean{LogicOverTheStep.Distinctions.Certificates.dualCheck\_sound}); at $c=1$ this fails, with a
counterexample that both checkers accept
(\lean{LogicOverTheStep.Distinctions.Certificates.Controls.postfixpoint\_above\_metric\_at\_one}).

Between a state that stays and one that leaks into a sink with probability $1/2$, the distance at
discount $9/10$ is exactly $9/11$
(\lean{LogicOverTheStep.Distinctions.Certificates.Controls.metric\_safe\_leaky}); a correct
certificate is accepted, and certificates with a wrong marginal, a kernel too small at one pair, or
a missing coupling are refused
(\lean{LogicOverTheStep.Distinctions.Certificates.Controls.correct\_accepted},
\lean{LogicOverTheStep.Distinctions.Certificates.Controls.wrongMarginal\_refused},
\lean{LogicOverTheStep.Distinctions.Certificates.Controls.tooSmall\_refused},
\lean{LogicOverTheStep.Distinctions.Certificates.Controls.missing\_refused}); no test of any depth
closes that pair, which the dual certificate decides
(\lean{LogicOverTheStep.Distinctions.Certificates.Controls.measureByTest\_never\_exact}). On the step,
a finite region of states closed under a stochastic policy is the step's own probabilistic system
there: its metric vanishes exactly on the step's Larsen--Skou bisimilarity, at $0<c\le 1$, and the
same checker decides certificates about it
(\lean{LogicOverTheStep.Distinctions.Certificates.behaviouralMetric\_region\_eq\_zero\_iff},
\lean{LogicOverTheStep.Distinctions.Certificates.check\_sound\_region}). On a small machine compiled to
an instance of the step, with a policy whose positive mass lies on firings, a state that fires and
one that refuses, with the same reading, are at distance exactly $9/10$ at discount $9/10$
(\lean{LogicOverTheStep.Distinctions.Certificates.StepControls.fires\_versus\_refuses}).

\emph{What does not hold.} The certificate is a Lean structure: no external solver's output is read,
and mapping concrete certificate records onto it is open. Kernel evaluation of this checker is shown
on chains of two to five states, and of the checker on classes on a region of the step of up to
eleven states (Section~\ref{sec:addendum-compact}); larger certificates need a compiled checker.
Undiscounted distances have no exact rule, and completeness of the checker is open.

\subsection{\texorpdfstring{\texttt{let}}{let} on every run, and the measured ledger}\label{sec:addendum-ledger}
On the Lean model of HE, take \mt{(let \$x d b)} with $d$ a data atom of one closed stored type and
$b$ an expression whose only variable is \mt{\$x}, called at HE's entry with empty bindings and a
plain environment. For any policy and any instance $b[x:=d]$ that meets two contracts, which the
settling class with \mt{if} nodes meets, a complete run exists, every complete run publishes the
answers of $b[x:=d]$, and every run finishes
(\lean{DialectsAsStrategies.TriangleFootprints.LetEveryRun.let\_every\_run\_he},
\lean{DialectsAsStrategies.TriangleFootprints.LetEveryRun.let\_settles\_every\_run}). The model of
PeTTa yields the same answers from some fuel on; with no hypothesis left on PeTTa's side this is
proved for \mt{(let \$x b (if \$x y n))}, with $b$ a Boolean and data branches
(\lean{DialectsAsStrategies.TriangleFootprints.LetPeTTaTemplates.let\_if\_agree}), and for other
templates PeTTa's side is assumed in two parts: that its binding of the renamed variable to the value
succeeds, and that the renamed template then runs under that binding to the same answers
(\lean{DialectsAsStrategies.TriangleFootprints.LetPeTTaTemplates.let\_agree\_under}).

The model of HE fixes the binding record of a matching row only up to its meaning, and so admits a
circular record at which a later type cast fails, where HE's interpreter merges syntactically and
succeeds. The variant matching step takes the syntactic merge itself: every frontier of the model is
a variant one up to how records are written, and every result of the model about all runs holds for
the variant machine unchanged
(\lean{DialectsAsStrategies.TriangleFootprints.SyntacticMerge.frontier\_represented},
\lean{DialectsAsStrategies.TriangleFootprints.SyntacticMerge.every\_run\_of\_model}). On it,
\mt{(pair (if True 1 2) (if True True False))} publishes \mt{(pair 1 True)} on every run, as HE's
interpreter does, and the variant machine and PeTTa agree on every run on trees with declared \mt{if}
(\lean{DialectsAsStrategies.TriangleFootprints.DeclaredBranchVariant.Controls.pair\_if\_every\_run},
\lean{DialectsAsStrategies.TriangleFootprints.DeclaredTrees.declared\_trees\_every\_run}). The
candidate of the model's primitive \mt{unify} admits the same circular record, which the variant does
not restrict (\lean{DialectsAsStrategies.TriangleFootprints.UnifyCandidateGap.unify\_cycle\_admitted}).
Both circular records are properties of the Lean model of HE, not of HE. \emph{Open:} \mt{let} at
frames with bindings and marked values, which \mt{let} inside bodies and branches needs; the
syntactic merge at \mt{unify}.

\emph{The measured ledger.} The statements of Section~\ref{sec:addendum-nov12} on HE and PeTTa are
theorems about the Lean models of the two dialects. Measured on hyperon-experimental 0.2.10 with three
later fixes and on CeTTa 1.5.0-dev in its HE and PeTTa modes, the differences there for an unmatched
call, a non-Boolean condition and a free body variable hold of the implementations as stated, and the
agreements hold as bags: HE publishes a nested argument's alternatives in another order than the
model, and the general theorems about all runs state HE's answers as bags
(\lean{DialectsAsStrategies.TriangleFootprints.NestedEveryRun.nested\_every\_run}). That HE checks
declared argument types and PeTTa does not is true of the PeTTa model only: CeTTa's PeTTa mode checks
a declared head and gives no answer when the check fails. The declared-type theorems therefore state
PeTTa's side on the program without the head's declaration, where the model gives the measured
answer, and the model answers the same with and without the declaration
(\lean{DialectsAsStrategies.TriangleFootprints.DeclaredEveryRun.declared\_every\_run},
\lean{DialectsAsStrategies.TriangleFootprints.TypedEveryRun.typed\_every\_run},
\lean{DialectsAsStrategies.TriangleFootprints.DeclaredHeads.headUndeclared},
\lean{DialectsAsStrategies.TriangleFootprints.DeclaredErrors.Controls.petta\_model\_reads\_no\_declaration}).

The agreement on integer arithmetic holds for every grounded table that meets a per-call contract
and answers the class's calls with PeTTa's values (Section~\ref{sec:addendum-tables}): for the model's
unbounded integers, and for the table of HE's grounded operations as measured where results fit in 64
bits. That table, in which division and remainder truncate and division by zero and overflow
are error answers (\lean{DialectsAsStrategies.TriangleFootprints.GroundedTables.measuredOutcome}),
agrees with the model's table on addition, subtraction, multiplication and comparison whose integer
results fit in 64 bits, and the smallest call is named for each place where the two part
(\lean{DialectsAsStrategies.TriangleFootprints.GroundedTables.measured\_agrees},
\lean{DialectsAsStrategies.TriangleFootprints.GroundedTables.tables\_part}). On HE's machine with that
table a complete run exists, and every complete run of \mt{(\% 7 0)} and of
\mt{(+ 9223372036854775807 1)} publishes an error; the PeTTa model faults on the first and answers
$2^{63}$ on the second, and so does HE's machine with the model's table on the second
(\lean{DialectsAsStrategies.TriangleFootprints.GroundedDifferences.Controls.zero\_divisor\_every\_run},
\lean{DialectsAsStrategies.TriangleFootprints.GroundedDifferences.Controls.overflow\_every\_run},
\lean{DialectsAsStrategies.TriangleFootprints.GroundedDifferences.Controls.model\_overflow\_every\_run}).
The measured table is HE's only on calls of the right arity whose arguments have no declared type or
the matching one: HE checks a call of a grounded operation against the operation's type before running
it, and the model gives grounded operations no types; measured, \mt{(+ 1 True)} is an argument-type
error and \mt{(+ 1)} an arity error on HE. That region is stated as a decidable predicate, with
controls on both sides (\lean{DialectsAsStrategies.TriangleFootprints.GroundedTables.TypedRegion},
Section~\ref{sec:addendum-tables}).
A further difference is named: on \mt{(= (g) 1)} and \mt{(= (g) (Error foo bar))} every complete HE
run publishes \mt{1} alone, while the PeTTa model yields \mt{1} and the error, so HE's answers are
not the union over its alternatives
(\lean{DialectsAsStrategies.TriangleFootprints.ErrorAlternatives.Controls.error\_dropped\_every\_run},
\lean{DialectsAsStrategies.TriangleFootprints.ErrorAlternatives.Controls.not\_union}). Two further
causes make the model's type query differ from the one measured on HE, each with its program: a
declaration with a variable subject, and a subtype row
(\lean{DialectsAsStrategies.TriangleFootprints.TypeQueryFragment.Controls.variable\_subject\_needed},
\lean{DialectsAsStrategies.TriangleFootprints.DeclaredErrors.Controls.subtype\_row\_needed}). The
declared-type theorems state that their programs have neither
(\lean{DialectsAsStrategies.TriangleFootprints.TypeQueryFragment.MeasuredTypeFragment}). That the
two queries agree on such programs, the types of grounded operations aside, is measured on those
theorems' programs and is not proved.

The Lean declarations cited in this addendum use \lean{propext}, \lean{Classical.choice} and
\lean{Quot.sound}, or fewer; none uses \lean{sorryAx}.

\section{Addendum: results of 13 November}\label{sec:addendum-nov13}
These results extend the earlier addenda and narrow open rows of Section~\ref{sec:open}, on
placements of rho, on an equivalence of programs and on programs at dependent frames, and the budgets
of Section~\ref{sec:addendum-budgets}; the earlier sentences they make stale are corrected where they
stand. The retained, revised and rejected rows are unchanged. Sections~\ref{sec:addendum-hosting}, \ref{sec:addendum-bags} and \ref{sec:addendum-congruence} are
stated on the step of a program (Definition~\ref{def:program-step}). No other result below is yet;
each is about one of the existing steps that the definition restricts. Where a result runs a set of rule instances that are the rules of a ground
program without census patterns, the step of that set is the first-order step restricted to the
program's fresh openings (\lean{RuleSyntax.OneFrame.instanceSet\_eq}), and so the step of the program
in the restricted sense stated after Definition~\ref{def:program-step}; when the rules have unknowns,
the set shares them with the caller and the two differ
(\lean{RuleSyntax.OneFrame.Controls.instanceSet\_captures}).

\subsection{A hosted source language, and runs corresponding both ways for its first fragment}\label{sec:addendum-hosting}
The source language has data, which are terms with binders at the declaration frames, and
computations (\lean{SourceLanguage.Comp}): return a datum; call a datum; run a datum as code; bind the
answers of one computation in another, the sequencing construct, written \mt{let}; choose between two
computations; fail; query, add to and remove from a space; make a new space; check a type; delay and
force a shared computation. A program is a list of equations, each with a left-hand datum and a
right-hand computation, together with a reading of data as code (\lean{SourceLanguage.Program}). A
configuration is a world that all alternatives share (two counters and the spaces) and a bag of
alternatives, each with its own store of equations, its focus and its pending frames
(\lean{SourceLanguage.Config}). A step moves one alternative (\lean{SourceLanguage.Step}), and the
answers of a configuration are a bag with one entry, a set of closed terms, for each finished
alternative whose store is solvable (\lean{SourceLanguage.Config.answers}).

A program is translated into a program of the step at the contextual declaration frames
(\lean{Hosting.translate}). Each equation gives one rule, which takes the row of an alternative whose
focus is a call of the equation's left side and makes the row of that alternative at the right side
(\lean{Hosting.schema}). Six further rules are the same for every program: a returned datum becomes a
value; a call becomes a value when no equation matches it; a choice becomes its left or its right
computation; a \mt{let} pushes a pending frame; and a value met by a pending frame resumes the frame's
body, telling the frame's pattern equal to the value. The translation of a union of programs is the
union of the translations (\lean{Hosting.translate\_union}). The start state of a query has no
knowledge, one row for the query's alternative, and one row for the left side of each equation
(\lean{Hosting.start}); the rule for a call that no equation matches asks, by an absence check, that
no such row matches the call (\lean{Hosting.unmatchedRule}). So a hosted program is present twice, as
rules and as rows, and the one check that is not positive is the one that makes an unmatched call a
value.

A state of the step \emph{represents} an alternative when its present rows are the alternative's row
and the equation rows, and the alternative's row, read at every solution of the knowledge into every
frame, is what the alternative reads under the solutions of its own store
(\lean{Hosting.Represents}). For programs and alternatives built from return, call, choice and
failure (\lean{Hosting.PlainComp}), every move of a represented alternative is matched, for each
alternative with a solvable store that the move produces, by one firing of the translated program
into a state that represents it (\lean{Hosting.simulation}), and every firing of the translated
program from a state that represents an alternative is a move of that alternative into one that the
resulting state represents (\lean{Hosting.reflection}). One state of the step holds one alternative.
A configuration is therefore matched by a \emph{frontier}, one represented state for each alternative
with a solvable store (\lean{Hosting.Frontier}), and the answers of the configuration are, as a bag,
the answers read off its frontier (\lean{Hosting.frontier\_answers}). Along runs, for a program and a
query of this fragment whose names are unopened: every configuration that a run of the source reaches
from the query has a frontier whose states are reached from the start state by runs of the step
(\lean{Hosting.runs\_forward}), and every state that a run of the step reaches from the start state
represents an alternative of a configuration that a run of the source reaches
(\lean{Hosting.runs\_reflected}). With sequencing added (\lean{Hosting.ArrowComp}) the same three
statements hold for programs whose equations are scoped, with pairwise distinct binder sites within
each equation and within the query (\lean{Hosting.arrow\_runs\_forward\_of\_scoped},
\lean{Hosting.arrow\_answers\_of\_scoped}, \lean{Hosting.arrow\_runs\_reflected\_of\_scoped}); under
those conditions every alternative that a run reaches keeps the names of its pending frames private
(\lean{Hosting.reachableHygienic\_of\_scoped}), which is what the rule that resumes a frame needs.
Two programs are run through the translation to their final bags: the one equation
\mt{(= (app (lam x. (\$F x)) \$a) (\$F \$a))}, a second-order pattern under a binder, used at two
arguments in the two alternatives of a choice
(\lean{Hosting.ExampleChoice.run\_through\_translation}), and a function applied twice, written with
\mt{let} (\lean{Hosting.ExampleTwice.run\_through\_translation}).

The matching is the step's own. The canonical firing of an equation's rule at a represented call is
admitted exactly when a fresh copy of the equation's left side unifies with the call in the
alternative's store (\lean{Hosting.schema\_admitted\_iff}). The frames' equality is equality of
terms, and no equation of the source is part of it (\lean{Hosting.contextual\_equality}); in the
first example the equation's rule is not admitted on \mt{(pair one one)}, a normal form of its right
side (\lean{Hosting.ExampleChoice.normal\_form\_refused}). For the tree's core of programs with
binders, embedded in the source language, rewriting from a closed call reaches exactly the closed
terms that a state reached by a run of the translated program shows
(\lean{Hosting.BinderCoreOnTheStep.binderCore\_on\_the\_step}).

The limits are these.
\begin{itemize}
\item \emph{The fragment.} Spaces, new spaces, types, delayed computations and text run as code are
  outside it. The language has no grounded operations.
\item \emph{Bags.} The bag of answers is read on a frontier, whose members are chosen along the run of
  the source. From the step of the translated program alone the bag is obtained for the first
  fragment only and under a hypothesis on the rules (Section~\ref{sec:addendum-bags}). Two equations
  that differ only in their binder site are one rule (\lean{Hosting.translate\_duplicate}).
\item \emph{A closed system, and its neighbours.} The correspondence is for the translated program
  run from its own start state. The rows of the start state belong to the program as well, and the
  rule for an unmatched call reads them by an absence check, so rows or rules added beside a hosted
  program can change what it answers. Which neighbours keep the answers is stated for the fragment
  with \mt{let}. Rows: if no other row reads as an alternative's row, every run of the translated
  program beside them is a run of the source; if moreover no other row matches an unmatched call, every run of the source is followed,
  and the answers shown are the source's, as a set. Rules: if, at every state the union reaches from
  the start, each firing of the other rules takes no hosted row, extends every solution of the
  knowledge while reading the hosted rows alike, and makes no row that reads as an alternative's
  row, then every run of the union is a run of the source with those firings silent, and the answers
  shown from the start are the source's (\lean{Hosting.Composition.AlongRuns.runs\_reflected\_along},
  \lean{Hosting.Composition.AlongRuns.answers\_along}); the source diverges exactly when, in some run
  of the union, the translated program fires infinitely often
  (\lean{Hosting.Divergence.BesideNeighbours.arrow\_diverges\_iff\_beside}). Knowledge of an equation
  row's private unknown can add answers. The bag of answers and termination are not kept: beside a
  neighbour that fires for ever, every union diverges.
\item \emph{The elaboration.} What is hosted is the source language as its elaboration of program
  text defines it. In that elaboration a call whose head is a variable does not evaluate its
  arguments; the example of a function applied twice is written with \mt{let} by hand.
\end{itemize}

\subsection{The bag of answers from the step alone, and spaces shared by alternatives}\label{sec:addendum-bags}
Every history of a translated program from the start state of a query whose names are unopened is a
chain: each firing takes the alternative's row and makes its next row. This holds for every
translated program, sequencing included
(\lean{Hosting.Alternatives.StepAlone.every\_history\_a\_chain}). A history follows a word of rules,
and two histories that follow one word show the same answer and are finished together
(\lean{Hosting.Alternatives.StepAlone.let\_words\_same\_answer}). Call a word \emph{complete} when
some history follows it to a state with no firing, and \emph{finished} when some history follows it
to a finished state. For a program and a query of the first fragment, with unopened names, whose
equations have pairwise distinct rules, and for every run of the source that ends: the complete
words, each counted once, form a bag, and the bag of answers of the final configuration is the bag of
the answers of its finished words (\lean{Hosting.Alternatives.bag\_of\_rule\_classes}). For such
programs the bag is a function of the translated program, the start state and the identity of the
query's row (\lean{Hosting.Alternatives.StepAlone.bag\_from\_step}), and two runs of the source that
both end show the same bag (\lean{Hosting.Alternatives.answers\_determined}). With words of names of
equations in place of words of rules the hypothesis on the rules is not needed
(\lean{Hosting.Alternatives.bag\_of\_classes}), and the statement holds at every configuration that a
run of the source reaches, not only where a run ends
(\lean{Hosting.Alternatives.StepAlone.partial\_bag}). The names are the positions of the equations in
the source program, which the translated program does not record.

The hypothesis is on rules, not on equations. The equation \mt{(= (f x) x)} written at two sites,
each with its variable named at its own site, gives two rules, and the program of the two satisfies
it (\lean{Hosting.Alternatives.StepAlone.two\_sites\_distinct}). Two copies of a closed equation give
one rule, because a program of the step is a set: one translated program then has two source programs
whose bags are $\{\mathit{one}\}$ and $\{\mathit{one},\mathit{one}\}$
(\lean{Hosting.Alternatives.StepAlone.same\_translation\_two\_bags}), with two classes of histories by
names and one word of rules
(\lean{Hosting.Alternatives.Controls.duplicate\_two\_classes\_one\_rule\_word}). Real HE and the
engine's HE lane answer twice for two copies of an equation, closed or not. So the step of a program
held as a set of rules fixes the bag of a hosted program only where no rule is repeated. Counting a
repeated rule needs rules that are told apart, which the definition of a program does not provide.

The source language's configuration has one world for all alternatives. Read so, the observation of
a query is a set of bags: for the race of an addition to a space with a query of that space, two
bags are observed (\lean{Hosting.Alternatives.Readings.shared\_race}) and no single bag is the
observation (\lean{Hosting.Alternatives.Readings.shared\_not\_one\_bag}). With one world for each
alternative, the moves of two distinct alternatives commute
(\lean{Hosting.Alternatives.Readings.sep\_commute}) and the race has one bag
(\lean{Hosting.Alternatives.Readings.sep\_race}). On that race, with the addition written first and
then with the query written first, real HE answers \mt{[1]} and \mt{[2, 1]} and the engine's HE lane
answers \mt{[1, 2]} and \mt{[1]}: both run the first reading under a fixed schedule, and their
schedules are opposite. Neither reading is hosted on the step here; the fragment of the hosting
theorem has no spaces.

\subsection{Rho with semantic substitution on the step}\label{sec:addendum-rho}
The reflective higher-order calculus is taken as the tree has it on closed processes
(\lean{Mettapedia.Languages.ProcessCalculi.RhoCalculus.LanguageDefGSLT.rhoLanguageDefGSLT}):
communication synchronizes up to name equivalence; substitution is semantic, a dropped received name
becoming the received process; quoted bodies are sealed, so a quote that holds the bound name under its
binder is outside the carrier; and reduction is taken modulo the equations. On the first-order step, whose admitted events these firings are
(\lean{FragmentsOfTheStep.ReflectiveCalculus.FirstOrderInstance.fires\_iff\_firstOrder}), three rule
schemas fixed before any process is given place it: communication takes an input row and an output
row whose channels are one unknown and makes the row $\mathit{react}\,K\,Q$; heating takes a bag and
makes its first component and the rest; clearing takes the null process
(\lean{FragmentsOfTheStep.ReflectiveCalculus.comm},
\lean{FragmentsOfTheStep.ReflectiveCalculus.schema\_renaming}). A store represents a closed process
when its knowledge is solvable and fixes the meaning of every row, every row means a closed canonical
process, and the rows' primes, counted with multiplicity, are the process's canonical bag. Through
this representation the step with the three schemas is a weak bisimulation of the source on all
closed processes, communication firings visible and the others silent: every closed process is
represented, heating and clearing keep the represented process, every communication firing is a
source step, every source step is matched by heating and clearing firings followed by one
communication, and a run of $k$ heating and clearing firings has $k$ at most the store's weight
(\lean{FragmentsOfTheStep.ReflectiveCalculus.reflective\_calculus\_on\_the\_step}). In this placement
the contractum is made by a term former whose meaning is the calculus's own substitution, not by
instantiating the rule. Two stores with no common occurrence and no common rule instance represent the
parallel composition of what they represent
(\lean{FragmentsOfTheStep.ReflectiveCalculus.union\_represents}), and an output barb is read off one row
once heating has split the store into primes
(\lean{FragmentsOfTheStep.ReflectiveCalculus.barbs\_after\_heating}).

The former $\mathit{react}\,K\,Q$ means the canonical form of the semantic substitution of $Q$ into
$K$ (\lean{FragmentsOfTheStep.ReflectiveCalculus.Tm.eval}), and two terms are equal in the knowledge
when they mean the same pattern under every grounding. The knowledge does the matching: a
communication firing on an input row and an output row of a represented store is admitted exactly
when their channels agree, which on canonical rows is name equivalence
(\lean{FragmentsOfTheStep.ReflectiveCalculus.Mechanisms.channel\_match\_is\_solvability}); in the source,
an input on \mt{@(*@0)} and an output on \mt{@0}, different names that are name-equivalent,
communicate, and an input on \mt{@(@0!(0))}, which is not name-equivalent to \mt{@0}, does not
(\lean{FragmentsOfTheStep.ReflectiveCalculus.Source.communicates\_up\_to\_name\_equivalence},
\lean{FragmentsOfTheStep.ReflectiveCalculus.Source.inequivalent\_channels\_do\_not\_communicate}). The
substitution, however, is the former's, not the rule's. Ordinary substitution gives another
process: \mt{for(y <- @0)\{*y\} | @0!(@0!(0))} communicates to \mt{@0!(0)}, ordinary substitution
gives \mt{*@(@0!(0))}, and beside a receiver on \mt{@0} the first communicates and the second is
stuck (\lean{FragmentsOfTheStep.ReflectiveCalculus.Source.semantic\_and\_ordinary\_differ}); the step
reaches the first, after heating, since no communication fires at that program's one-row store
(\lean{FragmentsOfTheStep.ReflectiveCalculus.Mechanisms.separating\_program\_on\_the\_step},
\lean{FragmentsOfTheStep.ReflectiveCalculus.Mechanisms.heating\_is\_needed}). With an input that binds
the received name and a process twin of it, the calculus's contractum is an ordinary substitution
instance (\lean{FragmentsOfTheStep.ReflectiveCalculus.TwinBinders.twin\_contractum}), and without the
twin ordinary substitution gives the inert drop of the code
(\lean{FragmentsOfTheStep.ReflectiveCalculus.TwinBinders.ordinary\_without\_twins}); this is proved for
substitution, not as a placement whose communication rule is instantiated ordinarily. One ground rule
per redex does not suffice: for the family \mt{for(y <- @0)\{*y\} | @0!(}$P_n$\mt{)}, with $P_n$ the
output on \mt{@0} nested $n$ deep, no finite list of ground rules has all the contracta
(\lean{FragmentsOfTheStep.ReflectiveCalculus.Mechanisms.ground\_placement\_is\_not\_finite}). With the
bag laws in the knowledge, splitting a row has two solutions with different first rows
(\lean{FragmentsOfTheStep.ReflectiveCalculus.Mechanisms.bag\_laws\_split\_two\_ways}), so the parallel
structure stays in the multiset of rows.

\paragraph{By ordinary instantiation, on the step of a program.} With an input that binds the received
name and a process twin of it, the calculus is a program of the step of
Definition~\ref{def:program-step}: four rules fixed before any process is given, at the contextual
frames of the twin-binder signature (\lean{FragmentsOfTheStep.ReflectiveCalculus.ByInstantiation.program},
\lean{FragmentsOfTheStep.ReflectiveCalculus.ByInstantiation.theStep}). Communication takes an input row and
an output row on one channel unknown and makes the continuation unknown applied to the quote of the
payload and to the payload; heating and two clearing rules do the housekeeping. Under every grounding
the made row is the term algebra's own instantiation, and no former computes a substitution
(\lean{FragmentsOfTheStep.ReflectiveCalculus.ByInstantiation.made\_row\_is\_instantiation}); read back, it is structurally congruent to the calculus's semantic
substitution and has the canonical form of its contractum (\lean{FragmentsOfTheStep.ReflectiveCalculus.ByInstantiation.contractum\_by\_instantiation},
\lean{FragmentsOfTheStep.ReflectiveCalculus.ByInstantiation.contractum\_by\_instantiation\_canonical}). Equality in the knowledge is the calculus's name equivalence on
closed names and syntax elsewhere; the bag laws are not in it (\lean{FragmentsOfTheStep.ReflectiveCalculus.ByInstantiation.bag\_laws\_not\_in\_knowledge}), and a
communication is admitted exactly when the canonical forms of the two channels agree
(\lean{FragmentsOfTheStep.ReflectiveCalculus.ByInstantiation.comm\_admitted\_iff}). Through a representation, the step of this program is a weak bisimulation of
the source on all closed processes, housekeeping silent and bounded
(\lean{FragmentsOfTheStep.ReflectiveCalculus.ByInstantiation.reflective\_calculus\_by\_instantiation}); the placement by the \mt{react} former is this one
followed by a normal-form map (\lean{FragmentsOfTheStep.ReflectiveCalculus.ByInstantiation.react\_is\_normal\_form}, \lean{FragmentsOfTheStep.ReflectiveCalculus.ByInstantiation.normal\_form\_represents}). The carrier
seals quotes: for a quote that holds the bound name under its binder, instantiation enters the quote
and the calculus's substitution does not (\lean{FragmentsOfTheStep.ReflectiveCalculus.ByInstantiation.outside\_the\_carrier}). This is a correspondence of
reductions; barbs and the criteria for encodings are in Section~\ref{sec:addendum-both-directions}.

\subsection{Absence outside every calculus with a parallel context rule}\label{sec:addendum-separation}
A calculus with a parallel composition has terms, a reduction, a composition meaningful on
composable pairs, a success observation and a static equivalence; reachable success is a weak barb.
It is \emph{persistent} when its reduction is closed beside every term, success persists beside every
term, and reachable success is read up to the static equivalence; it is \emph{switchable} when some
composable neighbour switches a term's reachable success off
(\lean{NullsOnTheStep.ParallelClosed.ParallelCalculus.Persistent},
\lean{NullsOnTheStep.ParallelClosed.ParallelCalculus.Switchable}). A \emph{uniform translation} is
homomorphic on parallel composition, up to the target's static equivalence, on every composable pair,
and preserves and reflects reachable success (\lean{NullsOnTheStep.ParallelClosed.UniformTranslation});
none goes from a switchable calculus into a persistent one
(\lean{NullsOnTheStep.ParallelClosed.no\_uniformTranslation}). Take the first-order step under the
program of two rules, one making \mt{T} where \mt{U} is absent and one reading \mt{U} and making
\mt{W}, with the union of configurations, on disjoint stores with jointly solvable knowledge, as its
composition. It is switchable: alone, \mt{T} is reached, and beside the store holding \mt{U} no run
reaches it (\lean{NullsOnTheStep.ParallelClosed.source\_switchable}), although runs exist there
(\lean{NullsOnTheStep.ParallelClosed.tick\_runs\_beside}); that a persistent blocker keeps success
unreachable along every run is proved over ground terms
(\lean{NullsOnTheStep.ParallelClosed.never\_succeeds\_of\_blocker}).

Every presentation in the tree's form that declares a parallel context rule, a rewrite of a
collection kind whose only premise reduces one element beside the rest, with three distinct names
(\lean{NullsOnTheStep.ParallelClosed.HasParallelContextRule}), has a reduction closed beside every
term, as given and modulo its equations
(\lean{NullsOnTheStep.ParallelClosed.closedBeside\_contextual},
\lean{NullsOnTheStep.ParallelClosed.closedBeside\_moduloEquations}), and is persistent
(\lean{NullsOnTheStep.ParallelClosed.presentationCalculus\_persistent}); rho's presentation is one
(\lean{NullsOnTheStep.ParallelClosed.rhoCalc\_hasParallelContextRule}), and a rule that authorizes the
parallel frame with two of its names equal does not close the reduction
(\lean{NullsOnTheStep.ParallelClosed.ContextRuleControls.withLoop\_not\_closed}). So no relation that a
parallel component switches off is the reduction of such a presentation, the absence step on
patterns among them (\lean{NullsOnTheStep.ParallelClosed.contextual\_not\_switchedOff},
\lean{NullsOnTheStep.ParallelClosed.rhoCalc\_not\_absenceStep}), and no uniform translation goes from
the step under that program into any such presentation read modulo its equations, into the reduction
of any rho rule set of the tree on processes modulo structural congruence, or into any calculus of
positive bag rules (\lean{NullsOnTheStep.ParallelClosed.no\_uniformTranslation\_presentation},
\lean{NullsOnTheStep.ParallelClosed.no\_uniformTranslation\_rho},
\lean{NullsOnTheStep.ParallelClosed.no\_uniformTranslation\_bag}). This separation is for translations
homomorphic on parallel composition only. A map defined on whole configurations escapes it: the one
that sends a configuration to \mt{T} when it can succeed and to nothing otherwise preserves and
reflects reachable success and is not homomorphic
(\lean{NullsOnTheStep.ParallelClosed.Controls.owner\_translation}); it is the characteristic
function of success, not an interpreter. And a presentation can test a whole store: one whose single
rule, applied to a store term, adds \mt{T} exactly when the name \mt{U} occurs nowhere in that
term has no parallel context rule
(\lean{NullsOnTheStep.ParallelClosed.OwnedStore.ownedStore\_step\_iff},
\lean{NullsOnTheStep.ParallelClosed.OwnedStore.ownedStore\_no\_parallelContextRule}); that test is
of the name in the whole term, which is coarser than the absence of a row, and no translation of
the step into this presentation is stated. Absence is not the only way a neighbour refuses a firing
of the step: admission of a firing whose checks are equalities is monotone under the union of
configurations given a solution of the firing's own equations that also solves the neighbour's
knowledge (\lean{NullsOnTheStep.Composition.positive\_admission\_monotone}), and the joint
solvability of the two knowledges, which composability asks, does not supply that solution. So the
separation is not one of absence. The step under a program with no check at all, one rule that reads
a row, tells an equation and makes a row
(\lean{NullsOnTheStep.ParallelClosed.PositiveTell.tellRule\_checks}), is switchable through an
unknown that the neighbour's knowledge also constrains
(\lean{NullsOnTheStep.ParallelClosed.PositiveTell.tellSource\_switchable}), and no uniform translation
carries it into those targets either
(\lean{NullsOnTheStep.ParallelClosed.PositiveTell.no\_uniformTranslation\_presentation},
\lean{NullsOnTheStep.ParallelClosed.PositiveTell.no\_uniformTranslation\_rho},
\lean{NullsOnTheStep.ParallelClosed.PositiveTell.no\_uniformTranslation\_bag}). What these theorems
separate is a calculus with a given composition under this criterion; they do not single out a
feature of the step, and they say nothing about translations that are not homomorphic on parallel
composition. The tree's earlier
encoding results, from bag rules with an absence condition into pattern-level rho and for any feature
that a neighbour can switch off
(\lean{Mettapedia.GSLT.Coordination.ConsensusNumbers.no\_parallel\_encoding},
\lean{Mettapedia.GSLT.Coordination.ConsensusNumbers.rho\_cannot\_encode\_absence},
\lean{Mettapedia.Languages.ProcessCalculi.RhoCalculus.Weaknesses.Monotone.rho\_no\_encoding\_of\_disableable}),
are generalized at both ends: with equality as the static equivalence, the two target hypotheses of
the first are persistence (\lean{NullsOnTheStep.ParallelClosed.persistent\_of\_eq\_iff}). The abstract
argument is short; what the theorems add is its instantiation on the step under a program and on
every presentation with such a rule. The rho targets are the tree's rule-set reductions and rho's
presentation read modulo its equations; the calculus of Section~\ref{sec:addendum-rho} is one of them
by Section~\ref{sec:addendum-both-directions}.

Besides homomorphism, the other two hypotheses on the translation are needed as well
(\lean{NullsOnTheStep.ParallelClosed.Controls.rows\_into\_maker},
\lean{NullsOnTheStep.ParallelClosed.Controls.rows\_into\_inert}), and so is persistence of the target:
the identity is a uniform translation of the step under that program into itself, which is not
persistent (\lean{NullsOnTheStep.ParallelClosed.Controls.identity\_into\_source},
\lean{NullsOnTheStep.ParallelClosed.Controls.source\_not\_persistent}). With one target context for parallel
composition in place of the composition itself, the separation holds when that context keeps
reachable success of its first hole (\lean{NullsOnTheStep.ParallelClosed.no\_compositionalTranslation}),
and a context that guards its hole does not keep it
(\lean{NullsOnTheStep.ParallelClosed.Controls.guarded\_context}); separations under criteria with a
context per operator and divergence reflection, under replacement freeness, or by a resource bound
are not proved. A parallel rule over a binary constructor and arbitrary authored frames are not
covered.

\subsection{The reflective calculus and the one-match step, both directions under one named criterion}\label{sec:addendum-both-directions}
The criterion is the uniform one of Section~\ref{sec:addendum-separation}: a translation homomorphic on
parallel composition up to the target's static equivalence, on composable pairs, that preserves and
reflects reachable success. It is Gorla's homomorphism of parallel composition together with his
success sensitiveness (\lean{FragmentsOfTheStep.ReflectiveCalculus.Encodings.uniform\_iff}), with three differences: it has no renaming clause, success is
an output barb on a closed channel that is not reserved, and homomorphism is asked on the pairs the
source composes. The step is read as a calculus of \emph{systems}, a program with a state, at the
frames of the twin-binder signature: a system reduces by a step of its program; two systems are
composed by the union of their programs and the tensor of their states, which keeps their unknowns
apart; and two systems are composable when their joint knowledge is solvable, which is when each one's
is (\lean{FragmentsOfTheStep.ReflectiveCalculus.Encodings.stepParallelSolvable}, \lean{FragmentsOfTheStep.ReflectiveCalculus.Encodings.tensor\_solvable\_iff}).

Under that criterion, for every closed channel, the reflective calculus embeds into the step and the
step does not embed into it (\lean{FragmentsOfTheStep.ReflectiveCalculus.Encodings.strict\_separation\_solvable}). The embedding sends a closed process to
the four-rule program of Section~\ref{sec:addendum-rho} with the state of the process: it is homomorphic
on parallel composition on the nose, the images of two processes are composable, and it is success
sensitive (\lean{FragmentsOfTheStep.ReflectiveCalculus.Encodings.rhoIntoStepSolvable}, \lean{FragmentsOfTheStep.ReflectiveCalculus.Encodings.encode\_par}, \lean{FragmentsOfTheStep.ReflectiveCalculus.Encodings.image\_composable},
\lean{FragmentsOfTheStep.ReflectiveCalculus.Encodings.encode\_successSensitive}); it has operational correspondence, reflects and preserves divergence, and
the contractum of one communication is correct (\lean{FragmentsOfTheStep.ReflectiveCalculus.Encodings.encode\_operationalCorrespondence},
\lean{FragmentsOfTheStep.ReflectiveCalculus.Encodings.encode\_divergenceReflecting}, \lean{FragmentsOfTheStep.ReflectiveCalculus.Encodings.divergence\_preserved}, \lean{FragmentsOfTheStep.ReflectiveCalculus.Encodings.contractum\_correct}); weak barbs
agree through the representation, and the translation is fully abstract for barbed bisimilarity indexed
by a set of closed names (\lean{FragmentsOfTheStep.ReflectiveCalculus.Encodings.weak\_barbs\_agree}, \lean{FragmentsOfTheStep.ReflectiveCalculus.Encodings.fully\_abstract}). It is not compositional at
output or at input in Gorla's sense (\lean{FragmentsOfTheStep.ReflectiveCalculus.Encodings.output\_not\_compositional}, \lean{FragmentsOfTheStep.ReflectiveCalculus.Encodings.input\_not\_compositional}), and
his renaming clause is not shown. In the other direction, the step is switchable in this composition:
a system with one rule that makes the success row where no row \mt{*@0} is present succeeds, and beside
the composable system that holds that row it never does (\lean{FragmentsOfTheStep.ReflectiveCalculus.Encodings.solvable\_switchable}); the reflective
calculus is persistent (\lean{FragmentsOfTheStep.ReflectiveCalculus.Encodings.rhoParallel\_persistent}); and no translation of the step is success
sensitive and homomorphic up to a relation that respects reachable success, whatever the relation
(\lean{FragmentsOfTheStep.ReflectiveCalculus.Encodings.no\_translation\_any}), while with the relation that relates every two processes the
characteristic function of success is such a translation (\lean{FragmentsOfTheStep.ReflectiveCalculus.Encodings.oracle\_with\_total}). So the calculus of
Section~\ref{sec:addendum-rho} is a target of the separation, for the two sources of
Section~\ref{sec:addendum-separation} as well (\lean{FragmentsOfTheStep.ReflectiveCalculus.Encodings.no\_uniformTranslation\_absence},
\lean{FragmentsOfTheStep.ReflectiveCalculus.Encodings.no\_uniformTranslation\_positiveTell}).

What switches a system off is located. In the restricted composition a system whose rules have only
equation checks and no census pattern is switched off by no composable neighbour, whatever the
neighbour's rows and rules (\lean{FragmentsOfTheStep.ReflectiveCalculus.Encodings.succeeds\_beside}). So in every switching pair the switched system has
a rule with an absence check, a count check or a census pattern (\lean{FragmentsOfTheStep.ReflectiveCalculus.Encodings.switched\_has\_guard}), the image of
a process is never switched off (\lean{FragmentsOfTheStep.ReflectiveCalculus.Encodings.image\_never\_switched}), and the calculus of systems without such
rules is not switchable (\lean{FragmentsOfTheStep.ReflectiveCalculus.Encodings.positiveParallel\_not\_switchable}): for it this argument gives no
separation. These are statements about success that may be reached: with success read on every
maximal run, a system with no such rule is switched off by a neighbour with none
(\lean{FragmentsOfTheStep.ReflectiveCalculus.Encodings.MustReading.must\_reading\_switched}). A census over a ground pattern is switched off by a neighbour
whose row has an unknown of its own, and not by a ground row
(\lean{FragmentsOfTheStep.ReflectiveCalculus.Encodings.GroundCensus.groundCensus\_switched\_by\_free}, \lean{FragmentsOfTheStep.ReflectiveCalculus.Encodings.GroundCensus.groundCensus\_beside\_ground}), and a
count check and a census pattern witness the separation as absence does
(\lean{FragmentsOfTheStep.ReflectiveCalculus.Encodings.Witnesses.separation\_by\_count}, \lean{FragmentsOfTheStep.ReflectiveCalculus.Encodings.Witnesses.separation\_by\_census}). The restriction of composition is
needed: with every pair composable, a neighbour whose knowledge has no solution switches off every
system, the image of a process included (\lean{FragmentsOfTheStep.ReflectiveCalculus.Encodings.junk\_switches}, \lean{FragmentsOfTheStep.ReflectiveCalculus.Encodings.image\_switched\_off}), and the same
criterion then separates the reflective calculus from the states of its own four-rule program
(\lean{FragmentsOfTheStep.ReflectiveCalculus.Encodings.strict\_separation\_fourRule}). Even so, the calculus of systems without such rules is not
persistent in the sense of Section~\ref{sec:addendum-separation}, which asks success beside every term,
composable or not (\lean{FragmentsOfTheStep.ReflectiveCalculus.Encodings.positiveParallel\_not\_persistent}).

Three limits. First, the criterion is homomorphism of parallel composition. For compositionality with
a coordinating context the separation holds when both holes of the context are at top level beside a
coordinating process (\lean{FragmentsOfTheStep.ReflectiveCalculus.Encodings.no\_coordinatedTranslation}, \lean{FragmentsOfTheStep.ReflectiveCalculus.Encodings.step\_not\_into\_rho\_coordinated},
\lean{FragmentsOfTheStep.ReflectiveCalculus.Encodings.positiveTell\_not\_into\_rho\_coordinated}), and the argument fails for a context that guards a hole
(\lean{FragmentsOfTheStep.ReflectiveCalculus.Encodings.guarded\_context\_fails}); that a valid encoding has both holes at top level is not proved for this
source and this target. Second, the composition of systems brings programs, which are sets of rules, and
under this criterion that shape alone separates. The systems with no absence check, count check or
census pattern do not embed into the reflective calculus with its equations either: by counting,
since there are as many such systems as sets of numbers and countably many processes
(\lean{FragmentsOfTheStep.ReflectiveCalculus.Encodings.Counting.succeeds\_iff\_mem}, \lean{FragmentsOfTheStep.ReflectiveCalculus.Encodings.Counting.positive\_not\_into\_rho}); and, with one finite program,
because a program beside itself is the program while the equations of the reflective calculus cancel a
common parallel component, so that every uniform translation forgets the program
(\lean{FragmentsOfTheStep.ReflectiveCalculus.Encodings.Idempotence.program\_blind}, \lean{FragmentsOfTheStep.ReflectiveCalculus.Encodings.Idempotence.positive\_not\_into\_rho\_by\_idempotence}). Barbed
bisimilarity does not cancel (\lean{FragmentsOfTheStep.ReflectiveCalculus.Encodings.Cancellation.bisimilarity\_not\_cancellative}), so for it only the counting
argument is available. Neither argument is about a feature of the step. A comparison that is needs
finite programs, the target's own bisimilarity, and a composition in which a program beside itself is
not the program; that comparison is open. Third, all of it is about the step that commits one match
(Section~\ref{sec:recommendation}); a step that commits several matches at once is not compared here.

\subsection{Bisimilarity of the program-free step, and of programs}\label{sec:addendum-program-bisimilarity}
For every substitution action, on configurations with consistent knowledge and given a fresh supply
of identities for every label, bisimilarity of the step that offers every semantic firing under its
rule is exactly equivalence modulo the knowledge
(\lean{LogicOverTheStep.ProgramFreeBisimilarity.ruleOffers\_bisimilar\_iff}); from bisimilarity to the
equivalence no supply is needed
(\lean{LogicOverTheStep.ProgramFreeBisimilarity.moduloEquivalent\_of\_bisimilar}). Both hypotheses are
needed: one row against two under the knowledge $0\doteq 1$ are bisimilar, since neither moves, and
not equivalent modulo the knowledge
(\lean{LogicOverTheStep.ProgramFreeBisimilarity.ProgramFreeControls.consistency\_needed}); and a pair
equivalent modulo the knowledge whose labels have no fresh supply is not bisimilar, since the rule that
makes two rows \mt{X} moves at one of them only
(\lean{LogicOverTheStep.ProgramFreeBisimilarity.ProgramFreeControls.freshSupply\_needed}). So the
congruence theorem of the program-free step (\lean{Semantic.ruleBisimilar\_tensor}) states that a
static equivalence is preserved by a store disjoint from both sides, and it follows from that static
fact (\lean{LogicOverTheStep.ProgramFreeBisimilarity.ruleBisimilar\_tensor\_of\_static}); it is not a
congruence result for programs.

A program here is a set of rule instances, and its step is the program-free step restricted to them.
Its bisimilarity contains equivalence modulo the knowledge, given a fresh supply
(\lean{LogicOverTheStep.ProgramFreeBisimilarity.programBisimilar\_of\_moduloEquivalent}), and for a
finite program it can be strictly coarser: for the program that takes one \mt{T}, the store with the
row \mt{T} and the store with the rows \mt{T} and \mt{U} are bisimilar and not equivalent modulo the
knowledge (\lean{LogicOverTheStep.ProgramFreeBisimilarity.program\_strictly\_coarser}). Over a ground
action and given a supply of identities, two consistent configurations are bisimilar for a program
whose rules only take and make exactly when every word of the program needs a bag below the bag of
labels of the first exactly when it needs one below that of the second, a word's need being the least
bag from which it fires (\lean{LogicOverTheStep.ProgramFreeBisimilarity.takeMake\_bisimilar\_iff}).
Bisimilarity for a program that takes one row of one label is preserved by a disjoint store with
consistent composites (\lean{LogicOverTheStep.ProgramFreeBisimilarity.takeOne\_tensor}); for a program
that takes two rows together it is not: under the rule that takes \mt{T} and \mt{U} together, one row
\mt{T} and no rows are bisimilar, and beside a disjoint row \mt{U} they are not
(\lean{LogicOverTheStep.ProgramFreeBisimilarity.takeTU\_not\_congruence}). These are data for an
equivalence of programs preserved in context, not that equivalence, which is the subject of
Section~\ref{sec:addendum-congruence}.

\subsection{An equivalence of programs that every context preserves, and what it is}\label{sec:addendum-congruence}
A system is a program with a state of its step (\lean{RuleSyntax.Equivalence.System}), and a
reduction of a system is an admitted firing of its program. Fix a family of observations of states;
the one used below is the may-barb: some present row and a pattern have a common instance at some
solution of the knowledge (\lean{RuleSyntax.Equivalence.MayBarb}). A basic context adds rules, tells
equations, extends the frame, or adds rows. Two systems are \emph{congruent} when some symmetric
relation contains them that preserves the observations, answers every reduction by a reduction into
the relation, and is closed under every basic context, the identities of added rows being chosen
afresh on the answering side (\lean{RuleSyntax.Equivalence.Candidate},
\lean{RuleSyntax.Equivalence.Congruent}). For every frame action, every type of extensions and every
family of observations this is an equivalence relation
(\lean{RuleSyntax.Equivalence.congruent\_equivalence}), preserved by adding rules, telling
equations, extending the frame, adding rows, every composite of these, and running beside another
system (\lean{RuleSyntax.Equivalence.Congruent.addRules}, \lean{RuleSyntax.Equivalence.Congruent.tell},
\lean{RuleSyntax.Equivalence.Congruent.extend}, \lean{RuleSyntax.Equivalence.Congruent.addRows},
\lean{RuleSyntax.Equivalence.Congruent.composite}, \lean{RuleSyntax.Equivalence.Congruent.beside}). It
is sound for contextual equivalence up to the identities of added rows, and contained in
bisimilarity without contexts (\lean{RuleSyntax.Equivalence.congruent\_contextualUpTo},
\lean{RuleSyntax.Equivalence.Congruent.bisimilar}). When every admitted firing of $P_1$ is answered at
the same state by an admitted firing of $P_2$ with a congruent result, and conversely, $P_1\cup Q$
and $P_2\cup Q$ are congruent at every state, for every $Q$
(\lean{RuleSyntax.Equivalence.replace\_congruent}).

At one frame of first-order terms without function symbols, with the natural numbers as identities
and may-barbs, the congruence is characterized. Two states are \emph{statically equivalent} when
both knowledges are unsolvable, or some renaming of the identities makes the present rows
correspond, the knowledges have the same solutions, each label equals the label of its image at
every solution, and the two states mention the same unknowns
(\lean{RuleSyntax.Equivalence.Characterization.Static},
\lean{RuleSyntax.Equivalence.Characterization.static\_iff\_clauses}). Beside a program that leaves a
constant unused, as every finite program does
(\lean{RuleSyntax.Equivalence.Characterization.bounded\_of\_finite}), two states are congruent
exactly when they are statically equivalent
(\lean{RuleSyntax.Equivalence.Characterization.congruent\_iff\_static}). Two such programs are
congruent at every state exactly when each answers every admitted firing of the other, at every
state, by an admitted firing with a statically equivalent result
(\lean{RuleSyntax.Equivalence.Characterization.programs\_iff}); the sufficient half holds for every
pair of programs and beside every further program
(\lean{RuleSyntax.Equivalence.Characterization.congruent\_of\_answersStatic}). The condition on the
programs is sufficient and not necessary: without it the characterization of states can fail, already
for positive programs (\lean{RuleSyntax.Equivalence.Characterization.bounded\_needed}), and it holds
beside two programs that do not satisfy it
(\lean{RuleSyntax.Equivalence.Characterization.Unbounded.unbounded\_characterized}). Two states whose
knowledges are both unsolvable are statically equivalent whatever their rows;
\lean{RuleSyntax.Equivalence.Characterization.unsolvable\_example} is such a pair whose rows no
renaming relates. No least enabling context exists, for a rule that checks absence and equally for a
positive program of two rules
(\lean{RuleSyntax.Equivalence.Characterization.no\_least\_enabling\_tell},
\lean{RuleSyntax.Equivalence.Characterization.LeastTell.no\_least\_enabling\_tell\_positive}); what is
proved in place of a semantics with least labels is the characterization over every reachable context
(\lean{RuleSyntax.Equivalence.Characterization.congruent\_iff\_reachable}).

Three limits say what this congruence is. First, it is strong, and it records every unknown that a
firing opened: a rule that takes a row and makes the same row is not congruent to the empty program
(\lean{RuleSyntax.Equivalence.Characterization.touch\_separates}), the result of its firing is not
statically equivalent to the state it fired at
(\lean{RuleSyntax.Equivalence.Characterization.Inlining.touch\_result\_not\_static}), and ``tell
$x\doteq 5$, make $x$'' is not congruent to ``make $5$''
(\lean{RuleSyntax.Equivalence.Characterization.Inlining.inline\_separates}). Second, every unknown is
public: a state that knows $w\doteq 5$ and holds the row $w$, and a state that knows $w\doteq w$ and
holds the row $5$, mention the same unknowns, have the same rows, have equal labels at every
solution of the first, and are congruent beside no pair of programs
(\lean{RuleSyntax.Equivalence.Characterization.Inlining.private\_unknown\_separates}). A context may
tell an equation on any unknown, and nothing here gives an unknown a scope. Third, it does not
count: a rule and a renamed copy of it are one program for the congruence, beside every program and
at every state
(\lean{RuleSyntax.Equivalence.Characterization.CongruentTranslations.variant\_congruent},
\lean{RuleSyntax.Equivalence.Characterization.CongruentTranslations.added\_variant\_congruent}). So it
does not keep bags. Two programs of the hosted source language, both within the hypotheses of the bag
theorem of Section~\ref{sec:addendum-bags}, have translations that are congruent at every state for
every family of observations, and answer $\{\mathit{one}\}$ and $\{\mathit{one},\mathit{one}\}$
(\lean{RuleSyntax.Equivalence.Characterization.CongruentTranslations.congruent\_translations\_two\_bags});
at one and the same state the two systems are congruent and one has a complete history, finished and
showing an answer, that the other lacks
(\lean{RuleSyntax.Equivalence.Characterization.CongruentTranslations.congruent\_systems\_one\_class\_less}).
The congruence answers the question whether a context can tell two systems apart by what may be
observed. The observation of a hosted MeTTa program is its bag of answers, and an equivalence of
hosted programs that keeps it must count; that equivalence is not built.

\subsection{Definitions by rows}\label{sec:addendum-definitions}
A definition over the vectors package of Section~\ref{sec:addendum-indexed} is written in a small
source form: a name, a declared telescope and result type, and rows of one pattern per parameter
with a body; the checker does not read MeTTa text. It elaborates each row into a firing of the
converting step, with the function at its declared type, the pattern variables and the call's index
as the frame, what the patterns teach told as the definition of the index, and the body as the made
row at the declared type; it synthesizes the body's typing, and answers \emph{admitted},
\emph{refused} at a row, or \emph{not decided}
(\lean{NativeTypesOnTheStep.Computing.Vectors.Checker.checkDefinition}). Admitted means that every row
elaborates, its firing is admitted by the computable converting step, and its made row is typed after
its match (\lean{NativeTypesOnTheStep.Computing.Vectors.Checker.checkDefinition\_admitted}); refused
means that the refused row's firing is refused by the converting step for every settler
(\lean{NativeTypesOnTheStep.Computing.Vectors.Checker.checkDefinition\_refused}). Elaboration covers one
call index per row, \mt{cons} with variable sub-patterns, vector lengths that are parameters matched by
variables, and no frame type that mentions the index; refutation reaches bodies that are variables,
constants or pairs. Outside these classes the verdict is not decided, which is never a refusal and can
hide a real error: the recursive call of \mt{append} with its two vectors swapped is not decided
(\lean{NativeTypesOnTheStep.Computing.Vectors.Checker.swap\_undecided}), and neither is the abstraction
body of \mt{konst}, whose row every settler admits
(\lean{NativeTypesOnTheStep.Computing.Vectors.Checker.konst\_undecided},
\lean{NativeTypesOnTheStep.Computing.Vectors.Checker.konst\_every\_settler}). By kernel evaluation at
budget 16, \mt{append} is admitted, \mt{(= (append \$xs \$ys) \$xs)} is refused at its first row, and
\mt{head} on a successor length and the interpreter's literal row are admitted
(\lean{NativeTypesOnTheStep.Computing.Vectors.Checker.append\_admitted},
\lean{NativeTypesOnTheStep.Computing.Vectors.Checker.keep\_refused},
\lean{NativeTypesOnTheStep.Computing.Vectors.Checker.head\_admitted},
\lean{NativeTypesOnTheStep.Computing.Vectors.Checker.eval\_admitted}); at budget 0 \mt{append} is not
decided (\lean{NativeTypesOnTheStep.Computing.Vectors.Checker.append\_budget}), and read as a budgeted
procedure the checker keeps a verdict once reached at every larger budget
(\lean{NativeTypesOnTheStep.Computing.Vectors.Checker.procedure\_resolves}). An empty list of rows is
admitted: admission says nothing about coverage, which is stated by the index. \mt{append} covers
every length and \mt{head} does not, and at length zero a ground call of \mt{head} has no answer by its
rows (\lean{NativeTypesOnTheStep.Computing.Vectors.Checker.append\_covers},
\lean{NativeTypesOnTheStep.Computing.Vectors.Checker.head\_not\_covers},
\lean{NativeTypesOnTheStep.Computing.Vectors.Checker.head\_zero\_unanswered}).

The guarantee is one firing deep. If a definition is admitted and the function's declared rows are
among its rows, every answer of a ground call by those rows has the call's declared type and is the
made row's only conversion normal form
(\lean{NativeTypesOnTheStep.Computing.Vectors.Checker.guarantee}); in a fired row the recursive call is
a call of the function variable, which the ground call fills with a closed term of the declared type.
In the witness, the cons row of \mt{append} on \mt{append num 1 1 [2] [3]}, with that variable filled
by the tree's eliminator program, answers \mt{[2, 3]}
(\lean{NativeTypesOnTheStep.Computing.Vectors.Checker.guarantee\_witness}). That the declared rows are
among the checked ones is an assumption about the declaration, and it is needed: with
\mt{(= (append \$xs \$ys) \$xs)} declared as a further, unchecked row, a ground call answers \mt{[]},
which is not of its declared type
(\lean{NativeTypesOnTheStep.Computing.Vectors.Checker.second\_row\_breaks}). The checker does not prove
that the rows define a function whose every ground call answers at its type.

A definition is then declared as a new constant whose rows join the package's computation. For every
list of definitions, each declared as a new constant, that enters no name twice and whose rules pass
the tree's declaration-time check together with the package's fourteen rows
(\lean{NativeTypesOnTheStep.Computing.Definitions.devCheck}; a name entered twice is refused,
\lean{NativeTypesOnTheStep.Computing.Definitions.twice\_devCheck}, so the declaration of a name and
its rules come from one entry, \lean{NativeTypesOnTheStep.Computing.Definitions.devCheck\_find}), the
extended package is Church--Rosser
(\lean{NativeTypesOnTheStep.Computing.Definitions.dev\_churchRosser}), every step of the package is a
step of the extension (\lean{NativeTypesOnTheStep.Computing.Definitions.pkg\_step\_dev}), and two terms
that mention no new constant and convert in the extension already convert in the package
(\lean{NativeTypesOnTheStep.Computing.Definitions.dev\_conservative}). Addition by rows with
\mt{append} over it passes the check and overlapping rows do not
(\lean{NativeTypesOnTheStep.Computing.Definitions.addAppend\_devCheck},
\lean{NativeTypesOnTheStep.Computing.Definitions.overlap\_devCheck}), and in the extended package
\mt{append num 1 1 [2] [3]} has exactly one answer, \mt{[2, 3]}
(\lean{NativeTypesOnTheStep.Computing.Definitions.append\_answer}). The extended package is not proved
qualified: typed right sides, termination and the algorithmic constants of the extension are missing,
so a call of a new constant has no type guarantee there yet. Its interpretation in the unextended
package is a compiled term built from the recursor of the numbers: for a definition whose verdict is a
function, the row checker admits it, the compiled term is certified at the declared type, and every
row holds as a conversion at every instance of its pattern
(\lean{NativeTypesOnTheStep.Computing.Definitions.qualify\_function}). Every answer of a call of a
closed checked function, at arguments typed along its telescope, is of the call's declared type, is
typed-equal to the call, and is the only answer
(\lean{NativeTypesOnTheStep.Computing.Definitions.call\_answer}). That such a call has an answer is
proved for addition by rows at a numeral and a normal number
(\lean{NativeTypesOnTheStep.Computing.Definitions.add\_exactly\_one}) and is not stated here in general.
An earlier reduction of it to one hypothesis, that every closed typed term of the package has an
answer, is withdrawn, because that hypothesis is false: an answer in this sense is a term at which no
step of the package's reduction starts, that reduction includes the equality of universe heads, which
is reflexive, and so a universe, a closed typed term, has no answer
(\lean{Interface.ComputationProfile.HeadBoundary.full\_answer\_totality\_false}). Existence is proved
on the directed reduction, the package's reduction without the equalities of heads
(\lean{NativeTypesOnTheStep.Computing.DirectedAnswers.DirectedAnswer}), at calls whose result is data.
A closed type is a data type when it is the numbers, the unit type, the empty type, a pair type whose
second component does not depend on the first and whose components are data types, or typed-equal to
one of these (\lean{NativeTypesOnTheStep.Computing.DirectedAnswers.DataType}); vectors of every closed
length over a small data type are data types
(\lean{NativeTypesOnTheStep.Computing.DirectedAnswers.vec\_dataType}). For a definition whose compiled
term is certified as a closed function, and closed arguments typed along its telescope at which the
declared result type is a data type, the compiled term applied to the arguments has exactly one
directed answer, a value of that type typed-equal to the call
(\lean{NativeTypesOnTheStep.Computing.DirectedAnswers.ground\_call\_answer}), and a numeral when the
declared result type is the numbers
(\lean{NativeTypesOnTheStep.Computing.DirectedAnswers.ground\_call\_numeral}). The condition is on the
call and not on the definition: the declared result type of \mt{append} is open, and its calls at a
data type of elements satisfy it
(\lean{NativeTypesOnTheStep.Computing.DirectedAnswers.appendR\_answer}); addition by rows is the first
instance (\lean{NativeTypesOnTheStep.Computing.DirectedAnswers.add\_answer}). The proof uses the
normalization model that the package's qualification provides and no hypothesis of strong
normalization (\lean{NativeTypesOnTheStep.Computing.DirectedAnswers.inductive\_reaches},
\lean{NativeTypesOnTheStep.Computing.DirectedAnswers.data\_answer}). The statement is about the
compiled term in the unextended package. The rows' own reduction in the extended package reaches the
same answer for a class of recursive definitions
(\lean{NativeTypesOnTheStep.Computing.DirectedAnswers.rows\_reach}): a definition whose compiled term is
certified as a closed function, whose rules for its constant are a rule at \mt{zero} and a rule at
\mt{succ k} on one call index, each joining its other side when the constant is read as the compiled
term, the rule at \mt{zero} reading as itself, and every occurrence of the constant in the rule at
\mt{succ k} a call with all the definition's parameters and \mt{k} at the index. At closed arguments
typed along the telescope at which the declared result type is a data type, and for a call whose
arguments are only convertible with them, some reduction of the call reaches a value; it is the only
normal form that any reduction of the call reaches, and it is the one answer of the compiled call in
the package. Nothing is said of reductions that do not end: an argument convertible with
\mt{zero} may itself have an infinite reduction. Addition, append, map and a sum with an accumulator
are instances. Separately, for every checked development whose rules join at the compiled terms,
every value that the rows reach from a closed term is the answer of that term's reading in the
package, with no typing asked (\lean{NativeTypesOnTheStep.Computing.DirectedAnswers.value\_of\_reading}).
At function types existence is not proved. In a constructor presentation whose heads are
related to themselves, an answer of the full reduction is exactly a directed answer without a head
(\lean{NativeTypesOnTheStep.Computing.DirectedAnswers.answer\_iff\_headFree}). Admitted rows
need not compile to
a certified term: $\mathit{big} : \Pi\,(n : \mathrm{num}).\,U_1$ has admitted rows, and its compiled
term is certified neither at budget 16 nor at 64
(\lean{NativeTypesOnTheStep.Computing.Definitions.big\_rows\_admitted},
\lean{NativeTypesOnTheStep.Computing.Definitions.big\_verdict}).

The comparison of conversion has no cliff where its reducer has none. For every package with the
facts under which the algorithm derives exactly the typed equalities, and every reducer that keeps its
answers and on every typed term eventually returns a weak-head normal reduct
(\lean{NativeTypesOnTheStep.Computing.Eventually.WhnfEventually}), the executable comparison accepts
every pair of equal typed terms at some budget, with the same certificate at every larger one
(\lean{NativeTypesOnTheStep.Computing.Eventually.compare\_eventually}). With the tree's normalizer as
the reducer this is proved for the tower and, for every complete root evaluator, for the executable
package (\lean{NativeTypesOnTheStep.Computing.Eventually.tower\_compare\_eventually},
\lean{NativeTypesOnTheStep.Computing.Eventually.executable\_compare\_eventually}); the tree defines
no root evaluator for the executable package, so the second statement is conditional on one. For the vectors
package it is proved under strong normalization of its typed terms, which is a hypothesis and is not
proved (\lean{NativeTypesOnTheStep.Computing.Eventually.vectors\_compare\_eventually}). Under that
hypothesis the round's conversion scan, and the computable admission of a firing without guards or
queries that the families settle and whose made rows have their frame types along the answer, stop
depending on the conversion budget from some budget on
(\lean{NativeTypesOnTheStep.Computing.Eventually.vectors\_scan\_eventually},
\lean{NativeTypesOnTheStep.Computing.Eventually.vecAdmits\_eventually}). That the hypothesis is needed is not
proved: in a tower with one constant whose row rewrites it to itself, the comparison answers an
equality at no budget and strong normalization fails, but that constant has no weak-head normal
form either, so the example shows only that normal forms must exist
(\lean{NativeTypesOnTheStep.Computing.Eventually.loop\_cliff},
\lean{NativeTypesOnTheStep.Computing.Eventually.loop\_not\_typedSN}). No fixed budget suffices: at the
type $(\lambda x.\,x)^k\,\mathrm{num}$ the pair $\mathrm{zero},\mathrm{zero}$ is unanswered at budget $k$
and answered from $k+3$ on (\lean{NativeTypesOnTheStep.Computing.Eventually.vectors\_cliff},
\lean{NativeTypesOnTheStep.Computing.Eventually.vectors\_cliff\_transient}), and for $k=16$ the kernel
evaluates it unanswered at the tree's fuel 16 and answered at 17
(\lean{NativeTypesOnTheStep.Computing.Eventually.cliff16}). The comparison cannot refuse: on typed
terms that are not equal it returns nothing at every budget, as on $U_0$ and $U_1$
(\lean{NativeTypesOnTheStep.Computing.Eventually.u0\_u1\_unanswered}), so a consumer that reads nothing
as different reads exhaustion as refusal; and the checker accepts the first projection of a pair
literal at no budget, although it is typed
(\lean{NativeTypesOnTheStep.Computing.Eventually.fst\_pair\_never\_checked}). The proofs give a
sufficient budget, not a computable one, except $k+3$ on that family, and the class of terms the
checker accepts is not characterized.

\subsection{Inductive types on the step}\label{sec:addendum-inductive}
On the ground first-order step, take as parties the tree's processes with rights and installed
schemas, a move of a party being a firing of one of its schemas formed under its \mt{send} rights. A
type given by constructors is then a region of rows. \emph{Membership}: when every party writes the
type only through its constructors, that is, wherever it holds \mt{send} in the type its rules make
rows of the type only as conclusions of declared constructors whose premises they take or read, every
row of the region in every state reached is generated, with any number of such writers
(\lean{LogicOverTheStep.InductiveTypes.membership}). The condition is sufficient and not exact: a
writer that makes the row of $2$ with no premise present keeps membership
(\lean{LogicOverTheStep.InductiveTypes.Numbers.withShortcut\_membership},
\lean{LogicOverTheStep.InductiveTypes.Numbers.shortcut\_skips\_premise}). \emph{Induction}: under the
same condition, a property closed under the constructors holds of every row of the region in every
state reached from the empty store (\lean{LogicOverTheStep.InductiveTypes.induction}); a writer that
makes a row no constructor makes breaks it
(\lean{LogicOverTheStep.InductiveTypes.Numbers.withJunk\_breaks\_induction}). \emph{Closure}, ``these
and no others'', has two exact statements. Every closure statement survives every other party exactly
when the region is owned, and a sole \mt{send} right in the type owns it
(\lean{LogicOverTheStep.InductiveTypes.closure\_of\_soleSender}); a given closure statement survives
every move, from every consistent ground state, of a party that writes through the constructors and
realizes them exactly when its list is closed under the constructors
(\lean{LogicOverTheStep.InductiveTypes.survives\_iff\_ruleClosedIn}). A second writer through the
constructors keeps membership and breaks the closure statement that the only number present is $0$
(\lean{LogicOverTheStep.InductiveTypes.Numbers.twoWriters\_membership},
\lean{LogicOverTheStep.InductiveTypes.Numbers.twoWriters\_breaks\_closure}). A party that realizes the
constructors reaches every generated row from every consistent ground state
(\lean{LogicOverTheStep.InductiveTypes.reachable\_of\_generated}), so none of these statements is
vacuous.

Recursion is proved for one function, addition by structural recursion, as three rules run under a
program on the first-order step. In a consistent ground state where the request is the only work row
and its arguments are generated rows, the numerals of $a$ and $b$, every run has at most $a+2$ moves,
every final state reached holds exactly one work row, the answer $a+b$, and a run of $a+2$ moves
reaches one
(\lean{LogicOverTheStep.InductiveTypes.Recursion.addition\_on\_generated}); read from the stage, a run
has at most $m+1$ moves when the first argument is in stage $m$
(\lean{LogicOverTheStep.InductiveTypes.Recursion.addition\_bounded\_by\_stage}); on a row no
constructor makes, the run stops after one move with no answer
(\lean{LogicOverTheStep.InductiveTypes.Recursion.junk\_request\_no\_answer}). A general recursor rule
set is not built. Against the model's numbers, in the tower with addition: the reading of rows into the
model is defined on numerals only, so that a row is generated exactly when its content reads as a
closed canonical number holds by the reading's construction
(\lean{LogicOverTheStep.InductiveTypes.Agreement.generated\_iff\_closedCanonical}); the content is
in what follows. Every closed term
typed at $\mathrm{num}$ is equal to a numeral (\lean{LogicOverTheStep.InductiveTypes.Agreement.canonicity}),
hence to the reading of exactly one generated row
(\lean{LogicOverTheStep.InductiveTypes.Agreement.closed\_term\_has\_row}); and on the rows of any two
closed terms of $\mathrm{num}$ the step's addition answers the numeral that their sum in the model
equals (\lean{LogicOverTheStep.InductiveTypes.Agreement.answer\_agrees\_closed}). Vectors are treated
on closed rows only, where the generated vector rows are exactly the generated lists of numbers at the
numeral of their length (\lean{LogicOverTheStep.InductiveTypes.Families.vecRow\_generated\_iff\_extrinsic}).
With logic variables in rows, and on the step of a program, these facts are not yet stated.

\subsection{Typed syntax and inductive families in the model}\label{sec:addendum-families}
A typed interpreter's term type has several constructors at one index, among them constructors whose
arguments sit at the index of their result, as addition of two number expressions does. A family
declared by recursion on its index cannot refer to itself there: a constant that a package does not
declare occurs in no derivable judgment
(\lean{Mettapedia.TypeTheory.Calculi.ParameterizedPiSigmaId.Presentation.TypedEquality.Normalization.TowerSyntaxLimit.free\_of\_undeclared}),
so at an index constructor without fields such a family steps to a term free of itself
(\lean{Mettapedia.TypeTheory.Calculi.ParameterizedPiSigmaId.Presentation.TypedEquality.Normalization.TowerSyntaxLimit.nullary\_unfolds}).
An extrinsic encoding, which serves as a control, gets several constructors at one index without a new
case of the model's logical relation: an expression of type $t$ is a pair of an untyped expression and
a proof of a typing relation computed by recursion on it, in a package whose model hypotheses are
proved. There \mt{exp tnum} has the two constructors \mt{lit} and \mt{add}, the interpreter is typed,
and \mt{eval (add (lit 2) (lit 3))} is $5$ by conversion
(\lean{Mettapedia.TypeTheory.Calculi.ParameterizedPiSigmaId.Presentation.TypedEquality.Normalization.TowerSyntaxModel.eval\_typed},
\lean{Mettapedia.TypeTheory.Calculi.ParameterizedPiSigmaId.Presentation.TypedEquality.Normalization.TowerSyntaxModel.eval\_two\_plus\_three}),
and the type former is injective at every index, open ones included
(\lean{Mettapedia.TypeTheory.Calculi.ParameterizedPiSigmaId.Presentation.TypedEquality.Normalization.TowerSyntaxModel.exp\_injective}).
It is not an inductive family: at a variable index it has an element $(\mathit{rnum}\,0, p)$ with a
neutral proof $p$
(\lean{Mettapedia.TypeTheory.Calculi.ParameterizedPiSigmaId.Presentation.TypedEquality.Normalization.TowerSyntaxModel.junk\_typed}),
and two typing proofs of one untyped expression give elements that are not equal by conversion
(\lean{Mettapedia.TypeTheory.Calculi.ParameterizedPiSigmaId.Presentation.TypedEquality.Normalization.TowerSyntaxModel.same\_raw\_not\_equal})
and are equal in the identity type in the one case proved
(\lean{Mettapedia.TypeTheory.Calculi.ParameterizedPiSigmaId.Presentation.TypedEquality.Normalization.TowerSyntaxModel.same\_raw\_propositionally\_equal}).
Its recursion principle is not dependent and its motive lies in the smallest universe, and its
language has numbers, Booleans, addition, a conditional and a zero test, with no binders. Read as
computing rows of the knowledge's conversion, the package is qualified
(\lean{NativeTypesOnTheStep.Computing.Syntax.syntax\_qualified}), and the interpreter's addition and
literal rows are admitted on the converting step and refused on the dependent full step (\lean{NativeTypesOnTheStep.Computing.Syntax.add\_admitted},
\lean{NativeTypesOnTheStep.Computing.Syntax.add\_present\_refused},
\lean{NativeTypesOnTheStep.Computing.Syntax.lit\_admitted},
\lean{NativeTypesOnTheStep.Computing.Syntax.lit\_present\_refused}).

The model's own former is begun as a first stage. One family pack, defined before the logical
relation, reads a family at its parameters: constructor forms, value fields, recursive fields at
indices computed from the fields, targets, and the packs of the index types, with reducibility defined
at all indices at once
(\lean{Mettapedia.TypeTheory.Calculi.ParameterizedPiSigmaId.Presentation.TypedEquality.Normalization.Families.familyPack}).
Both existing packs are instances of it, as equalities of separately defined packs: the simple
inductive pack, under the premises of the inductive case of the relation
(\lean{Mettapedia.TypeTheory.Calculi.ParameterizedPiSigmaId.Presentation.TypedEquality.Normalization.Families.simple\_pack\_eq}),
and the identity pack, read as the family with one constructor \mt{refl x} that stores its point,
under the premises of the identity case
(\lean{Mettapedia.TypeTheory.Calculi.ParameterizedPiSigmaId.Presentation.TypedEquality.Normalization.Families.identity\_pack\_eq}).
The laws a case of the relation asks hold for the pack with two hypotheses that families add: the
family's type respects reducible equality of indices, and reducibly equal fields give reducibly equal
recursive indices
(\lean{Mettapedia.TypeTheory.Calculi.ParameterizedPiSigmaId.Presentation.TypedEquality.Normalization.Families.TypeCongr},
\lean{Mettapedia.TypeTheory.Calculi.ParameterizedPiSigmaId.Presentation.TypedEquality.Normalization.Families.RecursiveExt};
for symmetry and transitivity,
\lean{Mettapedia.TypeTheory.Calculi.ParameterizedPiSigmaId.Presentation.TypedEquality.Normalization.Families.FamEqTm.symm},
\lean{Mettapedia.TypeTheory.Calculi.ParameterizedPiSigmaId.Presentation.TypedEquality.Normalization.Families.FamEqTm.trans}).
Declared as families, vectors have no canonical form at a variable length
(\lean{Mettapedia.TypeTheory.Calculi.ParameterizedPiSigmaId.Presentation.TypedEquality.Normalization.FamilyVectors.neutral\_at\_variable}),
and in a concrete setting \mt{lit x} is reducible at \mt{tnum} and the conditional at a variable type,
while a literal at a variable type is not
(\lean{Mettapedia.TypeTheory.Calculi.ParameterizedPiSigmaId.Presentation.TypedEquality.Normalization.FamilyExamples.lit\_reducible},
\lean{Mettapedia.TypeTheory.Calculi.ParameterizedPiSigmaId.Presentation.TypedEquality.Normalization.FamilyExamples.ite\_reducible\_at\_variable},
\lean{Mettapedia.TypeTheory.Calculi.ParameterizedPiSigmaId.Presentation.TypedEquality.Normalization.FamilyExamples.lit\_not\_reducible\_at\_variable}):
the element that the extrinsic encoding has at a variable index does not arise. The pack is not yet a
case of the logical relation, and no arm of the general development is written for it; the concrete
setting gives the family constants a rigid role, so it exercises the pack, not the relation; value
fields whose types mention earlier fields and indices that depend on earlier indices are excluded at
this stage; and that the vector eliminator's rules preserve typing is not proved.

\subsection{The models of HE and PeTTa, for any grounded table}\label{sec:addendum-tables}
The every-run theory of the models of HE and PeTTa (Section~\ref{sec:addendum-nov12}) is now stated
once for any grounded table, a record of the operators the table executes and its outcome relation,
under a per-call contract: on every call of an operator it executes, the table answers one data
value, declines, or errs
(\lean{DialectsAsStrategies.TriangleFootprints.GroundedTables.GroundedTable.Contract}). The model's
table and the table of HE's grounded operations as measured meet it
(\lean{DialectsAsStrategies.TriangleFootprints.GroundedTables.model\_contract},
\lean{DialectsAsStrategies.TriangleFootprints.GroundedTables.measured\_contract}). The grounded case of
the theorems asks the table to answer each grounded call of the class with PeTTa's value: the model's
table does so for every integer operation, the measured table when the result fits in 64 bits
(\lean{DialectsAsStrategies.TriangleFootprints.GroundedTables.measured\_answers}), and on
\mt{(+ 9223372036854775807 1)} the two tables' verdicts part
(\lean{DialectsAsStrategies.TriangleFootprints.GroundedTables.verdicts\_part}). For a program of the
class and a term that settles on a table in this sense, a complete run of HE's machine with the table
exists, and every complete run publishes, as a bag, the answers the PeTTa model gives from some fuel
on (\lean{DialectsAsStrategies.TriangleFootprints.GroundedEveryRun.answer\_bags\_every\_run\_G}); the
induction lemmas behind such theorems state one list, in the models' order, and HE's floats are not
modelled. The model's machine reads no argument type: for a table that meets the contract, two programs that leave
the operator undeclared publish the same list on every complete run of a call, whatever they declare
about its arguments
(\lean{DialectsAsStrategies.TriangleFootprints.GroundedDifferences.machine\_reads\_no\_argument\_types}).
HE checks a call of a grounded operation against the operation's type before running it, so the region
in which the model's untyped grounded calls can be HE's is stated as a decidable predicate read off
argument values, a conjunct of the theorems that use it
(\lean{DialectsAsStrategies.TriangleFootprints.GroundedTables.TypedRegion}). Eight calls outside it and
four inside it are theorems about the model's answers
(\lean{DialectsAsStrategies.TriangleFootprints.GroundedDifferences.Controls.outside\_region},
\lean{DialectsAsStrategies.TriangleFootprints.GroundedDifferences.Controls.inside\_region}); measured on
hyperon-experimental 0.2.10 with three later fixes, HE answers the eight with argument-type or arity
errors and the four as the model does. Agreement inside the region is measured, not proved.

Three further differences are theorems about the models, each beside a measured row. On a name with
two declared rows, HE's evaluator answers once for each applicable row and the PeTTa model runs the
call once (\lean{DialectsAsStrategies.TriangleFootprints.ThreeDifferences.two\_rows\_models}); measured,
HE answers once and CeTTa 1.5.0-dev in its PeTTa mode once for each row, so each model answers as the
other dialect's host. On partial application, every complete run of HE's machine on
\mt{(co7-wrap 1 2)} publishes the call itself, and with the equation \mt{(= (co7-wrap 1 2) added)}
added it publishes \mt{added}, as HE does
(\lean{DialectsAsStrategies.TriangleFootprints.ThreeDifferences.partial\_application\_he}); on
\mt{(co7-wrap 1 2)} without that equation the PeTTa model stops at a fault with no value, where CeTTa's
PeTTa mode answers \mt{old}, and with the equation it answers \mt{added}, so on that host the added
equation removes an answer
(\lean{DialectsAsStrategies.TriangleFootprints.ThreeDifferences.partial\_application\_petta}). An
\mt{Atom} parameter used twice with a choice as argument gives four pairs on HE's evaluator and
\mt{(pair (two) (two))} on the PeTTa model, as on their hosts
(\lean{DialectsAsStrategies.TriangleFootprints.ThreeDifferences.atom\_parameter\_models}); for that
program and for \mt{(co7-wrap 1)} only the answers of HE's evaluator are theorems
(\lean{DialectsAsStrategies.TriangleFootprints.ThreeDifferences.partial\_application\_models} for the
latter), not every run of HE's machine. All of this
relates Lean models of the dialects to each other and to the step restricted to one ground rule per
machine transition; a placement of a dialect by the rules of its specification is open
(Section~\ref{sec:open}).

\subsection{Evaluation modes decided by types}\label{sec:addendum-bytype}
Like Section~\ref{sec:addendum-evaluation-modes}, this subsection is about a separate core, not about
the step or any evaluator, and its results decide no design. In a typed refinement of that core every
position gets an expected type from the declarations alone: a term at a position typed \mt{Atom} or
\mt{(Code T)} is passed as written, and elsewhere it is evaluated. Four typed designs are given by two
choices: whether held text runs at a position typed \mt{\%Undefined\%} or by a type variable that does
not occur in the code's type, and whether plain \mt{Atom} text runs at value positions
(\lean{DialectsAsStrategies.EvaluationModes.ByType.TDesign}). The design \mt{byType} takes both. On
every program of the earlier study's surface, by-name and \mt{byType} elaborate to the same core
program, and every typed program run in \mt{byType} is the by-name run of its projection, with
\mt{(Code T)} read as \mt{Atom} (\lean{DialectsAsStrategies.EvaluationModes.ByType.byName\_eq\_byType},
\lean{DialectsAsStrategies.EvaluationModes.ByType.byType\_eq\_byName\_proj}); code value is likewise the
design \mt{codeOnly} (\lean{DialectsAsStrategies.EvaluationModes.ByType.codeValue\_eq\_codeOnly}).
Weakening a value result changes no run that refused nothing, for every program, exactly in the designs
whose default runs (\lean{DialectsAsStrategies.EvaluationModes.ByType.resultWeakening\_iff}); reading
every \mt{(Code T)} declaration as \mt{Atom} does so exactly in \mt{byType}
(\lean{DialectsAsStrategies.EvaluationModes.ByType.heldWeakening\_iff}); and the two hold together
exactly there (\lean{DialectsAsStrategies.EvaluationModes.ByType.weakening\_determines}). For any answer
function at all, \mt{\%Undefined\%} cannot weaken both \mt{Number} and \mt{(Code Number)} when the two
answer differently, and every design meets that premise
(\lean{DialectsAsStrategies.EvaluationModes.ByType.no\_cross\_mode\_weakening},
\lean{DialectsAsStrategies.EvaluationModes.ByType.cross\_mode\_premises}). So the bit that remains is
whether to accept the second law, which reads \mt{Atom} as code of unknown type: \mt{byType} keeps both
laws, and \mt{codeOnly} keeps the first and not the second
(\lean{DialectsAsStrategies.EvaluationModes.ByType.remaining\_bit}). If \mt{Atom} is read as data, that
change crosses modes, and no answer function weakens across modes; the results do not decide which
reading of \mt{Atom} holds. With alternatives, \mt{let} holds
in every design for every program (\lean{DialectsAsStrategies.EvaluationModes.ByType.keepsLet}); on a
choice between two answers used twice, a value parameter gives the two equal pairs and held text all
four pairs (\lean{DialectsAsStrategies.EvaluationModes.ByType.dup\_two\_rows}); and dropping a check can lose every
answer: under \mt{Number} one alternative is refused and the other answers \mt{done}, and under
\mt{\%Undefined\%} the run has no answer at any fuel
(\lean{DialectsAsStrategies.EvaluationModes.ByType.bag\_strong},
\lean{DialectsAsStrategies.EvaluationModes.ByType.bag\_weak\_none}). Read as alternatives rather than as
checks, a declared row that refuses nothing changes the answers
(\lean{DialectsAsStrategies.EvaluationModes.ByType.rows\_as\_alternatives\_break\_weakening}). By-need is
not in this core, and a tie from its runs to the step is open.

\subsection{Prices, checkers and regions}\label{sec:addendum-prices}
This extends Section~\ref{sec:addendum-types-classes}. ``Exactly'' below means an equivalence proved
in Lean without \lean{Classical.choice}, between a principle and a statement about a class of systems
in universe zero; it is not a theorem inside a subsystem of second-order arithmetic. On systems with
explicit successor lists, decided observations and countably many atoms and labels, finitary adequacy
is exactly LLPO (\lean{LogicOverTheStep.Prices.AdequacyPrice.listAdequacy\_iff\_llpo}). The reflection
principle that logical equivalence implies bisimilarity on every system image-finite up to its
equations is exactly excluded middle, already with decided observations
(\lean{LogicOverTheStep.Prices.AdequacyPrice.imageFiniteReflection\_iff\_em},
\lean{LogicOverTheStep.Prices.AdequacyPrice.coverAdequacy\_iff\_em}); this quantifies over all such
systems, and does not say that the adequacy of one given system needs excluded middle. A
distinguishing formula refutes bisimilarity with no principle
(\lean{LogicOverTheStep.Prices.AdequacyPrice.not\_bisimilar\_of\_apart}). On the chains built from binary
sequences, ``a search that cannot fail forever succeeds'' is exactly Markov's principle, ``reaching is
decided'' exactly LPO and ``avoiding is decided'' exactly WLPO
(\lean{LogicOverTheStep.Prices.SearchPrice.mp\_iff\_search},
\lean{LogicOverTheStep.Prices.SearchPrice.lpo\_iff\_search},
\lean{LogicOverTheStep.Prices.SearchPrice.wlpo\_iff\_safety}); on every search system with successor
lists and a decided goal these principles are upper bounds
(\lean{LogicOverTheStep.Prices.SearchPrice.reaches\_of\_not\_avoids},
\lean{LogicOverTheStep.Prices.SearchPrice.reaches\_decided}). For binary systems, a bound from ``every
run finishes'' in bar form is exactly the fan theorem for decidable bars, and from ``no run is
infinite'' the fan theorem with Markov's principle
(\lean{LogicOverTheStep.Prices.TerminationPrice.barBound\_iff\_fan},
\lean{LogicOverTheStep.Prices.TerminationPrice.runBound\_iff}), while a rank into the naturals or an
inductive proof of well-foundedness gives a bound with no principle
(\lean{LogicOverTheStep.Prices.TerminationPrice.bounded\_of\_rank},
\lean{LogicOverTheStep.Prices.TerminationPrice.bounded\_of\_acc}). General finite branching is open;
with infinite branching, runs have every length
(\lean{LogicOverTheStep.Prices.TerminationPrice.Controls.fan\_unbounded}); and in the computable
reading Markov's principle holds and the fan theorem fails
(\lean{LogicOverTheStep.Prices.TerminationPrice.computable\_reading}). No price is proved for
alternating fixed points.

Sound checkers, computable functions on the naturals that say yes only on the property, have no top
when the property is not decidable: every one has a strict extension and two incomparable ones, and
there are infinite strictly increasing chains
(\lean{LogicOverTheStep.Prices.CheckerTower.exists\_strict\_extension},
\lean{LogicOverTheStep.Prices.CheckerTower.exists\_incomparable\_extensions},
\lean{LogicOverTheStep.Prices.CheckerTower.exists\_strict\_chain}), as for least-fixpoint membership built
from halting (\lean{LogicOverTheStep.Prices.CheckerTower.reachesHalting\_no\_top}); a decidable property
has a greatest one (\lean{LogicOverTheStep.Prices.CheckerTower.greatest\_of\_computablePred}). This is
about which instances a checker answers, not about speed. For a computable certificate relation, the
search for a certificate converges on every instance that a sound checker accepts
(\lean{LogicOverTheStep.Prices.CheckerTower.certificate\_choice}).

For a computable certificate relation and a decidable region, the search for a certificate is
computably bounded on the region exactly when the property is decided there, and exactly when a
computable certificate exists there (\lean{LogicOverTheStep.Prices.Regions.region\_theorem},
\lean{LogicOverTheStep.Prices.Regions.certificate\_iff\_decided}), which can be the least one
(\lean{LogicOverTheStep.Prices.Regions.least\_certificate\_of\_decidedOn}). The region must be
decidable: on the certified instances of an undecidable property the property is decided and the
search is not bounded (\lean{LogicOverTheStep.Prices.Regions.decided\_not\_bounded}). For an
undecidable property every computable bound is exceeded on infinitely many certified instances
(\lean{LogicOverTheStep.Prices.Regions.exceeded\_infinite}), and every bounded region misses a certified
and an uncertified instance (\lean{LogicOverTheStep.Prices.Regions.misses\_certified},
\lean{LogicOverTheStep.Prices.Regions.misses\_uncertified}). A region does not find certificates sooner:
on a certified instance the tree's budgeted procedure, the ambient search and a search with a hint
answer at budget $k$ exactly when the least certificate is at most $k$
(\lean{LogicOverTheStep.Prices.Regions.same\_yes\_condition}). What a bounded region adds is a cost on
its uncertified instances, where the ambient search has none
(\lean{LogicOverTheStep.Prices.Regions.bounded\_against\_unbounded}). The cost is the least step index at
which Mathlib's bounded evaluation returns; it is more than the input, and under composition at most
the larger of the two costs (\lean{LogicOverTheStep.Prices.Regions.input\_lt\_cost},
\lean{LogicOverTheStep.Prices.Regions.cost\_comp\_le}, \lean{LogicOverTheStep.Prices.Regions.not\_additive}),
so it is not a count of steps, and no speed-up factor is claimed. A finite table of answers has a cost
bound fixed before the program (\lean{LogicOverTheStep.Prices.Regions.fast\_path}), the bound depending
on the table and its answers, and every program has equivalents above any bound
(\lean{LogicOverTheStep.Prices.Regions.padded}). Checking a supplied certificate is decidable exactly
when the relation is computable, having one is semidecidable, and the truth of ``this program never
halts on input $0$'' is not semidecidable, so a computable certificate relation that certifies only
true such claims misses one (\lean{LogicOverTheStep.Prices.Regions.check\_decidable\_iff},
\lean{LogicOverTheStep.Prices.Regions.certified\_re}, \lean{LogicOverTheStep.Prices.Regions.never\_halting\_not\_re},
\lean{LogicOverTheStep.Prices.Regions.sound\_misses\_never\_halting}). On the tree's own objects, the
unifier's run finishes exactly when its budget is at least its number of rounds, which is defined by
its rank and is not proved computable
(\lean{LogicOverTheStep.Prices.Regions.TreeUnifier.drive\_finished\_iff}), and the definition
checker answers its table of seven definitions from budget 16 on
(\lean{LogicOverTheStep.Prices.Regions.TreeChecker.table\_answered}); typed conversion on a qualified
package is decided as a proposition (\lean{NativeTypesOnTheStep.Computing.Qualified.decide}), and no
computable bound is proved for it.

\subsection{Numbers on the step}\label{sec:addendum-numbers}
The terms here are the module's own, not the tree's MeTTa atoms: stored numerals and two-argument
calls of grounded operations, with a table in which a refused call stays as written. The knowledge
identifies two terms when every instance of them evaluates to terms of equal value, and the step is
the tree's first-order step over that equality. So, by construction, the knowledge identifies $5$ and
$5.0$ for every table (\lean{LogicOverTheStep.Numbers.OnTheStep.numerals\_identified}). The theorems with
content: the knowledge identifies every pair of storage variants exactly when every operation of the
table respects the value lens (\lean{LogicOverTheStep.Numbers.OnTheStep.identifies\_variants\_iff});
with such a table, re-storing every numeral of a program, its labels, its knowledge and its query by an
injective value-preserving map changes no answer reachable by runs
(\lean{LogicOverTheStep.Numbers.OnTheStep.program\_answers\_iff}); and a table with one operation that
splits a pair of storages has a one-row configuration at which the firing that checks the call's count
is admitted and the firing that checks its storage variant is refused
(\lean{LogicOverTheStep.Numbers.OnTheStep.one\_row\_step\_converse}), as for a storage-dependent division
modelled on HE's measured \mt{/} (\lean{LogicOverTheStep.Numbers.OnTheStep.division\_one\_row}). Rows a
rule makes and labels are re-stored only by one map applied everywhere. The naturals, the whole numbers
and the rationals are types of the value lens, strictly included
(\lean{LogicOverTheStep.Numbers.Tower.naturals\_isConcept}, \lean{LogicOverTheStep.Numbers.Tower.tower\_strict}),
and ``stored as an integer'' is a type of neither the value lens nor the storage lens
(\lean{LogicOverTheStep.Numbers.Tower.intStored\_not\_value\_type},
\lean{LogicOverTheStep.Numbers.Tower.intStored\_not\_storage\_type}).

In a language of binary operations at one type and Boolean tests, whose operations' sets of types are
closed under joins, as on a chain, coherence holds exactly when conversions compose and operations and
tests commute with them (\lean{LogicOverTheStep.Numbers.Coherence.Language.coherent\_iff}); any failure
to commute gives an expression with two elaborations of different values
(\lean{LogicOverTheStep.Numbers.Coherence.Language.commute\_needed}), and with \mt{==} in scope a silent
conversion must be injective (\lean{LogicOverTheStep.Numbers.Coherence.Language.criterion}); division by
zero is not treated in that language. The tower
$\mathbb{N}\subset\mathbb{Z}\subset\mathbb{Q}$ with exact division and a separately named quotient is
coherent (\lean{LogicOverTheStep.Numbers.Coherence.named\_coherent}); with truncated subtraction, or with
the storage-dependent division, it is not
(\lean{LogicOverTheStep.Numbers.Coherence.truncated\_incoherent},
\lean{LogicOverTheStep.Numbers.Coherence.storageDiv\_incoherent}). Every map from the whole numbers
into a finite set of values identifies two different whole numbers, no addition on a finite set with
a nonzero element commutes with its inclusion into $\mathbb{Q}$, and a nearest rounding of a sum is the
sum exactly when the sum is representable (\lean{LogicOverTheStep.Numbers.Coherence.rounding\_merges},
\lean{LogicOverTheStep.Numbers.Coherence.no\_commuting\_addition},
\lean{LogicOverTheStep.Numbers.Coherence.roundedAdd\_exact\_iff}); floats are not modelled. A fast path
that may read storage is licensed when on its region it answers what the value operation answers;
licensed paths, in any order, compute the value operation, and adding a path changes no answer exactly
when it is licensed (\lean{LogicOverTheStep.Numbers.FastPaths.withFasts\_value},
\lean{LogicOverTheStep.Numbers.FastPaths.adding\_iff\_licensed}). The 64-bit paths for \mt{+}, \mt{-}
and \mt{*} are licensed (\lean{LogicOverTheStep.Numbers.FastPaths.add64\_licensed},
\lean{LogicOverTheStep.Numbers.FastPaths.sub64\_licensed},
\lean{LogicOverTheStep.Numbers.FastPaths.mul64\_licensed}), and the table of HE's grounded operations as
measured, which is HE's on calls that pass HE's type check of the operation
(Section~\ref{sec:addendum-tables}), is the 64-bit path on its region, with an overflow error outside it
(\lean{LogicOverTheStep.Numbers.FastPaths.fast\_is\_measured},
\lean{LogicOverTheStep.Numbers.FastPaths.measured\_add}, \lean{LogicOverTheStep.Numbers.FastPaths.measured\_sub},
\lean{LogicOverTheStep.Numbers.FastPaths.measured\_mul}); an erroring product splits $2^{62}$
stored as an integer from $2^{62}$ stored as a fraction
(\lean{LogicOverTheStep.Numbers.FastPaths.checkedMul\_not\_respects}). No work is counted and no speed is
claimed. No operation on values is both the truncating quotient of whole numbers and the division of
the rationals (\lean{LogicOverTheStep.Numbers.CanonicalStorage.no\_value\_operation\_both}). Truncated
and floored quotients both commute with the inclusion of the naturals and differ on negative arguments
(\lean{LogicOverTheStep.Numbers.Dialects.truncated\_commutes},
\lean{LogicOverTheStep.Numbers.Dialects.floored\_commutes},
\lean{LogicOverTheStep.Numbers.Dialects.quotients\_differ}); the choice between them is not made.

\subsection{Certificates on classes and on supports}\label{sec:addendum-compact}
The certificate of Section~\ref{sec:addendum-certificates} has two smaller forms, each sound by
reduction to it. Along a lumping of a chain onto a smaller one, a coupling bound of the quotient pulls
back to the chain at every discount $0\le c$
(\lean{LogicOverTheStep.Distinctions.Certificates.Compact.couplingBound\_pullback}), and the metric is
invariant under lumping at $0<c\le 1$
(\lean{LogicOverTheStep.Distinctions.Certificates.Compact.bisimulationMetric\_lump}). A certificate on
classes is a map onto classes, a chain on them and a certificate of that chain; its checker verifies
the lumping with sums and equalities and runs the earlier checker on the quotient, and if it accepts,
the metric is at most the kernel at the classes, at every $0\le c$
(\lean{LogicOverTheStep.Distinctions.Certificates.Compact.classCheck\_sound}). A supported certificate
adds, for each action and state, a set holding the successors of positive mass, and its checker reads
tables only on products of supports
(\lean{LogicOverTheStep.Distinctions.Certificates.Compact.supportedCheck\_sound}); through a quotient
with an error, the distance is bounded by the errors at the two states plus the quotient's kernel
(\lean{LogicOverTheStep.Distinctions.Certificates.Compact.abstractionCheck\_sound}). Two weaker lumping
checks, by class averages and by supports, each accept a certificate whose bound is false
(\lean{LogicOverTheStep.Distinctions.Certificates.Compact.Controls.average\_unsound},
\lean{LogicOverTheStep.Distinctions.Certificates.Compact.Controls.support\_unsound}), and a convex
combination of accepted certificates need not be a coupling bound
(\lean{LogicOverTheStep.Distinctions.Certificates.Compact.UpToControls.average\_not\_couplingBound}).

Costs are counted by a work function written beside each checker, which follows its code branch by
branch and counts products, additions, comparisons and tests; it is not derived from the checkers by a
theorem, it is not time, and reading a certificate's own data is free in it. In that count every
accepted certificate of the earlier form costs that checker at least $|A|\,n(n-1)(n^2+1)$ products and
$|A|\,n(n-1)(n^2+n)$ additions
(\lean{LogicOverTheStep.Distinctions.Certificates.Compact.prefixpointWork\_lower}); this is a bound for
the tree's checker of that form, not for every checker of it. On a family built for the purpose,
$m+1$ copies of each of the four states of a leaking chain, the checker on classes is at least four
times cheaper than every accepted earlier certificate for every $m\ge 1$
(\lean{LogicOverTheStep.Distinctions.Certificates.Compact.Family.four\_times\_cheaper}), and its
certificate four times smaller when the earlier certificate proves bounds at least as strong at the
stated pairs (\lean{LogicOverTheStep.Distinctions.Certificates.Compact.Family.four\_times\_smaller});
with one state per class the class form costs more
(\lean{LogicOverTheStep.Distinctions.Certificates.Compact.Family.loses\_without\_classes}). For every
chain, an explicit size condition gives the factor four for either form
(\lean{LogicOverTheStep.Distinctions.Certificates.Compact.four\_times\_cheaper\_of\_large},
\lean{LogicOverTheStep.Distinctions.Certificates.Compact.supported\_four\_times\_cheaper\_of\_large}); with
one action and one observation it is met from 9 states with four classes and from 10 states with
supports of at most two states
(\lean{LogicOverTheStep.Distinctions.Certificates.Compact.four\_classes\_from\_nine},
\lean{LogicOverTheStep.Distinctions.Certificates.Compact.two\_successors\_from\_ten}). On a closed region
of the step under the policy of Section~\ref{sec:addendum-certificates}, with 11 states that lump onto
five classes, the certificate on classes is accepted by kernel evaluation and shows, in the step's own
system, two refusing states with the same reading to be Larsen--Skou bisimilar
(\lean{LogicOverTheStep.Distinctions.Certificates.Compact.StepClasses.late\_bisimilar\_three}), and its
checker is at least four times cheaper than every accepted earlier certificate of the region
(\lean{LogicOverTheStep.Distinctions.Certificates.Compact.StepClasses.step\_four\_times\_cheaper}); the
region is shown at 7, 9 and 11 states, not at every size.

The Lean declarations cited in this addendum use \lean{propext}, \lean{Classical.choice} and
\lean{Quot.sound}, or fewer; none uses \lean{sorryAx}.

\section{Receipts}\label{sec:receipts}

The recorded Lean audit covers the metatheory declarations; the C implementation
and HOL machine-code theorem above have their separate evidence boundaries.
Unqualified Lean names in an example are in
the namespace in its title, and those in the instance columns of the feature table are in
\lean{InterfaceTable}. One audit file imports 229 modules that own the cited declarations, with
warnings as errors, and prints the axioms of all 492 of them: 429 use \lean{propext},
\lean{Classical.choice} and \lean{Quot.sound}, 42 use \lean{propext} and \lean{Quot.sound}, 5 use
\lean{propext} alone, and 16 use no axiom; none uses \lean{sorryAx}, and every name resolves to a
declaration of the \lean{Mettapedia} tree. Four further declarations cited in
Sections~\ref{sec:core} and~\ref{sec:map}, \lean{Interface.EffectCore}, \lean{Interface.Step.Guarded},
\lean{Interface.Step.Standard} and \lean{Race.TriggeredProgram.tensor\_ports}, are checked by a
second file with the same settings; each uses \lean{propext}, \lean{Classical.choice} and
\lean{Quot.sound}. The thirteen declarations of the step of a program cited in
Sections~\ref{sec:core} and~\ref{sec:recommendation} (namespace \lean{RuleSyntax}) are checked by a
third file with the same settings: twelve use those three axioms and one uses \lean{propext} and
\lean{Quot.sound}. The 270 declarations cited in Section~\ref{sec:addendum-nov13}, which include every
name that earlier sections take from it, are checked by a fourth file with the same settings:
173 use \lean{propext}, \lean{Classical.choice} and \lean{Quot.sound}, 64 use \lean{propext} and
\lean{Quot.sound}, 15 use \lean{propext} alone, and 18 use no axiom; none uses \lean{sorryAx}.

\section*{Vocabulary}
\begin{description}
\item[frame] an object of the frame action: a declaration family, a formed telescope, or the one
  object of a substitution action;
\item[take, read, make] what a firing removes from the space, observes there without removing, and
  adds to it with fresh identities;
\item[tell] an equation a firing adds to its information;
\item[guard] an ask of a firing at every solution of its information;
\item[knowledge] the equation store of a state; \textbf{information}: that of a firing at its target;
\item[space] the present occurrences of a state, or those at one \textbf{handle};
\item[rights] what a firing may do at a handle: \mt{receive} (take), \mt{peek} (read), \mt{send}
  (make), \mt{observe} (query);
\item[schema] a rule whose variables are opened in a fresh scope at each firing; its \textbf{capture
  set} lists the caller variables it shares;
\item[birth/guard cut] the information at which a firing's guards are asked;
\item[predicate query] a proposition about the occurrences in a region. An \textbf{exact query}
  additionally fixes the set of identities at every solution;
\item[dead names] names of a larger frame that no row, knowledge or ledger entry mentions;
\item[authority restriction] the firings a process may fire, by publication rights;
\item[fibre] a restriction of the frames with stronger guarantees.
\end{description}

\begin{thebibliography}{9}
\addcontentsline{toc}{part}{References}

\bibitem{meredith-radestock-rho}
L. Gregory Meredith and Matthias Radestock.
A reflective higher-order calculus.
\emph{Electronic Notes in Theoretical Computer Science}, 141(5), 49--67, 2005.
\href{https://doi.org/10.1016/j.entcs.2005.05.016}{doi:10.1016/j.entcs.2005.05.016}.

\bibitem{f1r3lang-compiler-draft}
F1R3FLY.
\emph{F1R3Lang compiler specification}, design frame and decision register.
\href{https://github.com/F1R3FLY-io/publications/blob/9e6d7381568072d294ef6c6b5810299650c4b772/rho2metal/f1r3lang-compiler-spec/body-1-frame.tex}{Source revision \texttt{9e6d73815680}}.

\bibitem{saraswat-rinard-ccp}
Vijay A. Saraswat and Martin Rinard.
Concurrent constraint programming.
\emph{Proceedings of POPL}, 232--245, 1990.
\href{https://doi.org/10.1145/96709.96733}{doi:10.1145/96709.96733}.

\bibitem{fruehwirth-chr}
Thom Fr\"uhwirth.
Theory and practice of constraint handling rules.
\emph{Journal of Logic Programming}, 37(1--3), 95--138, 1998.
\href{https://doi.org/10.1016/S0743-1066(98)10005-5}{doi:10.1016/S0743-1066(98)10005-5}.

\bibitem{mork-sources-sinks}
MORK project.
\href{https://github.com/trueagi-io/MORK/wiki/Sources-and-Sinks}{Sources and sinks}.
MM2 documentation.

\bibitem{ports-grittner-ssi}
Dan R. K. Ports and Kevin Grittner.
Serializable snapshot isolation in PostgreSQL.
\emph{Proceedings of the VLDB Endowment}, 5(12), 1850--1861, 2012.
\href{https://arxiv.org/abs/1208.4179}{arXiv:1208.4179}.

\end{thebibliography}

\end{document}
